\documentclass{article}

\def\npart {IA}
\def\nterm {Michaelmas}
\def\nyear {2023}
\def\nlecturer {Prof N.\ Peake}
\def\ncourse {Vector and Matrices}

\setcounter{tocdepth}{2}

\usepackage{alltt}
\usepackage{amsmath, amsfonts, amssymb, amsthm}
\usepackage{booktabs}
\usepackage{caption}
\usepackage{enumitem}
\usepackage{fancyhdr}
\usepackage{graphicx}
\usepackage{mathdots}
\usepackage{mathtools}
\usepackage{microtype}
\usepackage{multirow}
\usepackage{pdflscape}
\usepackage{pgfplots}
\usepackage{siunitx}
\usepackage{textcomp}
\usepackage{slashed}
\usepackage{tabularx}
\usepackage{tikz}
\usepackage{tkz-euclide}
\usepackage[normalem]{ulem}
\usepackage[all]{xy}

\pgfplotsset{compat=1.18}
\pagestyle{fancyplain}

\lhead{\emph{\nouppercase{\leftmark}}}
\rhead{
  \ifnum\thepage=1
  \else
    \npart\ \vline\ \ncourse
  \fi}

\def\nauthor{Marcus Ng}
\author{Based on lectures by \nlecturer \\\small Notes taken by \nauthor}
\date{\nterm\ \nyear}
\title{Part \npart\ --- \ncourse}

\newcommand*{\Cdot}{{\raisebox{-0.25ex}{\scalebox{1.5}{$\cdot$}}}}
\setlist[enumerate,1]{label={(\roman*)}}

\newcommand {\pd}[2][ ]{
  \ifx #1 { }
    \frac{\partial}{\partial #2}
  \else
    \frac{\partial^{#1}}{\partial #2^{#1}}
  \fi
}

% Theorems
\theoremstyle{definition}
\newtheorem*{aim}{Aim}
\newtheorem*{axiom}{Axiom}
\newtheorem*{claim}{Claim}
\newtheorem*{cor}{Corollary}
\newtheorem*{conjecture}{Conjecture}
\newtheorem*{defi}{Definition}
\newtheorem*{eg}{Example}
\newtheorem*{ex}{Exercise}
\newtheorem*{fact}{Fact}
\newtheorem*{law}{Law}
\newtheorem*{lemma}{Lemma}
\newtheorem*{notation}{Notation}
\newtheorem*{prop}{Proposition}
\newtheorem*{question}{Question}
\newtheorem*{problem}{Problem}
\newtheorem*{rrule}{Rule}
\newtheorem*{thm}{Theorem}
\newtheorem*{assumption}{Assumption}
\newtheorem*{assert}{Assertion}

\newtheorem*{remark}{Remark}
\newtheorem*{warning}{Warning}
\newtheorem*{exercise}{Exercise}

\newtheorem{nthm}{Theorem}[section]
\newtheorem{nlemma}[nthm]{Lemma}
\newtheorem{nprop}[nthm]{Proposition}
\newtheorem{ncor}[nthm]{Corollary}

\SetEnumitemKey{cases}{
  label=\underline{Case~\arabic*:},
  ref={\arabic*},
  align=left,
  leftmargin=\parindent
}

% \renewcommand{\labelitemi}{--}
% \renewcommand{\labelitemii}{$\circ$}
% \renewcommand{\labelenumi}{(\roman{*})}

\let\stdsection\section
\renewcommand\section{\newpage\stdsection}

% Strike through
\def\st{\bgroup \ULdepth=-.55ex \ULset}


%%%%%%%%%%%%%%%%%%%%%%%%%
%%%%% Maths Symbols %%%%%
%%%%%%%%%%%%%%%%%%%%%%%%%

% Logical symbols
\newcommand{\contradiction}{
  \begin{tikzpicture}[rotate=45,x=0.5ex,y=0.5ex]
  \draw[line width=.1ex] (0,2) -- (3,2) (0,1) -- (3,1) (1,3) -- (1,0) (2,3) -- (2,0);
  \end{tikzpicture}
}


\let\U\relax
\let\C\relax
\let\G\relax

% Matrix groups
% \newcommand{\GL}{\mathrm{GL}}
% \newcommand{\Or}{\mathrm{O}}
% \newcommand{\PGL}{\mathrm{PGL}}
% \newcommand{\PSL}{\mathrm{PSL}}
% \newcommand{\PSO}{\mathrm{PSO}}
% \newcommand{\PSU}{\mathrm{PSU}}
% \newcommand{\SL}{\mathrm{SL}}
% \newcommand{\SO}{\mathrm{SO}}
% \newcommand{\Spin}{\mathrm{Spin}}
% \newcommand{\Sp}{\mathrm{Sp}}
% \newcommand{\SU}{\mathrm{SU}}
\newcommand{\U}{\mathrm{U}}
% \newcommand{\Mat}{\mathrm{Mat}}

% Matrix algebras
% \newcommand{\gl}{\mathfrak{gl}}
% \newcommand{\ort}{\mathfrak{o}}
% \newcommand{\so}{\mathfrak{so}}
% \newcommand{\su}{\mathfrak{su}}
% \newcommand{\uu}{\mathfrak{u}}
% \renewcommand{\sl}{\mathfrak{sl}}

% Special sets
\newcommand{\C}{\mathbb{C}}
\newcommand{\CP}{\mathbb{CP}}
\newcommand{\GG}{\mathbb{G}}
\newcommand{\N}{\mathbb{N}}
\newcommand{\Q}{\mathbb{Q}}
\newcommand{\R}{\mathbb{R}}
\newcommand{\RP}{\mathbb{RP}}
\newcommand{\T}{\mathbb{T}}
\newcommand{\Z}{\mathbb{Z}}
\renewcommand{\H}{\mathbb{H}}

% Brackets
\newcommand{\abs}[1]{\left\lvert #1\right\rvert}
\newcommand{\bket}[1]{\left\lvert #1\right\rangle}
\newcommand{\brak}[1]{\left\langle #1 \right\rvert}
\newcommand{\braket}[2]{\left\langle #1\middle\vert #2 \right\rangle}
\newcommand{\bra}{\langle}
\newcommand{\ket}{\rangle}
\newcommand{\norm}[1]{\left\lVert #1\right\rVert}
\newcommand{\normalorder}[1]{\mathop{:}\nolimits\!#1\!\mathop{:}\nolimits}
\newcommand{\tv}[1]{|#1|}
\renewcommand{\vec}[1]{\boldsymbol{\mathbf{#1}}}

% not-math
% \newcommand{\bolds}[1]{{\bfseries #1}}
% \newcommand{\cat}[1]{\mathsf{#1}}
% \newcommand{\ph}{\,\cdot\,}
% \newcommand{\term}[1]{\emph{#1}\index{#1}}
% \newcommand{\phantomeq}{\hphantom{{}={}}}

% Probability
% \DeclareMathOperator{\Bernoulli}{Bernoulli}
% \DeclareMathOperator{\betaD}{beta}
% \DeclareMathOperator{\bias}{bias}
% \DeclareMathOperator{\binomial}{binomial}
% \DeclareMathOperator{\corr}{corr}
% \DeclareMathOperator{\cov}{cov}
% \DeclareMathOperator{\gammaD}{gamma}
% \DeclareMathOperator{\mse}{mse}
% \DeclareMathOperator{\multinomial}{multinomial}
% \DeclareMathOperator{\Poisson}{Poisson}
% \DeclareMathOperator{\var}{var}
% \newcommand{\E}{\mathbb{E}}
% \newcommand{\Prob}{\mathbb{P}}

% Algebra
% \DeclareMathOperator{\adj}{adj}
% \DeclareMathOperator{\Ann}{Ann}
% \DeclareMathOperator{\Aut}{Aut}
% \DeclareMathOperator{\Char}{char}
% \DeclareMathOperator{\disc}{disc}
% \DeclareMathOperator{\dom}{dom}
% \DeclareMathOperator{\fix}{fix}
% \DeclareMathOperator{\Hom}{Hom}
% \DeclareMathOperator{\id}{id}
% \DeclareMathOperator{\image}{image}
% \DeclareMathOperator{\im}{im}
% \DeclareMathOperator{\re}{re}
% \DeclareMathOperator{\tr}{tr}
% \DeclareMathOperator{\Tr}{Tr}
% \newcommand{\Bilin}{\mathrm{Bilin}}
% \newcommand{\Frob}{\mathrm{Frob}}

% Others
% \newcommand\ad{\mathrm{ad}}
% \newcommand\Art{\mathrm{Art}}
% \newcommand{\B}{\mathcal{B}}
% \newcommand{\cU}{\mathcal{U}}
% \newcommand{\Der}{\mathrm{Der}}
% \newcommand{\D}{\mathrm{D}}
% \newcommand{\dR}{\mathrm{dR}}
% \newcommand{\exterior}{\mathchoice{{\textstyle\bigwedge}}{{\bigwedge}}{{\textstyle\wedge}}{{\scriptstyle\wedge}}}
% \newcommand{\F}{\mathbb{F}}
% \newcommand{\G}{\mathcal{G}}
% \newcommand{\Gr}{\mathrm{Gr}}
% \newcommand{\haut}{\mathrm{ht}}
% \newcommand{\Hol}{\mathrm{Hol}}
% \newcommand{\hol}{\mathfrak{hol}}
% \newcommand{\Id}{\mathrm{Id}}
% \newcommand{\lie}[1]{\mathfrak{#1}}
% \newcommand{\op}{\mathrm{op}}
% \newcommand{\Oc}{\mathcal{O}}
% \newcommand{\pr}{\mathrm{pr}}
% \newcommand{\Ps}{\mathcal{P}}
% \newcommand{\pt}{\mathrm{pt}}
% \newcommand{\qeq}{\mathrel{``{=}"}}
% \newcommand{\Rs}{\mathcal{R}}
% \newcommand{\Vect}{\mathrm{Vect}}
% \newcommand{\wsto}{\stackrel{\mathrm{w}^*}{\to}}
% \newcommand{\wt}{\mathrm{wt}}
% \newcommand{\wto}{\stackrel{\mathrm{w}}{\to}}
% \renewcommand{\d}{\mathrm{d}}
% \renewcommand{\P}{\mathbb{P}}
%\renewcommand{\F}{\mathcal{F}}

\let\Im\relax
\let\Re\relax

% \DeclareMathOperator{\area}{area}
% \DeclareMathOperator{\card}{card}
% \DeclareMathOperator{\ccl}{ccl}
% \DeclareMathOperator{\ch}{ch}
% \DeclareMathOperator{\cl}{cl}
% \DeclareMathOperator{\cls}{\overline{\mathrm{span}}}
% \DeclareMathOperator{\coker}{coker}
% \DeclareMathOperator{\conv}{conv}
% \DeclareMathOperator{\cosec}{cosec}
% \DeclareMathOperator{\cosech}{cosech}
% \DeclareMathOperator{\covol}{covol}
% \DeclareMathOperator{\diag}{diag}
% \DeclareMathOperator{\diam}{diam}
% \DeclareMathOperator{\Diff}{Diff}
% \DeclareMathOperator{\End}{End}
% \DeclareMathOperator{\energy}{energy}
% \DeclareMathOperator{\erfc}{erfc}
% \DeclareMathOperator{\erf}{erf}
% \DeclareMathOperator*{\esssup}{ess\,sup}
% \DeclareMathOperator{\ev}{ev}
% \DeclareMathOperator{\Ext}{Ext}
% \DeclareMathOperator{\fst}{fst}
% \DeclareMathOperator{\Fit}{Fit}
% \DeclareMathOperator{\Frac}{Frac}
% \DeclareMathOperator{\Gal}{Gal}
% \DeclareMathOperator{\gr}{gr}
% \DeclareMathOperator{\hcf}{hcf}
\DeclareMathOperator{\Im}{Im}
% \DeclareMathOperator{\Ind}{Ind}
% \DeclareMathOperator{\Int}{Int}
% \DeclareMathOperator{\Isom}{Isom}
% \DeclareMathOperator{\lcm}{lcm}
% \DeclareMathOperator{\length}{length}
% \DeclareMathOperator{\Lie}{Lie}
% \DeclareMathOperator{\like}{like}
% \DeclareMathOperator{\Lk}{Lk}
% \DeclareMathOperator{\Maps}{Maps}
% \DeclareMathOperator{\orb}{orb}
% \DeclareMathOperator{\ord}{ord}
% \DeclareMathOperator{\otp}{otp}
% \DeclareMathOperator{\poly}{poly}
% \DeclareMathOperator{\rank}{rank}
% \DeclareMathOperator{\rel}{rel}
% \DeclareMathOperator{\Rad}{Rad}
\DeclareMathOperator{\Re}{Re}
% \DeclareMathOperator*{\res}{res}
% \DeclareMathOperator{\Res}{Res}
% \DeclareMathOperator{\Ric}{Ric}
% \DeclareMathOperator{\rk}{rk}
% \DeclareMathOperator{\Rees}{Rees}
% \DeclareMathOperator{\Root}{Root}
% \DeclareMathOperator{\sech}{sech}
% \DeclareMathOperator{\sgn}{sgn}
% \DeclareMathOperator{\snd}{snd}
% \DeclareMathOperator{\Spec}{Spec}
% \DeclareMathOperator{\spn}{span}
% \DeclareMathOperator{\stab}{stab}
% \DeclareMathOperator{\St}{St}
% \DeclareMathOperator{\supp}{supp}
% \DeclareMathOperator{\Syl}{Syl}
% \DeclareMathOperator{\Sym}{Sym}
% \DeclareMathOperator{\vol}{vol}
\numberwithin{equation}{section}

\begin{document}
\maketitle{
    \small 
    \noindent\textbf{Complex numbers}\\
    Review of complex numbers, including complex conjugate, inverse, modulus, argument and Argand diagram. Informal treatment of complex logarithm, $n$-th roots and complex powers. de Moivre's theorem.\hspace*{\fill}[2]

    \vspace{10pt}
    \noindent\textbf{Vectors}\\
    Review of elementary algebra of vectors in $\R^3$, including scalar product. Brief discussion of vectors in $\R^n$ and $\C^n$; scalar product and the Cauchy-Schwarz inequality. Concepts of linear span, linear independence, subspaces, basis and dimension.

    \vspace{5pt}
    \noindent Suffix notation: including summation convention, $\delta_{ij}$ and $\varepsilon_{ijk}$. Vector product and triple product: definition and geometrical interpretation. Solution of linear vector equations. Applications of vectors to geometry, including equations of lines, planes and spheres.\hspace*{\fill}[5]

    \vspace{10pt}
    \noindent\textbf{Matrices}\\
    Elementary algebra of $3\times 3$ matrices, including determinants. Extension to $n\times n$ complex matrices. Trace, determinant, non-singular matrices and inverses. Matrices as linear transformations; examples of geometrical actions including rotations, reflections, dilations, shears; kernel and image.\hspace*{\fill}[4]

    \vspace{5pt}
    \noindent Simultaneous linear equations: matrix formulation; existence and uniqueness of solutions, geometric interpretation; Gaussian elimination.\hspace*{\fill}[3]

    \vspace{5pt}
    \noindent Symmetric, anti-symmetric, orthogonal, hermitian and unitary matrices. Decomposition of a general matrix into isotropic, symmetric trace-free and antisymmetric parts.\hspace*{\fill}[1]

    \vspace{10pt}
    \noindent\textbf{Eigenvalues and Eigenvectors}\\
    Eigenvalues and eigenvectors; geometric significance.\hspace*{\fill}[2]

    \vspace{5pt}
    \noindent Proof that eigenvalues of hermitian matrix are real, and that distinct eigenvalues give an orthogonal basis of eigenvectors. The effect of a general change of basis (similarity transformations). Diagonalization of general matrices: sufficient conditions; examples of matrices that cannot be diagonalized. Canonical forms for $2 \times 2$ matrices.\hspace*{\fill}[5]

    \vspace{5pt}
    \noindent Discussion of quadratic forms, including change of basis. Classification of conics, cartesian and polar forms.\hspace*{\fill}[1]

    \vspace{5pt}
    \noindent Rotation matrices and Lorentz transformations as transformation groups.\hspace*{\fill}[1]}

\tableofcontents

%%%%%%%%%%%%%%%%%%%%%%
% Lecture 1: 6/10/2023
%%%%%%%%%%%%%%%%%%%%%%
\section{Complex Numbers}
\subsection{Basic properties}
\begin{defi}[Complex number]
    A \emph{complex number} $z \in \C$ is
    \[
        z = a + ib, \quad a, b \in \R, \quad i = \sqrt{-1}  
    \]
    We define two other functions, $\Re()$ and $\Im()$ such that,
    \[
        \Re(z) = a \quad \Im(z) = b 
    \]
    Clearly, we may note that $\R \subset \C$ in the case where $b = 0$
\end{defi}

We can define algebraic manipulation over the complex numbers,
\begin{align*}
    z_1 \pm z_2 &= a_1 \pm a_2 + i(b_1 \pm b_2) \\
    z_1z_2 &= a_1a_2 - b_1b_2 + i(a_1b_2 + a_2b_1)
\end{align*}
We can define the inverse of a complex number
\begin{equation}
    z^{-1} = \frac{1}{a + ib} = \frac{a - ib}{a^2 + b^2}\label{eq:1}
\end{equation}
this demonstrates that $\C$ is \emph{closed} under addition and multiplication. We also define a complex conjugate of $z = a + ib$ as $\overline{z}$ or $z^*$.
\[
    \overline{z}, z^* = a - ib
\]
The modulus of $z$ as $\abs{z}$
\begin{align*}
    \abs{z} &= \abs{a + ib} = \sqrt{a^2 + b^2} \\
    z\overline{z} &= a^2 + b^2 = \abs{z}^2 \\
    z^{-1} &= \frac{\overline{z}}{\abs{z}^2} \tag{see \eqref{eq:1}}
\end{align*}
We also introduce the Argand Diagram, where the coordinate axes are the Real and Imaginary axes. We can see that $\overline{z}$ is the reflection of $z$ in the $x$-axis.

% draw an argand diagram in tikz
\begin{defi}[Triangle inequality]
    For $z_1, z_2 \in \C$, the triangle inequality states that
    \begin{equation}
        \abs{z_1 + z_2} \leq \abs{z_1} + \abs{z_2}\label{eq:2}
    \end{equation}
\end{defi}

\begin{prop}
    The triangle inequality has another common alternative form,
    \begin{equation}
        \abs{z_1 - z_2} \geq \abs{\abs{z_1} - \abs{z_2}}\label{eq:3}
    \end{equation}
\end{prop}

\begin{proof}
    Write $z_1 = z^{\prime}_1 - z^{\prime}_2$, and $z_2 = z^{\prime}_2$ and substitute them into \eqref{eq:2}. Then,
    \begin{align*}
        &\abs{z^{\prime}_1} \leq \abs{z^{\prime}_1 - z^{\prime}_2} + \abs{z^{\prime}_2} \\
        \Rightarrow& \abs{z^{\prime}_1} - \abs{z^{\prime}_2} \leq\abs{z^{\prime}_1 - z^{\prime}_2}\\
        \Rightarrow& \abs{z^{\prime}_1 - z^{\prime}_2} \geq \abs{z^{\prime}_1} - \abs{z^{\prime}_2} \tag{$\star$}
    \end{align*}
    Notice that our labelling of $z_1$ and $z_2$ was arbitrary, and so we can freely swap the assignment. This gives,
    \[
        \abs{z^{\prime}_1 - z^{\prime}_2} \geq \abs{z^{\prime}_2} - \abs{z^{\prime}_1} \tag{$\star\star$}
    \]
    Since $\abs{z^{\prime}_1 - z^{\prime}_2} \geq (\star) \ \text{and} \ (\star\star)$, hence, \eqref{eq:3}
\end{proof}

\begin{defi}[Polar Representation]
    For a complex number $z = a + ib$ we use the argand diagram to write $z = r(\cos{\theta} + i\sin{\theta})$, where
    \[
        r = \abs{z} \quad \text{and} \quad \theta = \arctan\left(\frac{b}{a}\right) = \arg{(z)}  
    \]
\end{defi}

\begin{remark}
    Note that the mapping $(r, \theta) \rightarrow z$ is unique. However, the reverse $r \rightarrow (r, \theta)$ is \emph{NOT} unique since for each value of $\theta$, $\theta + 2n\pi$ is equally valid.

    In order to make this mapping single valued, we will by convention restrict $\theta$ to the principle value, where
    \[
        -\pi < \theta \leq \pi  
    \]

    This multivalued property becomes quite interesting when we consider complex functions. We can imagine taking the multivalued imaginary $\ln$ function as a multistory carpark, walking round and round up and down. To fix this what we do is to build a brick wall on the upramp (``branch cut'').
\end{remark}

Using complex numbers in this form allows us to view complex arithmetic geometrically. For instance, let
\begin{align*}
    z_k &= x_k + iy_k \quad \text{for} \quad k = 1, 2 \\
    &= r_k(\cos{\theta_k} + i\sin(\theta_k))
\end{align*}
\begin{align*}
    z_1z_2 &= r_1(\cos{\theta_1} + i\sin{\theta_1})r_2(\cos{\theta_2} + i\sin{\theta_2}) \\
    &= r_1r_2[(\cos{\theta_1}\cos{\theta_2} - \sin{\theta_1}\sin{\theta_2}) + i(\cos{\theta_1}\sin{\theta_2} + \sin{\theta_1}\cos{\theta_2})] \\
    &= r_1r_2(\cos{(\theta_1 + \theta_2)} + i\sin{(\theta_1 + \theta_2)}) \numberthis \label{eq:4}
\end{align*}

\subsection{Complex exponential function}
\begin{defi}[Complex exponential function]
    \begin{align*}
        \exp(z) &= 1 + z + \frac{z^2}{2!} + \cdots + \frac{z^n}{n!} + \cdots \\
        &=\sum_{n=0}^{\infty}{\frac{z^n}{n!}} \numberthis \label{eq:5}
    \end{align*}
\end{defi}
This function is ``uniformly convergent'' (something that we will cover next term in analysis). That means that we are able to do lots of algebraic manipulation without too much hesitation. =

\begin{prop}[Multiplication of $\exp()$]
    \[
        \exp(z_1)\exp(z_2) = \exp(z_1 + z_2)  
    \]
\end{prop}
\begin{proof}
    We will consider the double sum
    \begin{align*}
        \sum_{m,n=0}^{\infty}{a_mn} &= a_{00} + a_{01} + a_{02} + a_{03} + \cdots \\
         &+ a_{10} + a_{11} + a_{12} + \cdots \\
         &+ a_{20} + a_{21} + a_{22} + \cdots \\
         &+ a_{30} + a_{31} + a_{32} + \cdots \\
         &\vdots
    \end{align*}
    Then by collecting terms along the diagonal we can regroup the sum into
    \[
        =\sum_{r=0}^{\infty}\sum_{m=0}^{r}{a_{r-m, m}}
    \]
    Hence, 
    \begin{align*}
        \exp(z_1) \cdot \exp(z_2) &= \sum_{r=0}^{\infty}\sum_{m=0}^{r}{\frac{z_1^{r-m}}{(r-m)!}\frac{z_2^{m}}{m!}} \\
        &= \sum_{r=0}^{\infty}{\frac{1}{r!}}\sum_{m=0}^{r}{\frac{r!}{(r-m)!m!}z_1^{r-m}z_2^{m}} \\
        &= \sum_{r=0}^{\infty}{\frac{1}{r!} (z_1 + z_2)^r} \\
        &= \exp(z_1 + z_2)
    \end{align*}
\end{proof}

%%%%%%%%%%%%%%%%%%%%%%
% Lecture 2: 8/10/2023
%%%%%%%%%%%%%%%%%%%%%%
Similarly, we can define the complex $\sin(z)$ and $\cos(z)$ functions.
\begin{defi}[Complex sine and cosine]
    \begin{align*}
        \sin{z} &= \sum_{n=0}^{\infty}{(-1)^n\frac{z^{2n+1}}{(2n+1)!}} \\
        \cos{z} &= \sum_{n=0}^{\infty}{(-1)^n\frac{z^{2n}}{(2n)!}} \numberthis \label{eq:6}
    \end{align*}
\end{defi}

From \eqref{eq:5} we can derive that
\begin{align*}
    e^{iz} &= \sum_{n=0}^{\infty}{\frac{i^nz^n}{n!}} \\
    &= \sum_{n=0}^{\infty}{(-1)^n\frac{z^{2n+1}i^{2n+1}}{(2n+1)!}} + \sum_{n=0}^{\infty}{(-1)^n\frac{z^{2n}i^{2n}}{(2n)!}} \\
    &= i\sin{z} + \cos{z} \numberthis \label{eq:7}
\end{align*}

We may consider $\theta \in \R$ which gives,
\begin{equation}
    z = re^{i\theta} = r(\cos{\theta} + i\sin{\theta}) \label{eq:8}
\end{equation}

\subsection{Roots of unity}
We seek to solve equations of the form $z^n = 1$ where $n \in \Z$. Since this is a polynomial of degree $n$ we expect to find $n$ solutions. By applying the exponential form we found earlier,
\begin{align*}
    z^n = 1 &= \exp(i2 \pi k), \quad k \in \Z \\
    z &= \exp(i2 \pi \frac{k}{n})
\end{align*}
Notice that the roots will repeat after $k = n-1$. Hence the solutions are
\[
    z = \exp(i2 \pi \frac{k}{n}), \quad k = 0, 1, 2, \cdots, n-1
\]

\begin{eg}[$ z^7 = 2$]
    \[
        z = 2^{\frac{1}{7}} \exp(i 2\pi \frac{m}{7}), \quad m = 0, 1, 2, \cdots, 6
    \]
\end{eg}

An interesting point to note is that for the polynomial $z^n - 1 = 0$, if we arbitrarily choose a root $w = \exp(\frac{2\pi i}{n})$, then all the remaining roots will be
\[
    1, w, w^2, w^3, \cdots, w^{n-1}
\]
Further, if we note that since the coefficient of $z^{n-1}$ is $0$, that implies that the sum of the roots must be $0$.
\[
   1 + w + w^2 + w^3 + \cdots + w^{n-1} = 0
\]

\subsection{Complex Logarithm and Powers}
\begin{defi}[Complex Log]
    The \emph{complex log}, $\log{z}$, is a solution of the following
    \[
        e^w = z  
    \]
    Hence, $\exp()$ and $\log()$ are inverse functions, i.e;
    \[
        e^{\log{z}} = z
    \]
    If we write $z = re^{i\theta}$ from \eqref{eq:8} then,
    \begin{align*}
        \log{z} = \log{re^{i\theta}} &= \log{r} + \log{e^{i\theta}} \\
        &= log{\abs{z}} + i\arg{(z)} \numberthis \label{eq:9}
    \end{align*}
    It is important to note that this is not a single valued function. In order to define a single valued $\log$ we will insist that
    \[
        -\pi < \arg{(z)} \leq \pi
    \]
    Hence,
    \begin{equation}
        \log{z} = log{\abs{z}} + i\theta, \quad -\pi < \theta \leq \pi \label{eq:10}
    \end{equation}
\end{defi}

\begin{eg}
    \begin{align*}
        \log{(2i)} &= \log{2e^{i\frac{\pi}{2}}} \\
        &=\log{2} + i\frac{\pi}{2}
    \end{align*}
\end{eg}

\begin{defi}[Complex Power]
    Since we have defined the complex logarithm we can now define the complex power, $z^\alpha$ for $z, \alpha \in \C$ as
    \begin{equation}
        z^\alpha = e^{\alpha \log{z}} \label{eq:11}
    \end{equation}
    Again, this function is multivalued. Specifically using \eqref{eq:10},
    \[
        z^\alpha = e^{\alpha \log{z}}e^{i \alpha \theta}e^{i \alpha 2\pi n}, \quad n \in \Z
    \]
    However, unlike the roots of unity, there is no reason why this series will repeat and terminate. If $\alpha$ is real and rational then will there be finite solutions. Once again, we will make this function single valued by insisting on using the principle value for $\theta$
    \[
        -\pi < \arg{(z)} \leq \pi
    \]
\end{defi}

\begin{ex}\leavevmode
    \begin{enumerate}
        \item \begin{align*}
            (1 + i)^\frac{1}{2} &= e^{\frac{1}{2}\sqrt{2}}e^\frac{i\pi}{8}e^{in\pi}, \quad n = 0, 1 \\
            &= 2^\frac{1}{4}e^\frac{i\pi}{8}, 2^\frac{1}{4}e^\frac{i9\pi}{8} \\
            &= 2^\frac{1}{4}e^\frac{i\pi}{8} \tag{by insiting on the principle value}
        \end{align*}
        \item \begin{align*}
            i^i &= \exp(i \log{(e^\frac{i\pi}{2})}) \tag{from \eqref{eq:11}} \\
            &= \exp(i \frac{i\pi}{2}) \\
            &= \exp(-\frac{\pi}{2})
        \end{align*}
    \end{enumerate}
\end{ex}

\subsection{De Moivre's Theorem}
\begin{thm}[De Moivre's Theorem]
    \begin{equation}
        \cos{n\theta} + i\sin{n\theta} = (\cos{\theta} + i\sin{\theta})^n, \quad n \in \Z \label{eq:12}
    \end{equation}
\end{thm}

\begin{proof}
    By induction, we will first consider the case where $n > 0$. Clearly, $n=1$ is true. Assume $n=k$ is true for $k > 1$.
    \begin{align*}
        (\cos{\theta} + i\sin{\theta})^{k+1} &= (\cos{k\theta} + i\sin{k\theta})(\cos{\theta} + i\sin{\theta}) \\
        &= \cos{k\theta}\cos{\theta} - \sin{k\theta}\sin{\theta} + i(\cos{k\theta}\sin{\theta} + \sin{k\theta}\cos{\theta}) \\
        &= \cos{(k+1)\theta} + i\sin{(k+1)\theta}
    \end{align*}
    Now when $n < 0$, we write $m = -n$ such that $m > 0$
    \begin{align*}
        (\cos{\theta} + i\sin{\theta})^{-m} &= (\cos{m\theta} + i\sin{m\theta})^{-1} \\
        &= \cos{m\theta} - i\sin{m\theta} \tag{by ($\star$)}\\
        &= \cos{(-m\theta)} - i\sin{(-m\theta)}
    \end{align*}
    Note $(\star)$:
    \[
        z^{-1} = \frac{z^*}{\abs{z}^2}    
    \]
\end{proof}

Using De Moivre's theorem we can work out multiple angle formulae by comparing real and imaginary parts:
\begin{eg}
    \begin{align*}
        \cos{5 \theta} + i\sin{5 \theta} &= (\cos{\theta} + i\sin{\theta})^5 \\
        &\vdots \\
        \cos{5 \theta} &= 5\cos{\theta} - 20 \cos^3{\theta} + 16\cos^5{\theta}\\
        \sin{5 \theta} &= 5\cos{\theta} - 20 \cos^3{\theta} + 16\cos^5{\theta}\\
    \end{align*}
\end{eg}
Note that for non-integer powers of $n$ we will have to insist on the principle value \eqref{eq:12}

\subsection{Lines and Circles in $\C$}
\begin{defi}[Straight Line]
    A line $l$ through $z_0 \in \C$ parallel to $w \in \C$ can be generally written
    \[
        z = z_0 + \lambda w, \quad \lambda \in \R    
    \]
    However, it would be better to get rid of the parameterisation. We can use the fact that $\lambda$ must be real and hence its conjugate is equal to itself.
    \begin{align*}
        \lambda &= \frac{z - z_0}{w} \in \R\\
        &\Rightarrow \frac{z - z_0}{w} = \frac{\overline{z} - \overline{z_0}}{\overline{w}} \\
        &\Leftrightarrow z \overline{w} - \overline{z}w = z_0 \overline{w} - \overline{z_0} w
    \end{align*}
\end{defi}

\begin{defi}[Circle]
    A circle can be defined from its center $c \in \C$ and radius $\rho \in \R$. By definition, if $z$ lies on the circle,
    \begin{align*}
        &\abs{z - c} = \rho \\
        \Rightarrow& (z - c)(\overline{z} - \overline{c}) = \rho^2 \\
        \Leftrightarrow& z\overline{z} - \overline{c}z - c\overline{z} = \rho^2 - c\overline{c}
    \end{align*}
\end{defi}

%%%%%%%%%%%%%%%%%%%%%%
% Lecture 3: 11/10/2023
%%%%%%%%%%%%%%%%%%%%%%

\section{Vectors}
\subsection{Definition \& basic properties}
\begin{defi}[Vector]
    A vector $\vec{x}$ has a positive length and a direction. If,
    \[
        \abs{\vec{x}} = 0 \Rightarrow \vec{x} = \vec{0}  
    \]
    If $\abs{\vec{x}} = 1$ then we write $\uvec{x}$ as \emph{unit vector}. Note the distinction between scalar $0$ and vector $\vec{0}$.
\end{defi}

\begin{defi}[Vector Field]
    A vector field $F = F(\vec{x})$ is a vector function of position and maybe time.
\end{defi}

\begin{defi}[Vector addition]
    % insert paralellogram diagram.
    We can define the addition of vectors using the paralellogram rule,
    \begin{align*}
        \vec{a} + \vec{b} &= \vec{c} \\
        \pvec{OA} + \pvec{OB} &= \pvec{OC}
    \end{align*}
    More abstractly, addition will have the following properties
    \begin{enumerate}[eqn]
        \item $\vec{a} + \vec{b} = \vec{b} + \vec{a}$ \hspace*{\fill}(commutativity) \label{eq:2-1}
        \item $\vec{a} + (\vec{b} + \vec{c}) = (\vec{a} + \vec{b}) + \vec{c}$ \hspace*{\fill}(associativity) \label{eq:2-2}
        \item $\exists \vec{0}, \quad \vec{a} + \vec{0} = \vec{a}$ \hspace*{\fill}(unique zero) \label{eq:2-3}
        \item $\vec{a} + -\vec{a} = 0$ \hspace*{\fill}(unique additive inverse) \label{eq:2-4}
    \end{enumerate}
\end{defi}
\begin{defi}[Multiplication by a scalar]
    For $\lambda \in \R$ (we will do complex scalars later in the course), $\lambda\vec{a}$ has magnitutde $\abs{\lambda}\abs{\vec{a}}$ and is parallel/anti-parallel to $\vec{a}$ if $\lambda<0$ \textbackslash $\lambda>0$. 
    Likewise, we have the following properties
    \begin{enumerate}[eqn]
        \item $(\lambda + \mu)\vec{a} = \lambda\vec{a} + \mu\vec{a}$ \label{eq:2-5}
        \item $\lambda(\vec{a} + \vec{b}) = \lambda\vec{a} + \lambda\vec{b}$ \label{eq:2-6}
        \item $\lambda(\mu\vec{a}) = (\lambda\mu)\vec{a}$ \label{eq:2-7}
        \item $1\vec{a} = \vec{a}$ \label{eq:2-8}
    \end{enumerate}
\end{defi}

Implicit in the two definitions is that the set of vectors is closed under addition and scalar multiplication.

\begin{defi}[Vector Space]
    $V$ is a \emph{vector space} if it is a set of vectors which satisfy properties \eqref{eq:2-1} -- \eqref{eq:2-8} inclusive and is closed under addition and multiplication. Note that for now, we will be considering the vector space $V$ over the real numbers $\R$, but we will consider complex fields later. 
\end{defi}

\begin{eg}[$\R^n$ is a vector space]
\end{eg}
\begin{proof}
    \[
        \R^n = \{\vec{x} = (x_1, x_2, \ldots, x_n) \mid x_j \in \R, \ 1 \leq j \leq n\}
    \]
    Consider $\vec{x} , \vec{y} \in \R$
    \[
        \vec{x} + \vec{y} = (x_1 + y_1, x_2 + y_2, \ldots, x_n + y_n,)
    \]
    From this we see that property \eqref{eq:2-1} and \eqref{eq:2-2} is satisfied as we can inherit them from the Real numbers.

    \begin{enumerate}
        \item $\vec{0} = (0, 0, \ldots, 0)$ \hspace*{\fill}\eqref{eq:2-3}
        \item $-\vec{x} = (-x_1, -x_2, \ldots, -x_n)$ \hspace*{\fill}\eqref{eq:2-4}
        \item $\lambda\vec{x} = (\lambda x_1, \lambda x_2, \ldots, \lambda x_n)$ \hspace*{\fill}\eqref{eq:2-5} - \eqref{eq:2-8}
    \end{enumerate}
\end{proof}

Geometrically, we may interprit $\R^1$ as a line, and $\R^2$ as a plane. However, notice that the converse is not always true. Due to \eqref{eq:2-3} the space has to contain $0$. Therefore, if the line or plane does not go through the origin, it is not a vector space.

\subsection{Scalar Product}
\begin{defi}[Geometric definition in $\R^2$, $\R^3$]
    % insert diagram of dot product
    Given two vectors $\vec{a}, \vec{b}$ we define
    \[
        \vec{a} \cdot \vec{b} = \abs{\vec{a}}\abs{\vec{b}}\cos{\theta}  
    \]
    where $\theta$ is the angle between the two vectors.
\end{defi}
From this, we can see that
\begin{enumerate}
    \item $\vec{a} \cdot \vec{b} = \vec{b} \cdot \vec{a}$ 
    \item $\vec{a} \cdot \vec{a} = 0 \Leftrightarrow \vec{a} = \vec{0}$
    \item $\vec{a} \cdot \vec{b} = 0 \Leftrightarrow$ $\vec{a}$ and $\vec{b}$ are orthogonal
    \[\numberthis\]
\end{enumerate}

\begin{defi}[Projection]
    % insert diagram of projection
    $\vec{b}^\perp$ is the projection of $\vec{b}$ onto $\vec{a}$. Specifically, 
    \begin{equation}
        \vec{b}^\perp = \abs{\vec{b}}\cos{\theta}\frac{\vec{a}}{\abs{\vec{a}}} = \frac{(\vec{b}\cdot\vec{a})\vec{a}}{(\vec{a}\cdot\vec{a})} = (\uvec{a} \cdot \vec{b})\uvec{a} 
    \end{equation}
\end{defi}

We will move to a most abstract defintion of a scalar product
\begin{defi}[Inner product]
    We define $\vec{x} \cdot \vec{y} = \braket{\vec{x}}{\vec{y}} \in \R$ as an \emph{inner product}, if it satisfies
    \begin{enumerate}[eqn]
        \item $\vec{x} \cdot \vec{y} = \vec{y} \cdot \vec{x}$ \hspace*{\fill}(Symmetry) \label{eq:2-11}
        \item $\vec{x}, \vec{y}, \vec{z} \in V, \quad, \lambda, \mu \in \R$ \hspace*{\fill}(Linearity in the second argument) \label{eq:2-12}
        \[
            \vec{x} \cdot (\lambda\vec{y} + \mu\vec{z}) = \lambda\vec{x}\vec{y} + \mu\vec{x}\vec{z}  
        \]
        \item $\vec{x} \cdot \vec{x} \geq 0 $ and $\vec{x} \cdot \vec{x} = 0 \Leftrightarrow \vec{x} = \vec{0}$ \hspace*{\fill}(Positive definite) \label{eq:2-13}
    \end{enumerate}

    Note that we can use \eqref{eq:2-11} together with \eqref{eq:2-12} to derive linearlity in the first argument,
    \begin{equation}
        (\lambda\vec{y} + \mu\vec{z}) \cdot \vec{x} = \lambda\vec{x}\vec{y} + \mu\vec{x}\vec{z} \label{eq:2-15}
    \end{equation}
    We will see later in the course that linearity in the first argument does not hold for vector spaces over the complex field.

    We can also define the \emph{norm} of a vector $\vec{a}$ as 
    \[
        (\vec{a} \cdot \vec{a})^\frac{1}{2} = \abs{\vec{a}} = \norm{\vec{a}}  
    \]
\end{defi}

\begin{remark}
    Using this abstracted definition of an inner product, we can define an inner product between functions. For example define for function $f$ and $g$
    \[
        \braket{f}{g} = \int_{0}^{1}{fg \ \diffd x}
    \]
    This definition satisfies \eqref{eq:2-11} and \eqref{eq:2-12} and we can check \eqref{eq:2-13}
    \[
        \braket{f}{f} = \int_{0}^{1}{f^2 \ \diffd x} \geq 0
    \]
    and further if 
    \[
        \int_{0}^{1}{f^2 \ \diffd x} = 0 \Rightarrow f = 0
    \]
\end{remark}

\subsection{Cauchy-Schwartz Inequality}
\begin{thm}[CS Inequality]
    Given two vectors $\vec{x}$, $\vec{y}$
    \begin{equation}
        \vec{x} \cdot \vec{y} \leq \abs{\vec{x}}\abs{\vec{y}} \quad \vec{x}, \vec{y} \in \R^n
    \end{equation}
    with equality if and only if $\vec{x} = 0$ or $\vec{y} = 0$ or $\vec{x} = \lambda \vec{y}$ for some $\lambda \in \R$.
\end{thm}
\begin{proof}
    Consider $\abs{\vec{x} - \lambda \vec{y}}^2 \geq 0$
    \begin{align*}
        &(\vec{x} - \lambda \vec{y}) \cdot (\vec{x} - \lambda \vec{y}) \geq 0 \\
        \Rightarrow& \lambda^2 \vec{y} \cdot \vec{y} - 2\lambda\vec{x}\vec{y} + \vec{x} \cdot \vec{x} \geq 0
    \end{align*}
    Suppose $\vec{x}, \vec{y} \neq \vec{0}$, we have a quadratic equation in $\lambda$. Therefore since the quadratic is greater than or equal to $0$, the discriminant must be less than or equal to $0$.
    \begin{align*}
        \Delta &= [2(\vec{x} \cdot \vec{y})]^2 - 4 (\vec{y} \cdot \vec{y})(\vec{x} \cdot \vec{x}) \leq 0 \\
        \Rightarrow& [2(\vec{x} \cdot \vec{y})]^2 \leq 4 (\vec{y} \cdot \vec{y})(\vec{x} \cdot \vec{x}) \\
        \Rightarrow& (\vec{x} \cdot \vec{y})^2 \leq (\vec{y} \cdot \vec{y})(\vec{x} \cdot \vec{x}) \\
        \Rightarrow& (\vec{x} \cdot \vec{y})^2 \leq \abs{\vec{x}}^2\abs{\vec{y}}^2
    \end{align*}
    Furthermore, we have equality when
    \[
        \abs{\vec{x} - \lambda \vec{y}}^2 = 0 \Leftrightarrow \vec{x} = \lambda \vec{y}
    \]
    Hence the result
\end{proof}

\begin{eg} $\vec{x} = (\alpha, \beta, \gamma), y = (1, 1, 1)$
    \begin{align*}
        \alpha + \beta + \gamma &\leq \sqrt{\alpha^2 + \beta^2 + \gamma^2} \sqrt{3} \tag{\ref{eq:2-15}} \\
        \alpha^2 + \beta^2 + \gamma^2 &\geq \alpha\beta + \beta\gamma + \gamma\alpha \tag{rearranging}
    \end{align*}
\end{eg}

Note that we must not use the geometric definition of the dot product to prove the CS inequality because $\cos{\theta}$ is not defined for higher dimensions. However, we can use CS to define an angle in $\R^n$
\[
    \theta = \arccos{\left(\frac{\vec{x}\vec{y}}{\abs{\vec{x}}\abs{\vec{y}}}\right)}    
\]
since CS gives us that the fraction is less than or equal to $1$.

%%%%%%%%%%%%%%%%%%%%%%
% Lecture 4: 13/10/2023
%%%%%%%%%%%%%%%%%%%%%%

A further consequence of \eqref{eq:2-15} is the triangle inequality \eqref{eq:2}. We will consider the following
\begin{align*}
    \abs{\vec{x} + \vec{y}}^2 &= (\vec{x} + \vec{y}) \cdot (\vec{x} + \vec{y}) \\
    &= \vec{x} \cdot \vec{x} + 2 \vec{y} \cdot \vec{x} + \vec{y} \cdot \vec{y} \tag{collect terms over the reals} \\
    &= \abs{\vec{x}}^2 + \abs{\vec{y}}^2 + 2 \vec{y} \cdot \vec{x} \\
    &\leq \abs{\vec{x}}^2 + \abs{\vec{y}}^2 + 2 \abs{\vec{y} \cdot \vec{x}} \\
    &\leq \abs{\vec{x}}^2 + \abs{\vec{y}^2 + 2\abs{\vec{x}}}\abs{\vec{y}} \tag{By CS}\\
    &= (\abs{\vec{x}} + \abs{\vec{y}})^2 
\end{align*}
We may then take the positive square root on both sides to arive at
\begin{equation}
    \abs{\vec{x} + \vec{y}} \leq \abs{\vec{x}} + \abs{\vec{y}}
\end{equation}

\subsection{Vector Product}
We will see later in the course that this extends to the wedge product in higher dimensions. For now, we will define the product in $\R^3$.
\begin{defi}[Vector Product]
    Consider $\vec{a}, \vec{b} \in \R^3$. Define the vector product
    \[
        \vec{a} \times \vec{b} = \abs{\vec{a}}\abs{\vec{b}}\sin{\theta}\uvec{n}
    \]
    Where the direction of $\uvec{n}$ is such that $\vec{a}$, $\vec{b}$ and $\uvec{n}$ for a \emph{right handed system}. That is, when I curl my right hand fingers in the direction of $\theta$, the thumb points in the direction of $\uvec{n}$.
\end{defi}

Here are some properties of the vector product
\begin{enumerate}
    \item $\vec{a} \times \vec{b} - = \vec{b} \times \vec{a}$ \hspace*{\fill}(Anti-symmetric)
    \item $\vec{a} \times \vec{a} = \vec{0}$
    \item If $\vec{a} \times \vec{b} = \vec{0} \Rightarrow \vec{a} = \mu\vec{b}, \quad \text{for} \mu \in \R$
    \item $\vec{a} \times (\lambda \vec{b}) = \lambda (\vec{a} \times \vec{b})$
    \item[] \hspace*{\fill}\ilnumberthis
\end{enumerate}

\begin{defi}[Vector Area]
    We define the \emph{vector area} between two vectors $a$ and $b$ to be
    \[  
        \text{Vector Area} \equiv \frac{1}{2}(\vec{a} \times \vec{b})
    \]

    We may consider as an example the regular area of  a triangle, which can be found
    \begin{align*}
        \text{A} &= \frac{1}{2}\abs{\pvec{OB}}\abs{\pvec{BN}} \\
        &=  \frac{1}{2}\abs{\pvec{OB}}\abs{\pvec{OB}}\sin{\theta} \\
        &= \frac{1}{2}\abs{\vec{a} \times \vec{b}}
    \end{align*}
    
    Note that to turn this into a vector area, we omit the absolute value and use the vector itself to represent the area. An interesting consequence of this is that although the surface area of a sphere is $4\pi r^2$ the it's vector area is $\vec{0}$!
\end{defi}

\subsection{Scalar Triple Product}

\begin{defi}[Scalar Triple Product]
    Given three vectors $\vec{a}, \vec{b}, \vec{c} \in \R^3$, the triple product is defined as
    \begin{equation}
        \vec{a} \cdot (\vec{b} \times \vec{c}) = (\vec{b} \times \vec{c}) \cdot a = -a \cdot (c \times b)
    \end{equation}
    clearly, this is a scalar. Geometrically, we can interprit the triple product as the volume of a parallelopiped formed by $a$, $b$ and $c$, where the three vectors form a right handed system.
    
    % insert diagram of parallelopiped
    \begin{align*}
        \text{Volume} \ &= \ \text{Base Area} \ \times \perp \text{height} \\
        &\Rightarrow \abs{\vec{b} \times \vec{c}} \abs{\vec{a}} \sin{\theta} \\
        &= \abs{\vec{b} \times \vec{a}}\abs{\vec{a}}\cos{\bkt{\frac{\pi}{2} - \theta}} \\
        &= \abs{(\vec{b} \times \vec{a}) \cdot \vec{a}}
    \end{align*}
    Furthermore, since we specified that the vectors form a right handed system,
    \[
        \vec{a} \cdot (\vec{b} \times \vec{c}) \geq 0
    \]
\end{defi}

\begin{notation}
    We may instead write the tripple product using square brackets as follows
    \begin{equation}
        [\vec{a}, \vec{b}, \vec{c}] = [\vec{b}, \vec{c}, \vec{a}] =  [\vec{c}, \vec{a}, \vec{b}]
    \end{equation}

    where all even permutations are the equal, (even permutation meaning permutations performing an even number of pariwise swaps)
    \begin{equation}
        [\vec{a}, \vec{b}, \vec{c}] = - [\vec{b}, \vec{a}, \vec{c}] = - [\vec{a}, \vec{c}, \vec{b}] = - [\vec{c}, \vec{b}, \vec{a}]
    \end{equation}
    and all odd permutations are negative of the triple product.
\end{notation}

\begin{remark}
    The triple product can tell us if the three vectors are coplanar. Specifically,
    \begin{equation}
        \vec{a}, \vec{b}, \vec{c} \ \text{coplanar} \Leftrightarrow [\vec{a}, \vec{b}, \vec{c}] = 0  
    \end{equation}
\end{remark}

\begin{cor}
    Any triple product where two of the vectors are the same is equal to $0$.
    \[
        [\vec{a}, \vec{a}, \vec{c}] = [\vec{a}, \vec{b}, \vec{b}] = \ldots = 0  
    \]
\end{cor}

\subsection{Spanning sets \& Basis}
We will build up from the smaller example of $\R^2$, to $\R^3$, before working in $\R^n$. In general, spanning sets together with linear independance gives us a basis.

\subsubsection*{Working in $\R^2$}
\begin{defi}[Spanning sets]
    Consider $2$ vectors $\vec{a}, \vec{b} \neq 0 \in \R^2$ with $\vec{a} \times \vec{b} \neq 0$ i.e; $\vec{a}$ and $\vec{b}$ are not parallel. Then, any other vector $\vec{r}$ in this plane can be written as 
    \begin{equation}
        \vec{r} = \lambda \vec{a} + \mu \vec{b}, \quad \text{for} \lambda, \mu \in \R \label{eq:2-22}
    \end{equation}
    Then we say that the set $\{\vec{a}, \vec{b}\}$ spans $\R^2$, and every other vector in the space can be written as a linear combination of $\vec{a}$ and $\vec{b}$.
\end{defi}

\begin{prop}
    \eqref{eq:2-22} produces a unique representation.
\end{prop}

\begin{proof}
    Suppose 
    \[
        r = \lambda \vec{a} + \mu \vec{b} = \lambda' \vec{a} + \mu' \vec{b}    
    \]
    a useful trick to employ it to apply cross product with one of the terms, allowing us to remove it. Taking $\times \vec{b}$ to both sides of the equation we get
    \[
        (\lambda - \lambda')\vec{a} \times \vec{b} = \vec{0}  
    \]
    finally, as we have defined $\vec{a}$ and $\vec{b}$ to be non parallel,
    \[
        \therefore \lambda - \lambda' = 0
    \]

    Similarly, we can multiply $\times \vec{a}$ to arrive at
    \[
        \mu - \mu' = 0  
    \]
    Hence \eqref{eq:2-22} is unique.
\end{proof}

\begin{defi}[Linear Independance]
    Two vectors $\vec{a}, \vec{b}$ are said to be \emph{linearly independant} if, $\forall \alpha, \beta \in \R$
    \begin{equation}
        \alpha \vec{a} + \beta \vec{b} = 0 \Rightarrow \alpha = \beta = 0 \label{eq:2-23}
    \end{equation}
    Later on, we will see that this is generalised into $\R^n$
\end{defi}

\begin{remark}
    $\vec{a}$ and $\vec{b}$ are linearly dependant in $\R^2$ if $\vec{a} \times \vec{b} \neq 0$. To see that it is the case, take
    \begin{align*}
        \vec{b} \times \eqref{eq:2-23} &\Rightarrow \alpha (\vec{b} \times \vec{a}) = 0 \\
        \vec{a} \times \eqref{eq:2-23} &\Rightarrow \beta (\vec{a} \times \vec{b}) = 0 \\
    \end{align*}
    In both cases, since $\vec{a} \times \vec{b} \neq 0$, $\alpha = \beta = 0$
\end{remark}

\begin{defi}[Basis]
    A set $\{\vec{a}, \vec{b}\}$ is a basis of $\R^n$ if it spans $\R^2$ and it is linearly independant.
\end{defi}

\begin{remark}
    Note that the number of dimensions of a space in fact follows from the size of the set of minimal basis vectors.
\end{remark}

\begin{eg}
    \begin{enumerate}
        \item $\{\uvec{i}, \uvec{j}\}$ is a basis of $\R^2$
        \item $\{\uvec{i}, \uvec{j}, \vec{e}\}$ where $\vec{e} = \uvec{i} + \uvec{j}$ certainly spans $\R^2$, but it does not form a basis since the three vector system is not linearly independant. To see take,
        \[
            \uvec{i} + \uvec{j} + \vec{e} = \vec{0}    
        \]
    \end{enumerate}
\end{eg}

\begin{prop}
    A basis is the smallest spanning set
\end{prop}

\begin{proof}
    See Example sheet exercise
\end{proof}

\subsubsection*{Working in $\R^3$}
\begin{defi}[Basis in $\R^3$]
    If $\vec{a}, \vec{b}, \vec{c}$ are non-coplanar, then $\{\vec{a}, \vec{b}, \vec{c}\}$ is a basis of $\R^3$. We can write any vector $r$ as 
    \begin{equation}
        \vec{r} = \lambda \vec{a} + \mu \vec{b} + \nu \vec{c}, \quad \text{for}, \lambda, \mu, \nu \in \R \label{eq:2-24}
    \end{equation}
\end{defi}
\begin{proof}
    To show that it is uniquely defined, we will proceed by construction of the variables. Using a similar trick, using the tripple product to get rid of terms as repeated vectors in the triple product go to zero. $\eqref{eq:2-24} \cdot (\vec{b} \times \vec{c})$
    \begin{align*}
        \vec{r} \cdot (\vec{b} \times \vec{c}) &= \lambda \vec{a} \cdot (\vec{b} \times \vec{c}) \\
        &\Rightarrow \lambda = \frac{\vec{r} \cdot (\vec{b} \times \vec{c})}{\vec{a} \cdot (\vec{b} \times \vec{c})} = \frac{\vec{r} \cdot (\vec{b} \times \vec{c})}{[\vec{a}, \vec{b}, \vec{c}]}\\
        \vec{r} \cdot (\vec{c} \times \vec{a}) &= \mu \vec{b} \cdot (\vec{c} \times \vec{a}) \\
        &\Rightarrow \mu = \frac{\vec{r} \cdot (\vec{c} \times \vec{a})}{\vec{b} \cdot (\vec{c} \times \vec{a})}  = \frac{\vec{r} \cdot (\vec{c} \times \vec{a})}{[\vec{a}, \vec{b}, \vec{c}]} \\
        \vec{r} \cdot (\vec{a} \times \vec{b}) &= \nu \vec{c} \cdot (\vec{a} \times \vec{b}) \\
        &\Rightarrow \nu = \frac{\vec{r} \cdot (\vec{a} \times \vec{b})}{\vec{c} \cdot (\vec{a} \times \vec{b})}  = \frac{\vec{r} \cdot (\vec{a} \times \vec{b})}{[\vec{a}, \vec{b}, \vec{c}]} \\
    \end{align*}
    Hence, $\lambda, \mu, \nu$ can be determined for any $\vec{r}$ and $\{\vec{a}, \vec{b}, \vec{c}\}$ spans $\R^3$.
%%%%%%%%%%%%%%%%%%%%%%
% Lecture 5: 16/10/2023
%%%%%%%%%%%%%%%%%%%%%%
    Further to show linear independance, let 
    \[
        \vec{r} = \alpha \vec{a} + \beta \vec{b} + \gamma \vec{c} \quad \text{for} \ [\vec{a}, \vec{b}, \vec{c}] \neq 0
    \]
    We apply the same trick, using the triple product
    \begin{align*}
        \vec{r} \cdot (\vec{b} \times \vec{c}) &\Rightarrow \alpha [\vec{a}, \vec{b}, \vec{c}] = 0 \\
        &\Rightarrow \alpha = 0 \\
        \vec{r} \cdot (\vec{c} \times \vec{a}) &\Rightarrow \beta [\vec{a}, \vec{b}, \vec{c}] = 0 \\
        &\Rightarrow \beta = 0 \\
        \vec{r} \cdot (\vec{a} \times \vec{b}) &\Rightarrow \gamma [\vec{a}, \vec{b}, \vec{c}] = 0 \\
        &\Rightarrow \gamma = 0
    \end{align*}
    Hence, $\{\vec{a}, \vec{b}, \vec{c}\}$ form a basis of $\R^3$
\end{proof}

\begin{remark}
    Note that $\vec{a}, \vec{b}, \vec{c}$ need not be unique! As a standard basis in $\R^3$ we use,
    \[
        \{\uvec{i}, \uvec{j}, \uvec{k}\}
    \]
    Where they are
    \begin{enumerate}
        \item \[
            \uvec{i} \cdot \uvec{j} = \uvec{j} \cdot \uvec{k} = \uvec{k} \cdot \uvec{i} = 0  
        \] \hspace*{\fill}(orthogonal)
        \item \[
            \uvec{i} \cdot \uvec{i} = \uvec{j} \cdot \uvec{j} = \uvec{k} \cdot \uvec{k} = 1 
        \] \hspace*{\fill}(unit vectors)
    \end{enumerate}
    Together, we say that these form an \emph{orthonormal} basis.

    Another important point to keep in mind is that the choice of axis is highly arbitrary. Therefore, when we apply vectors to physical laws, we should expect the laws to be invariant with respect to changes of axis. We will cover more about this later in the course when we discuss tensors.
\end{remark}

\subsubsection*{Working in $\R^n$}
\begin{defi}[Linear independance in $\R^n$]
    A set of vectors $\{\vec{v_1}, \vec{v_2}, \cdots, \vec{v_n}\}$ where $v_i \in \R^n$ is \emph{linearly independant} if for scalars $\lambda_i \in \R$
    \begin{equation}\label{eq:2-25}
        \sum_{j=1}^{m}{\lambda_j \vec{v_j}} = \vec{0} \Leftrightarrow \lambda_j = 0 \quad \forall j = 1, 2, \cdots, n
    \end{equation}
\end{defi}

\begin{defi}[Spanning sets in $\R^n$]
    A set of vectors $\{\vec{u_1}, \vec{u_2}, \cdots, \vec{u_m}\}$ where $u_i \in \R^n$ is a \emph{spanning set} if for any $\vec{x} \in \R^n, \quad \exists \lambda_j$ such that,
    \begin{equation}\label{eq:2-26}
        \vec{x} = \sum_{j=1}^{m}{\lambda_j \vec{u_j}}
    \end{equation}
\end{defi}

\begin{remark}
    Note that we cannot span $\R^n$ if $m < n$.
\end{remark}

\begin{defi}[Basis in $\R^n$]
    A linearly independant set of vectors which spans $\R^n$ is a \emph{basis} of $\R^n$. A standard basis that we use
    \[
        \{\vec{e_1}, \vec{e_2}, \cdots, \vec{e_n}\}
    \]
    where
    \[
        e_i = \{00 \cdots 1 \cdots 00\}    
    \]
    where $1$ is in the $i^{th}$ position.
\end{defi}

\begin{defi}[Dimension]
    The \emph{dimension} of a vector space is the number of vectors in the basis.
\end{defi}

\begin{remark}
    We can prove that adding a vector to a basis does not result in a basis, and that if 2 basis exist, then they have the same number of elements. Therefore the dimension of a vector space is well defined.
\end{remark}

We can also define an inner product on this space, for two vectors $\vec{x}, \vec{y} \in \R^n$,
\begin{align*}
    \vec{x} = (x_1, \cdots, x_n) = \sum_{j=1}^{n}{x_i \vec{e_i}} \\
    \vec{y} = (y_1, \cdots, y_n) = \sum_{j=1}^{n}{y_i \vec{e_i}}
\end{align*}
Then the inner product of these two vectors is
\begin{equation}
    \vec{x} \cdot \vec{y} = \braket{x}{y} = \sum_{j=1}^{n}{x_iy_i}
\end{equation}

Note that for this we use the fact that $e_j$ form an orthonormal basis, such that
\[
    e_i \cdot e_j = \begin{cases}
        1 & \text{if} \ i = j \\
        0 & \text{otherwise}
    \end{cases}
\]

\subsubsection*{Working in $\C^n$}

We define the vector space $\C^n$
\[
    \C^n = \{(z_1, z_2, \cdots, z_n) \mid z_i \in \C\}  
\]
as well as a standard basis, which are still
\[
    (1, 0, 0, 0, \cdots), \quad (0, 1, 0, 0, \cdots), \quad (0, 0, 1, 0, \cdots)
\]

\begin{defi}[Scalar product in $\C^n$]
    In $\C^n$ the scalar product is defined, for $\vec{u}, \vec{v} \in \C^n$,
    \[
        \vec{u} \cdot \vec{v} = \braket{u}{v} = \sum_{i=1}^{n}\overline{u_i}v_i \in \C 
    \]
    Note the conjugation of the first argument. We have the following properties
    \begin{enumerate}[eqn]
        \item $\vec{u} \cdot \vec{v} = \overline{\vec{v} \cdot \vec{u}}$ \hspace*{\fill}(Conjugte symmetry)
        \item $\vec{u} \cdot (\lambda \vec{v} + \mu \vec{w}) = \lambda \vec{u} \cdot \vec{v} + \mu \vec{u} \cdot \vec{w}$ \label{eq:2-29}\hspace*{\fill}(Linearity in the $2^{nd}$ argument)
        \item $\vec{u} \cdot \vec{u} = 0 \Leftrightarrow \vec{u} = 0$
    \end{enumerate}

    To solve for the linearity of the first argument we consider
    \begin{align*}
        (\lambda \vec{u} + \mu \vec{v}) \cdot \vec{w} &= \overline{\vec{w} \cdot (\lambda \vec{u} + \mu \vec{v})} \\
        &= \overline{\sqbkt{\lambda \vec{w} \cdot \vec{u} + \mu \vec{w}\vec{v}}} \\
        &= \overline{\lambda} \overline{\vec{w} \cdot \vec{u}} + \overline{\mu} \overline{\vec{w} \cdot \vec{v}} \\
        &= \overline{\lambda} \vec{u} \cdot \vec{w} + \overline{\mu} \vec{v} \cdot \vec{w} \numberthis \label{eq:2-31}
    \end{align*}
\end{defi}

\begin{eg}
    To practice this we will try and expand for $\vec{x}, \vec{y} \in \C^n$   
    \begin{align*}
        &\sum_{k=1}^{4}{(-i)^k \abs{\vec{x} + i^k \vec{y}}} \\
        &= \sum_{k=1}^{4}{(-i)^k \braket{\vec{x} + i^k \vec{y}}{\vec{x} + i^k \vec{y}}} \\
        &= \sum_{k=1}^{4}{(-i)^k \crbkt{\braket{\vec{x} + i^k \vec{y}}{\vec{x}} + i^k\braket{\vec{x} + i^k \vec{y}}{\vec{y}}}} \tag{\ref{eq:2-29}}\\
        &= \sum_{k=1}^{4}{(-i)^k \crbkt{\braket{\vec{x}}{\vec{x}} + \overline{i^k}\braket{\vec{y}}{\vec{x}} + i^k\braket{\vec{x}}{\vec{y}} + i^k \overline{i^k}\braket{\vec{y}}{\vec{y}}}} \tag{\ref{eq:2-31}}\\
        &= \braket{\vec{x}}{\vec{x}}\sum_{k=1}^{4}{(-i)^k}  + \braket{\vec{y}}{\vec{x}}\sum_{k=1}^{4}{(-i)^k\overline{i^k}} + \braket{\vec{x}}{\vec{y}}\sum_{k=1}^{4}{(-i)^ki^k} + \braket{\vec{y}}{\vec{y}}\sum_{k=1}^{4}{(-i)^k i^k \overline{i^k}} \\
        &= \braket{\vec{x}}{\vec{x}}\sum_{k=1}^{4}{(-i)^k}  + \braket{\vec{y}}{\vec{x}}\sum_{k=1}^{4}{(-1)^k} + \braket{\vec{x}}{\vec{y}}\sum_{k=1}^{4}{(1)^k} + \braket{\vec{y}}{\vec{y}}\sum_{k=1}^{4}{(-i)^k}\\
        &= 4\braket{\vec{x}}{\vec{y}}
    \end{align*}
\end{eg}

\subsection{Vector subspaces}
\begin{defi}[Vector subspace]
    A non-empty subset $u$ of vectors of a vector space $V$ is a subspace of $V$ if $u$ is itself a vector space under the same operations as $V$.
\end{defi}

\begin{remark}
    Note that a subspace always includes the additive identity $0$. For instance, a plane is a subspace of 3D space only if it contains the origin.
\end{remark}

\begin{eg}[Trivial examples]
    $V$ and $\{0\}$ are subsapces of $V$. All others are called propper subspaces.
\end{eg}

\begin{thm}
    A subset $u$ of a vector space $V$ is a subspace of $V$ if and only if under the same operations as $V$,
    \begin{enumerate}
        \item $\vec{x}, \vec{y} \in u \Rightarrow \vec{x} + \vec{y} \in u$ \hspace*{\fill}(Closed under addition)
        \item $\vec{x} \in u \Rightarrow \lambda\vec{x}\in u$ for scalar $\lambda$ \hspace*{\fill}(Closed under scalar multiplication)
        \item $\vec{0} \in u$
    \end{enumerate}
    We can condense these three conditions into a single statement. $u$ is a subspace if and only if
    \begin{equation}\label{eq:2-32}
        \lambda \vec{x} + \mu \vec{y} \in u \quad \forall \vec{x}, \vec{y} \in u  
    \end{equation}
    To verify (i) we set $\lambda = \mu = 1$, to verify (ii) we set $\mu = 0$, to verify (iii) we set $\lambda = \mu = 0$
\end{thm}

\begin{eg}\leavevmode
    \begin{enumerate}
    \item Given $\{\vec{a}, \vec{b}, \vec{c}\}$ as a basis of a 3D space, $\{\vec{a} + \vec{c}, \vec{b} + \vec{c}\}$ is a basis of a 2D space.
    Suppose, $\vec{x}, \vec{y} \in \spnset{\vec{a} + \vec{c}, \vec{b} + \vec{c}}$. Then,
    \begin{align*}
        \vec{x} = \alpha_1 (\vec{a} + \vec{c}) + \beta_1 (\vec{b} + \vec{c}) \\
        \vec{y} = \alpha_2 (\vec{a} + \vec{c}) + \beta_2 (\vec{b} + \vec{c})
    \end{align*}
    Using \eqref{eq:2-32} we get,
    \begin{align*}
        \lambda\vec{x} + \mu \vec{y} =& \lambda\alpha_1 (\vec{a} + \vec{c}) + \lambda\beta_1 (\vec{b} + \vec{c}) \\
        &+ \mu\alpha_2 (\vec{a} + \vec{c}) + \mu\beta_2 (\vec{b} + \vec{c}) \\
        =& (\lambda\alpha_1 + \mu\alpha_2)(\vec{a} + \vec{c}) + (\lambda\beta_1 + \mu\beta_2)(\vec{b} + \vec{c}) \\
        & \in \spnset{\vec{a} + \vec{c}, \vec{b} + \vec{c}} \quad \therefore \text{subspace}
    \end{align*}
    To show linear independance,
    \begin{align*}
        \alpha (\vec{a} + \vec{c}) + \beta (\vec{b} + \vec{c}) &= 0 \\
        \Rightarrow \alpha\vec{a} + \beta \vec{b} + (\alpha + \beta)\vec{c} &= 0
    \end{align*}
    However since we chose $\{\vec{a}, \vec{b}, \vec{c} \}$ to be a 3D basis, we have linear independance of those vectors,
    \[
        \Rightarrow \alpha = \beta = \alpha + \beta = 0  
    \]
%%%%%%%%%%%%%%%%%%%%%%
% Lecture 6: 18/10/2023
%%%%%%%%%%%%%%%%%%%%%%
    \item $u = \{ \vec{x} \in \R^n \mid \sum_{i=1}^{n}{\alpha_ix_i} = 0, \ \alpha_i \in \R\}$ forms a subspace in $\R^n$.
    Consider $\vec{x}, \vec{y} \in u$.
    \begin{align*}
        \lambda \vec{x} + \mu \vec{y} &= \sum_{i=1}^{n}{\alpha_i(\lambda x_i + \mu y_i)} \\
        &= \lambda\sum_{i=1}^{n}{\alpha_ix_i} + \mu\sum_{i=1}^{n}{y_i} \\
        &= 0
    \end{align*}
    Therefore, $\lambda \vec{x} + \mu \vec{y} \in u$ and $u$ forms a subsapce of $\R^n$. We can think of this subspace as being $n-1$ dimensional. 
    If we think about the restriction imposed by the sum, we are able to freely choose $n-1$ variables before the $n^{th}$ variable is fixed by the sum condition.
    Also if we consider $\alpha = (1, 0, 0, 0, \cdots) = \vec{e_1}$. Then the equation for $u$ is $x_1 = 0$ after which we are free to change all other variables.
    Therefore a basis for this subspace could be $\{\vec{e_2}, \vec{e_3}, \cdots, \vec{e_n}\}$

    \item $w = \{ \vec{x} \in \R^n \mid \sum_{i=1}^{n}{\alpha_ix_i} = 1\}$ is \emph{NOT} a subspace.
    consdier the same linear combination,
    \begin{align*}
        \lambda \vec{x} + \mu \vec{y} &= \sum_{i=1}^{n}{\alpha_i(\lambda x_i + \mu y_i)} \\
        &= \lambda\sum_{i=1}^{n}{\alpha_ix_i} + \mu\sum_{i=1}^{n}{y_i} \\
        &= \lambda + \mu \\
        &\neq 1 \quad \forall \lambda, \mu \in \R
    \end{align*}
    This plane does not pass through the origin, and therefore is not a subspace.
    \end{enumerate}
\end{eg}

\subsection{Einstein Suffix Notation}

Consider a vector $\vec{v} \in \R^3$, we can write
\[
    \vec{v} = v_1\vec{e_1} + v_2\vec{e_2} + v_3\vec{e_3} = (v_1, v_2, v_3)
\]
In such notation, we write the $i^{th}$ component of $\vec{v}$ as $v_i$. Using Einstein notation, we will aim to write out vector equations in terms of these suffixes alone. For instance,
\begin{align*}
    \vec{a} = \vec{c} &\longrightarrow a_i = b_i \\
    \vec{c} = \alpha \vec{a} + \beta \vec{b} &\longrightarrow c_i = \alpha a_i + \beta b_i
\end{align*}
for the scalar product we may write
\[
    \vec{x} \cdot \vec{y} = \sum_{i=1}^{3}{x_iy_i} \longrightarrow x_iy_i
\]
In this case, the presence of a repeated \emph{free} suffix $i$ indicates that we sum over it.

\begin{defi}[Summation Convention]
    When we are writing sums, summation over free suffixes is understood.
    \[
        \vec{x} \cdot \vec{y} = x_iy_i
    \]
    Here $i$ is repeated and called a \emph{dummy} suffix. Since we are summing over it, it matters not what the dummy suffix is called.
\end{defi}

Here are the rules for using Einstein suffix notation.
\begin{enumerate}
    \item If the suffix appears once on the same term: it is a free suffix
    \item If the suffix appears twice in the same term: it is a dummy suffix and we sum over it
    \item If the suffix appears more than twice: this is wrong!
\end{enumerate}

\begin{eg}
    \begin{align*}
        [(\vec{a} \cdot \vec{b})\vec{c} - (\vec{a}\cdot\vec{c}\vec{b})]_i \tag{Note the single free suffix $i$}\\
        =(\vec{a} \cdot \vec{b})\vec{c_i} - (\vec{a}\cdot\vec{c})\vec{b_i} \\
        =a_kb_kc_i - a_kc_kb_i \tag{Note the use of dummy $k$}
    \end{align*}
    Note that for any equation the free suffixes must balance on the left and right. Here we see that we have one free $i$ on the left and right.
\end{eg}

\begin{defi}[Kronecker Delta]
    \[
        \delta_{ij} = \begin{cases}
            1 & \text{if} \ i = j \\
            0 & \text{otherwise}
        \end{cases}
    \]
    It is a symmetric symbol
    \[
        \delta_{ij} = \delta_{ji}
    \]
\end{defi}

\begin{remark}
    Here are some results from using the Kronecker Delta, 
    \begin{align*}
        a_i\delta_{i1} = a_1
    \end{align*}
    In this example we are summing over dummy $i$, however, all terms will be $0$ except for when $i = 1$. In this case the Kronecker delta, pulled out $1$ from the sum. In general,
    \begin{equation}\label{eq:2-33}
        a_i\delta_{ij} = a_j 
    \end{equation}
    We also have
    \begin{equation}\label{eq:2-34}
        \delta_{ij}\delta_{jk} = \delta_{ik}
    \end{equation}
    Where in this example, we have summed the product over all $j$. Note that this sum is only $1$ if $k = i$
    We can also sum over the Kronecker Delta itself,
    \begin{equation}
        \delta_{ii} = 3  
    \end{equation}

    We may compute an example,
    \[
        a_p \delta_{pq} b_q = a \cdot b
    \]
\end{remark}

\begin{defi}[Levi - Cevita (Alternating) Symbol]
    Consider the arrangements of $\{1, 2, 3\}$. We can classify them into two groups
    \[
        \{1, 2, 3\} \quad \{2, 3, 1\} \quad \{3, 1, 2\}
    \]
    where we perform an even number of swaps, and 
    \[
        \{2, 1, 3\} \quad \{1, 3, 2\} \quad \{3, 2, 1\}
    \]
    where we perform an odd number of swaps. These are known as even and odd permutations respectively. The Levi-Cevita symbol is:
    \[
        \varepsilon_{ijk} = \begin{cases}
            +1 & \text{if} \ i, j, k \ \text{is an even permutation of} \ 1, 2, 3 \\
            -1 & \text{if} \ i, j, k \ \text{is an odd permutation of} \ 1, 2, 3 \\
            0 & \text{otherwise}
        \end{cases}
    \]

    We note that $\varepsilon_{ijk}$ has three fixed suffixes.
\end{defi}

\begin{remark}\leavevmode
    \begin{enumerate}
        \item $\varepsilon_{ijk}\delta_{jk} = 0$ since the Kronecker delta is forcing two of the suffixes to be the same, in which case the Levi-Cevita symbol will be $0$.
        \item If $a_{jk} = a_{kj}$ then we say that $a_{jk}$ is symmetric. We can show that
        \[
            \varepsilon_{ijk}a_{jk} = \varepsilon_{ijk}a_{kj} = - \varepsilon_{ikj}a_{kj}
        \]
        Note that we first used the fact that $a_{jk}$ is symmetric, then the fact that when we swap the suffices of $\varepsilon$ once, we move from and even permutation to an odd and vice versa. Lastly, we recall that the names of dummy variables don't matter! Therefore, we get that
        \[
            \varepsilon_{ijk}a_{jk} = 0
        \]
        \item Recall the vector product, 
        \begin{align*}
            \vec{a} \times \vec{b} &= (a_1\uvec{i} + a_2\uvec{j} + a_1\uvec{k}) \times (b_1\uvec{i} + b_2\uvec{j} + b_1\uvec{k}) \\
            &= (a_2b_3 - a_3b_2)\uvec{i} + \cdots \\
            &= \begin{vmatrix}
                \uvec{i} & \uvec{j} & \uvec{k} \\
                a_1 & a_2 & a_3 \\
                b_1 & b_2 & b_3
            \end{vmatrix}
        \end{align*}
        We consider
        \[
            \varepsilon_{1jk}a_jb_k = \sum_{j=1}^{3}\sum_{k=1}^{3} \varepsilon_{1jk}a_jb_k 
        \]
        In this sum, the only non-zero terms are $j=2, k=3$ or $j=3, k=2$. Therefore we get,
        \begin{align*}
            = \varepsilon_{123}a_2b_3 + \varepsilon_{132}a_3b_2 \\
            = a_2b_3 - a_3b_2
        \end{align*}
        In general,
        \begin{equation}
            (\vec{a} \times \vec{b})_i = \varepsilon_{ijk}a_jb_k
        \end{equation}
    \end{enumerate}
\end{remark}

%%%%%%%%%%%%%%%%%%%%%%
% Lecture 7: 20/10/2023
%%%%%%%%%%%%%%%%%%%%%%

\begin{remark}[Important Identity]
    \begin{equation}
        \varepsilon_{ijk} \varepsilon_{ipq} = \delta_{jp} \delta_{kq} - \delta_{jq} \delta_{kp}   
    \end{equation}
    We can remember this as first pair, second pair, then mix.
\end{remark}

\begin{proof}
    We will proceed by exhaustion. On the right hand side,
    \[
    = \begin{cases} 
    +1 & \text{if} \ j = p \ \text{and} \ k = q \\
    -1 & \text{if} \ j = q \ \text{and} \ k = p \\
    0 & \text{otherwise}
    \end{cases}    
    \]
    On the left hand side, summing over $i$ the only non-zero times are when $j, k \neq i$ and $p, q \neq i$. So for $i = 1, 2, 3$
    \begin{align*}
        \text{if} \ j = p \ \text{and} \ k = q, & \ \text{LHS} \ = (-1)^2 \ \text{or} \ (1)^2 = 1 \\
        \text{if} \ j = q \ \text{and} \ k = p, & \ \text{LHS} \ = -1 \times 1 \ \text{or} \ 1 \times -1 = -1 \\
        \text{otherwise}, \qquad & \ \text{LHS} \ = 0
    \end{align*}
    Equally, 
    \begin{align*}
        \varepsilon_{ijk} \varepsilon_{pqk} &= \delta_{ip} \delta_{jq} - \delta_{iq} \delta_{jp} \\
        \varepsilon_{ijk} \varepsilon_{pjq} &= \delta_{ip} \delta_{kq} - \delta_{iq} \delta_{kp}   
    \end{align*}
\end{proof}

Using this notation we can now simply the scalar triple product,
\[
    [\vec{a}, \vec{b}, \vec{c}] = a_i \cdot (\vec{b} \times \vec{c})_i = \varepsilon_{ijk}a_ib_jc_k
\]
As well as the vector triple product
\begin{align*}
    [\vec{a} \times (\vec{b} \times \vec{c})] &= \varepsilon_{ijk}a_j(\vec{b} \times \vec{c})_k \\
    &= \varepsilon_{ijk}a_j \varepsilon_{kpq}b_pc_q \\
    &= \varepsilon_{ijk}\varepsilon_{pqk}a_j b_pc_q \\
    &= [\delta_{ip} \delta_{jq} - \delta_{iq} \delta_{jp}]a_j b_pc_q \\
    &= a_jb_ic_j - a_jb_jc_i \\
    &= (\vec{a} \cdot \vec{c})b_i - (\vec{a} \cdot \vec{b}) c_i
\end{align*}
Therefore we have
\begin{equation}
    \vec{a} \times (\vec{b} \times \vec{c}) = (\vec{a} \cdot \vec{c})\vec{b} - (\vec{a} \cdot \vec{b})\vec{c}
\end{equation}
and 
\begin{equation}
    \vec{a} \times (\vec{b} \times \vec{c}) \neq (\vec{a} \times \vec{b}) \times \vec{c}
\end{equation}

\begin{eg}[Past Tripos question - Spherical Triangle]\leavevmode
    \begin{enumerate}[label=(\alph*)]
        \item Prove that 
        \begin{equation}
            (a \times b) \cdot (a \times c) = (a \cdot a)(b \cdot c) - (a \cdot b)(a \cdot c)
        \end{equation}
        \begin{align*}
            \varepsilon_{ijk}a_jb_k \varepsilon{ipq}a_pc_q &= \varepsilon_{ijk}\varepsilon{ipq}a_jb_ka_pc_q \\
            &= (\delta_{jp} \delta_{kq} - \delta_{jq} \delta(kp))a_jb_ka_pc_q \\
            &= a_p b_q a_p c_q  - a_q b_p a_p c_q \\
            &= (a \cdot a)(b \cdot c) - (a \cdot c)(b \cdot a)
        \end{align*}
        \item Consider the unit sphere with center $O$, points $\vec{a}, \vec{b}, \vec{c}$ on the surface.
        Therefore $\abs{a} = \abs{b} = \abs{c} = 1$.
        Consider the greate circle containing $O$, $A$ and $B$.
        Let the arc length be $\delta(A, B)$
        \[
            \therefore a \cdot b = \cos(\delta(A, B))
        \]
        Now, the angle $\alpha$ in the spherical triangle is the angle between the planes containing $OAB$ and $OAC$.
        Therefore
        \begin{align*}
            \cos \alpha = \frac{a \times b}{\abs{a \times b}} \times \frac{a \times c}{\abs{a \times c}} \\
            \Rightarrow & \cos \alpha = \frac{(a \cdot a)(b \cdot c) - (a \cdot b)(a \cdot c)}{\abs{a \times b} \abs{a \times c}}
        \end{align*}
        Solving,
        \begin{align*}
            \abs{a \times c} = \abs{a} \abs{c} \sin(\delta(A, C)) \\
            \abs{a \times b} = \abs{a} \abs{b} \sin(\delta(A, B))
        \end{align*}
        \begin{align*}
            \Rightarrow \cos \alpha &= \frac{(b \cdot c) - (a \cdot b)(a \cdot c)}{\sin(\delta(A, C)) \sin(\delta(A, B))} \\
            &= \frac{\cos(\delta(B, C)) - \cos(\delta(A, B))\cos(\delta(A, C))}{\sin(\delta(A, C)) \sin(\delta(A, B))}
        \end{align*}
        \item Now consider equilateral spherical triangles, with side length $\delta$ and angle $\alpha$.
        \begin{align*}
            \cos \alpha \sin^2 \delta &= cos \delta - cos^2 \delta \\
            \cos \alpha &= \frac{cos \delta - cos^2 \delta}{\sin^2 \delta} \\
            &= \frac{cos \delta - cos^2 \delta}{1 - \cos^2 \delta} \\
            &= \frac{cos \delta}{1 + \cos \delta} \\
            &= 1 - \frac{1}{1 + \cos \delta}
        \end{align*}
        We have that $\cos \delta \leq 1$
        \[
            \Rightarrow \cos \alpha \leq \frac{1}{2} \Rightarrow \alpha \geq \frac{\pi}{3} 
        \]
        with equlity in the trivial case $\delta = 0$. Therefore for a non trivial equilateral triangle,
        \[
            \alpha \geq \frac{\pi}{3}  
        \]
    \end{enumerate}
\end{eg}

%%%%%%%%%%%%%%%%%%%%%%
% Lecture 8: 23/10/2023
%%%%%%%%%%%%%%%%%%%%%%

\subsection{Geometry}
\subsubsection{Lines}
A line $L$ through a point $A$ parallel to $\vec t$. For a general point $\vec x$ on $L$ is
\begin{subequations}
    \begin{equation}\label{eq:2-41a}
        \vec x = \vec a + \lambda \vec t \qquad \lambda \in \R
    \end{equation}
    If we take cross products we arrive at the alternative form
    \begin{equation}\label{eq:2-41b}
        (\vec x - \vec a) \times \vec t = \vec 0 
    \end{equation}
\end{subequations}
\subsubsection{Plane}
A plane with normal vector $\vec n$ containing the point $\vec b$. For a general point $\vec x$,
$\vec x - \vec b$ must lie in the plane. Therefore, $(\vec x - \vec b) \perp \vec n$.
\begin{subequations}
    \begin{equation}\label{eq:2-42a}
        (\vec x - \vec b) \cdot \vec n = 0 \Rightarrow \vec x \cdot \vec n = \vec b - \vec n
    \end{equation}\label{eq:2-42b}
    If $\abs{n} = 1$, $\vec b \cdot \uvec n$ is the perpendicular distance $p$ from the origin to the plane. Hence
    \begin{equation}
        \vec x \cdot \uvec n = p
    \end{equation}
\end{subequations}

Alternatively, if $\vec a, \vec b, \vec c$ lie on the plane, then (Sheet 3 Question 3)
\[
    (\vec x - \vec a) \cdot [(\vec b - \vec a) \times (\vec c - \vec a)] = 0
\]

\begin{eg}
    \begin{enumerate}[label=(\alph*)]
        \item Point of intersection of line $(\vec x - \vec a) \times \vec t = 0$ and plane $\vec x \cdot \vec n = \vec b \cdot \vec n$.
        Take $\times \vec n$ of \eqref{eq:2-41b}
        \begin{align*}
            &(\vec x \times \vec t) = (\vec a \times \vec t) \tag{$\times \vec n$} \\
            \Rightarrow& (\vec x \cdot \vec n)\vec t - (\vec t \cdot \vec n)\vec x = (\vec a \times \vec t) \times \vec n
        \end{align*}
        Eliminate $\vec x \cdot \vec n$ using \eqref{eq:2-42a}
        \[
            (\vec t \cdot \vec n)\vec x = (\vec b \cdot \vec n)\vec t - (\vec a \times \vec t)\times \vec n  
        \]
        We need to be careful that $\vec t \cdot \vec n \neq 0$ before dividing by it. 
        If $\vec t \cdot \vec n \neq 0$ then the line and the plane are not parallel,
        \[
            \therefore \vec x = \frac{\vec b \cdot \vec n}{\vec t \cdot \vec n} \vec t - \frac{(\vec a \times \vec t)\times \vec n}{\vec t \cdot \vec n}
        \]
        If $\vec t \cdot \vec n = 0$ then the line and the plane are parallel. There are two cases
        \begin{itemize}
            \item Disjoint and no solutions
            \item Coincident and all points on the line are solutions
        \end{itemize}
        In general, if we plug in values and arrive at a contradiction then there are no solutions, otherwise line and plane are coincident.

        \item Shortest distance between lines.
        \begin{align*}
            L_1 : (\vec x - \vec a_1) \times \vec t = 0 \\
            L_2 : (\vec x - \vec a_1) \times \vec t = 0
        \end{align*}
        The distance of closet approach will be along the line perpendicular to both lines, therefore the line will be $\parallel$ to $t_1 \times t_2$.
        % Insert diagram
        \begin{align*}
            S = \abs{a_1 - a_2} \cos \theta \\
            \Leftrightarrow S = \abs{\frac{(a_1 - a_2) \cdot (t_1 \times t_2)}{\abs{t_1 \times t_2}}}
        \end{align*}
        If $t_1$ and $t_2$ were parallel, then we would be dividing by $0$. Hence a necessary but not sufficient condition for lines to intersect
       \[
            (a_1 - a_2) \cdot (t_1 \times t_2) = 0
        \]
        The lines could be parallel and never intersect.
    \end{enumerate}
\end{eg}

\subsection{Vector Equations}
\begin{eg}
    Solve for $\vec x$: 
    \begin{equation}
        \vec x - (\vec x \times \vec a) \times \vec b = \vec c
    \end{equation}
    Our strategy is to take appropriate dot and cross products with suitable vectors.
    \[
        \vec x - (\vec x \cdot \vec b)\vec a + (\vec a \cdot \vec b) \vec x = \vec c \tag{$\star$}
    \]
    Take ($\star$) and dot with $\vec b$
    \begin{align*}
        (\vec x \cdot \vec b) - (\vec x \cdot \vec b)(\vec a \cdot \vec b) + (\vec x \cdot \vec b)(\vec a \cdot \vec b) = \vec c \cdot \vec b \\
        \Rightarrow \vec x \cdot \vec b = \vec c \cdot \vec b
    \end{align*}
    Then sub this back into ($\star$)
    \begin{align*}
        \vec x - (\vec c \cdot \vec b)\vec a + (\vec a \cdot \vec b) \vec x = \vec c \\
        \Rightarrow \vec x (1 + \vec a \cdot \vec b) = c + (\vec c \cdot \vec b)\vec a
    \end{align*}
    If $1 + \vec a \cdot \vec b \neq 0$ the
    \[
        \vec x = \frac{c + (\vec c \cdot \vec b)\vec a}{1 + \vec a \cdot \vec b} 
    \]
    otherwise, 
    \[
        1 + \vec a \cdot \vec b = 0 \Rightarrow c + (\vec c \cdot \vec b)\vec a = 0
    \]
    So if that does not hold, then there is a contradiction and there are no solutions. Otherwise
    \[
        c = - (\vec c \cdot \vec b)\vec a
    \]
    Then looking above, we have that
    \[
        \vec x \cdot \vec b = \vec c \cdot \vec b
    \]
    and we have a plane of solutions.
\end{eg}

\section{Matrices and Linear Maps}
\subsection{Examples}
\subsubsection{Rotation in $3$D}
Let the map $R: \R^3 \rightarrow \R^3$ represent a rotation of $\theta$ about the $z$-axis, such that $\vec x' = R(\vec x)$ is the rotated vector.
\begin{align*}
    \vec x &= (r \cos \phi, r \sin \phi, z) \\
    \vec x' &= (r \cos (\phi + \theta), r \sin (\phi + \theta), z) \\
    &= (x \cos \theta - y \sin \theta, x \sin \theta + y \cos \theta, z) 
\end{align*}
In suffix notation,
\[
    x'_i = R_{ij}x_j  
\]
where
\begin{equation}\label{eq:3-1}
    R = \begin{pmatrix}
        \cos \theta & -\sin \theta & 0 \\
        \sin \theta & \cos \theta & 0 \\
        0 & 0 & 1
    \end{pmatrix}
\end{equation}

\begin{remark}[Non-examinable]
    We will consider a rotation by $\theta$ about an arbitrary unit vector $\uvec n$.
    % Insert both diagrams

    Let's get $\vec x'$ in terms of $\vec x$. 
    \[
        \vec x' = \pvec OB + \pvec BC + \pvec CA'
    \]
    We have that
    \begin{align*}
        \pvec OB &= (\vec x \cdot \uvec n) \uvec n \\
        \pvec BC &= \pvec BA \cos \theta \\
        &= (\pvec BO + \pvec OA) \cos \theta \\
        &= (\vec x - (\vec x \cdot \uvec n)\uvec n) \cos \theta \\
        \abs{\pvec CA'} &= \abs{\pvec BA'} \sin \theta = \abs{\uvec n \times \vec x} \sin \theta \\
        \therefore \pvec CA' &= (\uvec n \times \vec x) \sin \theta \\
    \end{align*}
    Therefore,
    \[
        \vec x' = x \cos \theta + (1 - \cos \theta)(\uvec n \cdot \vec x)\uvec n + (\uvec n \times \vec x) \sin \theta 
    \]
    in suffix notation,
    \begin{align*}
        \vec x'_i &= x_i \cos \theta + (1 - \cos \theta)n_jx_jn_i - \varepsilon_{ijk} x_jn_k \sin \theta \\
        &= R_{ij}x_j
    \end{align*}
    where
    \begin{equation}\label{eq:3-2}
        R_{ij} = \delta_{ij}\cos \theta + (1 - \cos \theta)n_jn_i - \varepsilon_{ijk} n_k \sin \theta
    \end{equation}
\end{remark}

%%%%%%%%%%%%%%%%%%%%%%
% Lecture 9: 25/10/2023
%%%%%%%%%%%%%%%%%%%%%%

\subsubsection{Reflecetions in $\R^3$}
We reflect on a plane with normal $\uvec n$ through the origin .
\begin{align*}
    \vec x' &= \vec x - 2 (\vec x \cdot \uvec n) \uvec n \\
    \vec x'_i &= \vec x_i - 2 \vec x_j \vec n_j \vec n_i \\
    &= R_{ij}x_j \\
\end{align*}
where
\begin{equation}\label{eq:3-3}
    R_{ij} = \delta_{ij} - 2 n_i n_j
\end{equation}
We also have the projection of $\vec x$ onto the plane
\begin{equation}\label{eq:3-4}
    \vec x' = \vec x - (\vec x \cdot \uvec n) \uvec n
\end{equation}

\subsection{Linear Maps}
Consider sets $A$, $B$ an a map $T: A \rightarrow B$ such that $\vec x \in A$ is mapped to a unique $\vec x' = T(\vec x) \in B$. 
We say that
\begin{itemize}
    \item $A$ is the domain of $T$
    \item $B$ is the codomain of $T$
\end{itemize}
\begin{remark}[Image]
    $T(A)$ is the \emph{image} - the set of points of form $T(\vec x) \quad \forall \vec x \in A$.
    \[
        T(A) \subset B  
    \]
    but not necessarily the entire of $B$.
\end{remark}


\begin{defi}[Linear Map]
    Typically,
    \[
        T: R^n \rightarrow R^m \text{ or } T: \C^n \rightarrow \C^m
    \]
    Let $V, W$ be the real (complex) vector spaces and $T: V \rightarrow W$. 
    Then, $T$ is a \emph{Linear Map} if
    \begin{enumerate}
        \item $T(\vec a + \vec b) = T(\vec a) + T(\vec b) \quad \forall \vec a, \vec b \in V$
        \item $T(\lambda \vec a) = \lambda T(\vec a) \quad \forall \lambda \in \R \text{ or } (\C)$
    \end{enumerate}
    Or equivalently,
    \begin{equation}\label{eq:3-5}
        T(\lambda \vec a + \mu \vec b) = \lambda T(\vec a) + \mu T(\vec b)
    \end{equation}
\end{defi}
\begin{eg}
    \begin{enumerate}
        \item Translation
        \begin{align*}
            T: \R^3 &\rightarrow \R^3 \\
            \vec x &\mapsto \vec x + \vec a
        \end{align*}
        This is not a linear map since,
        \begin{align*}
            T(\lambda \vec x + \mu \vec y) &= \lambda \vec x + \mu \vec y + \vec a \\
            \lambda T(\vec x) + \mu T(\vec y) &= \lambda \vec x + \mu \vec y + (\lambda + \mu) \vec a \\
            \therefore&  T(\lambda \vec x + \mu \vec y) \neq \lambda T(\vec x) + \mu T(\vec y)
        \end{align*}
        \item We can check that reflections and rotations are linear maps.
        \item Define a map
        \begin{align*}
            S: \R^3 &\rightarrow \R^2 \\
                (x, y, z) &\mapsto (x + y, 2x - z)
        \end{align*}
        We can check that the map is linear
        \begin{align*}
            S(\lambda \vec x + \mu \vec y) &= (\lambda x_1 + \mu y_1 + \lambda x_2 + \mu y_2,  2(\lambda x_1 + \mu y_1) - (\lambda x_3 + \mu y_3)) \\
            &= \lambda (x_1 + x_2, 2x_1 - x_3) + \mu (y_1 + y_2, 2y_1 - y_3) \\
            &= \lambda S(\vec{x}) + \mu S(\vec y)
        \end{align*}
        Therefore the map is linear.
    \end{enumerate}
\end{eg}
\begin{remark}[Image and Kernel]
    Let's consider the effect of $S$ on the standard basis vectors in $\R^3$.
    \begin{align*}
        S(1, 0, 0) = (1, 2) \\
        S(0, 1, 0) = (1, 0) \\
        S(0, 0, 1) = (0, -1) \\
    \end{align*}
    These $3$ vectors are not linearly independant, but they certainly span $\R^2$.
    \[
        \therefore S(\R^3) = \R^2  
    \]
    and we say that the image of $\R^3$ under $S$ in $\R^2$.

    What if $S(x, y, z) = \vec 0$?
    \begin{align*}
        x + y = 0 \\
        2x - z = 0 \\
        \Rightarrow \vec x = (t, -t, 2t)
    \end{align*}
    Therefore any vector parallel to $(1, -1, 2)$
    The set of all $x \in \R^3$ which map to $\vec 0 \in \R^2$ is the \emph{Kernel} of $S$.
\end{remark}
\begin{defi}[Image \& Kernel]
    In general, for a map $f: U \rightarrow V$. 
    \begin{itemize}
        \item The \emph{image} of $f$ is the subset of $V$
        \[
            \im(f) = \crbkt{f(\vec u) \mid \forall \vec u \in U}  
        \]
        \item The \emph{kernel} of $f$ is the subset of $U$
        \[
            \ker(f) = \crbkt{\vec u \in U \mid f(\vec u) = \vec 0}  
        \]
    \end{itemize}
\end{defi}
We can see geometrically, that
\begin{enumerate}
    \item The image of the rotation map is the whole of $\R^3$ and the kernel is $\crbkt{0}$
    \item For the projection \eqref{eq:3-4}, the image is the plane through $\vec 0$ with normal $\vec n$, and the kernel is the line $\parallel$ to $\vec n$ through $\vec 0$, i.e;
    \[
        \spnset{\vec n}  
    \]
\end{enumerate}

\begin{thm}[Image \& Kernel are subspaces.]
    Consider a linear map $f: U \rightarrow V$ where $U, V$ are vector spaces.
    Then $\im (f)$ is a subspace of $V$ and $\ker (f)$ is a subspace of $U$.
\end{thm}
\begin{proof}
    \begin{subequations}
        By \eqref{eq:2-32}, if
        \[
            \vec x, \vec y \in \im(f) \Rightarrow (\exists \vec a, \vec b \in U \mid \vec x = f(\vec a) \text{ and } \vec y = f(\vec b))
        \]
        then,
        \begin{align*}
            \lambda \vec x + \mu \vec y &= \lambda f(\vec a) + \mu f(\vec b) \\
            &= f(\lambda \vec a + \mu \vec b) \tag{$f$ is linear}
        \end{align*}
        Since $U$ is a vector space, $\lambda \vec a + \mu \vec b \in U$ and  
        \begin{equation}
            \therefore \quad f(\lambda \vec x + \mu \vec y) \in \im (f)
        \end{equation}
        Therefore the image is a subspace of $V$.
    
        If $\vec x, \vec y \in \ker(f)$
        \begin{align*}
            f(\lambda \vec x + \mu \vec y) &= \lambda f(\vec x) + \mu f(\vec y) \\
            &= \vec 0 \in \ker(f)
        \end{align*}
        Therefore the linear combination lies in the kernel,
        \begin{equation}
            \lambda \vec x + \mu \vec y \in \ker(f)
        \end{equation}
        Therefore the kernel is a subspace of $U$.
    \end{subequations}
\end{proof}

\subsection{Rank and Nulity}
Let $f$ be a linear map $f: U \rightarrow V$. 

\begin{defi}[Rank and Nullity]
    The \emph{rank} of $f$ is denoted
    \begin{equation}\label{eq:3-7}
        r(f) = \dimset{f(U)}
    \end{equation}

    The \emph{nullity} of $f$ is denoted
    \begin{equation}\label{eq:3-8}
        n(f) = \dimset{\ker(f)}
    \end{equation}
\end{defi}
\begin{eg}
    For the projection \eqref{eq:3-4}
    \begin{itemize}
        \item Image is a plane $\therefore r(f) = 2$
        \item Kernel is the vector perpendicular to the plane $\therefore n(f) = 1$
    \end{itemize}
\end{eg}

\begin{thm}[Rank-Nullity]
    For a linear map $f: U \rightarrow V$ where $U, V$ are vector spaces,
    \begin{equation}\label{eq:3-9}
        r(f) + n(f) = \dim U
    \end{equation}
\end{thm}
\begin{proof}
    Write $\dim U = N$ and $n(f) = m$ with $m < N$. 
    If $m = N$, everything is mapped to $\vec 0$ and hence $f$ is the zero map.

    Suppose $(e_1, e_2, \cdots, e_m)$ is a basis for the $\ker(f)$. 
    We extend this basis to the whole of $U$, 
    \[
        \crbkt{e_1, e_2, \cdots, e_m, e_{m+1}, e_{m+2}, \cdots, e_{N}}  
    \]
    To prove the rank nullity theorem, we will show that
    \[
        \crbkt{f(e_{m+1}), f(e_{m+2}), \cdots, f(e_{N})}
    \]
    is a basis for $\im(f)$
    \begin{description}
        \item[Spanning]. The basis spans $\im(f)$.
        Take $y \in \im(f)$ this means that $\exists \vec x \in U$ such that
        \begin{align*}
            \vec y &= f(\vec x) \\
            &= f(\alpha_1 x_1 + \cdots \alpha_N e_N) \\
            &= \alpha_1f(x_1) + \cdots \alpha_N f(e_N) \\
            &= \alpha_{m+1}f(e_{m+1}) + \alpha_{m+2}f(e_{m+2}) + \cdots + \alpha_Nf(e_{N})
        \end{align*}
        Therefore
        \[
            \crbkt{f(e_{m+1}), f(e_{m+2}), \cdots, f(e_{N})}
        \]
        spans the image of $f$.
        \item[Linear indedpendance] Suppose
        \[
            \alpha_{m+1}f(e_{m+1}) + \alpha_{m+2}f(e_{m+2}) +\cdots, \alpha_Nf(e_{N}) = 0 \tag{$\star$}
        \]
        \begin{align*}
            \Rightarrow f(\alpha_{m+1}f(e_{m+1}), \alpha_{m+2}f(e_{m+2}), \cdots, \alpha_Nf(e_{N})) = 0 \\
            \therefore \alpha_{m+1}f(e_{m+1}), \alpha_{m+2}f(e_{m+2}), \cdots, \alpha_Nf(e_{N}) \in \ker(f) \\
            \Rightarrow \alpha_1 x_1 + \cdots + \alpha_N e_m = \alpha_{m+1}e_{m+1} + \cdots + \alpha_{N}e_{N}
        \end{align*}
        However, $e_1, e_2, \cdots, e_n$ is a basis for $U$ and are therefore linearly independant. 
        Therefore the only solution to this is
        \[
            \alpha_1 = \cdots = \alpha_m = \alpha_{m+1} = \cdots = \alpha_N = 0  
        \]
        Now back to ($\star$), since all $\alpha$'s are $0$, 
        \[
            \crbkt{f(e_{m+1}), f(e_{m+2}), \cdots, f(e_{N})}    
        \]
        are linearly independant.
    \end{description}
\end{proof}

%%%%%%%%%%%%%%%%%%%%%%
% Lecture 10: 27/10/2023
%%%%%%%%%%%%%%%%%%%%%%

\begin{eg}
    Calculate the kernel and the image of the map
    \begin{align*}
        f: \R^3 &\rightarrow R^3 \\
        (x, y, z) &\mapsto (x + y + z, 2x - y, x + 2z)
    \end{align*}

    \begin{enumerate}[steps]
        \item Find the kernel,
        \begin{equation*}
            \begin{aligned}
                x + y + z = 0 \\
                2x - y = 0 \\
                x + 2z = 0 \\
                \hline
                \Rightarrow x + 2z = 0
            \end{aligned}
        \end{equation*}
        Therefore, kernel is of the form $(-2z, z, z)$
        \[
            \ker(f) = \spnset{(-2, 1, 1)} \Rightarrow n(f) = 1  
        \]
        \item Extend basis of kernel to basis of the whole space
        \[
            \R^3 = \spnset{(-2, ,1, 1), (0, 1, 0), (0, 0, 1)}  
        \]
        \item Apply $f$ to this basis.
        \begin{align*}
            f(-2, ,1, 1) = \vec 0 \\
            f(0, 1, 0) = (1, -1, 0) \\
            f(0, 0, 1) = (1, 5, 2)
        \end{align*}
        \item Find a basis for $\im(f)$
        \[
            \im(f) = \spnset{(1, -1, 0), (1, 5, 2)} \therefore r(f) = 2
        \]
    \end{enumerate}

    Note that $r(f) + n(f) = 2 + 1 = 3$ agrees with the rank nullity theorem.
    Let's get the image as a plane.
    The plane contains $(1, -1, 0)$ and $(1, 5, 2)$. 
    \begin{align*}
        \vec n &= (1, -1, 0) \times (1, 5, 2) \\
        &= (1, 1, -3)\lambda
    \end{align*}
    We also know that $\im(f)$ contains the origin. Therefore the plane must pass through the origin,
    \begin{align*}
        \vec r \cdot (1, 1, -3) &= 0 \\
       \Rightarrow x + y - 3z &= 0
    \end{align*}
\end{eg}

\subsection{Matrices}
Consider a linear map (applies to $\C$) as well.
\begin{align*}
    \alpha: R^n &\rightarrow \R^m \\
    \vec x &\mapsto \vec x'
\end{align*}

For a basis of $\R^n$, $(e_1, e_2, \cdots, e_n)$, write
\begin{align*}
    \vec x &= \sum_{j=1}^{n}{x_je_j} \\
    \vec x' &= \alpha(\sum_{j=1}^{n}{x_je_j}) = \sum_{j=1}^{n}{x_j\alpha(e_j)} \tag{by \ref{eq:3-5}}\\
    x'_i &= \sum_{j=1}^{n}{x_j [\alpha(e_j)]_i} \qquad i = 1, 2, \cdots, m
\end{align*}
In suffix notation $A_{ij} = [\alpha(e_j)]_i$, 
\begin{equation}\label{eq:3-10}
    x'_j = A_{ij}x_j
\end{equation}
Where
\[
    A = \crbkt{A_{ij}} = \begin{pmatrix}
        A_{11} & \cdots & \cdots & \cdots & A_{1n} \\
        \vdots &        &        &        & \vdots \\
        \vdots &        & A_{ij} &        & \vdots \\
        \vdots &        &        &        & \vdots \\
        A_{m1} & \cdots & \cdots & \cdots & A_{mn} \\
    \end{pmatrix}_{m \times n}
\]
And $\vec x' = A \vec x$
\[
    \Rightarrow (\vec x')_{m \times 1} = (A)_{m \times n}(\vec x)_{n \times 1}    
\]
Each linear map $\alpha$ has an associated matric $A$ which is \underline{basis dependant}.

\subsubsection{Examples}
\begin{enumerate}[label=\arabic*.]
    \item Rotation of $\theta$ about $z$-axis \eqref{eq:3-1}
    \item Reflection in the plane \eqref{eq:3-3}
    \item Dilation 
    \begin{align*}
        \alpha: \R^3 &\rightarrow \R^3 \\
        (x, y, z) &\mapsto (\lambda x, \mu y, \gamma z)
    \end{align*}
    Has a dilation matrix
    \[
        \begin{pmatrix}
            \lambda & 0 & 0 \\
            0 & \mu & 0 \\
            0 & 0 & \gamma
        \end{pmatrix}
    \]
    \item Sheer in the $x$-axis
    \begin{align*}
        S: \R^3 &\rightarrow \R^3 \\
        (x, y, z) &\mapsto (x + \lambda y, y, z)
    \end{align*}
    Has a sheer matrix
    \[
        \begin{pmatrix} 
        x' \\ y' \\ z'
        \end{pmatrix}
        = 
        \begin{pmatrix}
            1 & \lambda & 0 \\
            0 & 1 & 0 \\
            0 & 0 & 1
        \end{pmatrix}
        \begin{pmatrix}
            x \\ y \\ z
        \end{pmatrix}
    \]
\end{enumerate}

\subsubsection{Matrix Algebra}
\begin{enumdescript}
    \item[Addition] Consider $2$ maps $\alpha, \beta: \R^n \rightarrow \R^m$.
    The sum of these maps 
    \[
        (\alpha + \beta)(\vec x) \equiv \alpha(\vec x) + \beta(\vec x)  
    \]
    We can check that the associated matrices do the same,
    \[
        (A+ B)_{ij}x_j = A_{ij}x_j + B_{ij}x_j
    \]
    \begin{equation}\label{eq:3-11}
        \Rightarrow (A + B)_{ij} = A_{ij} + B_{ij}
    \end{equation}
    \item[Scalar Multiplication] For $\lambda$ a scalar,
    \[
        (\lambda \alpha)\vec x = \lambda \alpha(\vec x)
    \]
    \begin{equation}\label{eq:3-12}
        \therefore (\lambda A)_{ij} = \lambda A_{ij}
    \end{equation}
    \item[Matrix Multiplication] Consider $2$ maps
    \begin{align*}
        \alpha: \R^l \rightarrow \R^n \\
        \beta: \R^n \rightarrow \R^m
    \end{align*}
    Where the associated matrices are
    \[
        A_{n \times l} = \crbkt{A_{ij}} \text{ and } B_{m \times n} = \crbkt{B_{ij}}
    \]
    The composition map, where we perform $\alpha$ followed by $\beta$ is
    \[
        \beta\alpha: \R^l \rightarrow \R^m  
    \]
    We will see that the definition for matrix multiplication comes natrually from the composition of maps.
    Take $\vec x \in \R^l \mapsto \vec x'' \in \R^m$
    \begin{align*}
        \vec x'' &= BA(\vec x) \\
        &= B\vec x'
    \end{align*}
    Where $\vec x' = A\vec x$. Then
    \begin{align*}
        x''_i = (B\vec x')_i &= B_{ik}x'_k \\
        &= B_{ik}A_{kj}x_j \\
        &= (BA)_{ij}x_j
    \end{align*}
    \begin{equation}\label{eq:3-13}
        (BA)_{ij} = B_{ik}A_{kj}
    \end{equation}
    \begin{remark}\leavevmode
        \begin{enumerate}
            \item Column number of $B$ $=$ row number of $A$. 
            If $l = m$ then we can form $BA$ and $AB$, but 
            \item $AB \neq BA$ not commutative in general
            \item $A(BC) = (AB)C$ only associative
        \end{enumerate}
    \end{remark}
    \item[Transposition] Let $A$ be an $m \times n$ matrix. Then $A^T$ is an $n \times m$ matrix defined by
    \[
        (A^T)_{ij} = A_{ji}  
    \]
    Clearly, $(A^T)^T = A$. We also have for column vectors,
    \[
        \vec x = \begin{pmatrix}
                x_1 \\ \vdots \\ x_n
        \end{pmatrix}
        \Rightarrow 
        \vec x^T = \begin{pmatrix}
            x_1 & \cdots & x_n
        \end{pmatrix}
    \]
    \begin{prop}
        \begin{equation}\label{eq:3-14}
            (AB)^T = B^TA^T
        \end{equation}
    \end{prop}
    \begin{proof}
        \begin{align*}
            (AB)^T_{ij} = (AB)_{ji} &= A_{jk}B_{ki} \tag{from \eqref{eq:3-13}} \\
            &= B_{ki}A_{jk} \\
            &= (B^T)_{ik}(A^T)_{kj} \\
            &= (B^TA^T)_{ij}
        \end{align*}
    \end{proof}

    \begin{defi}[Hermition Conjugate]
        \begin{equation}\label{eq:3-15}
            A^\dagger = (A^T)^* = (A^*)^T
        \end{equation}
        i.e; take the complex conjugate and then transpose.
    \end{defi}
    
    %%%%%%%%%%%%%%%%%%%%%%
    % Lecture 11: 30/10/2023
    %%%%%%%%%%%%%%%%%%%%%%
    \begin{defi}[Symmetric Matrix]
        A \emph{symmetric} matrix is such that 
        \[
            A^T = A
        \]
        We see that $A$ is neccesarily square.
    \end{defi}
    
    \begin{defi}[Hermition Matrix]
        A \emph{hermition} matrix is such that
        \[
            A^\dagger = A
        \]
        We see that the diagonal elements of $A$ are neccesarily real.
    \end{defi}
    
    \begin{defi}[Anti-Symmetric Matrix]
        A \emph{anti-symmetric} matrix is such That        
        \[
            A^T = -A
        \]
        We see that its diagonal elements are neccesarily $0$.
    \end{defi}
    
    \begin{defi}[Skew-Hermition Matrix]
        A \emph{skew-hermition} matrix is such that
        \[
            A^\dagger = -A
        \]
        We see that the diagonal elements of $A$ are neccesarily purely imaginary.
    \end{defi}

    \item[Trace] The trace of a matrix $A$ is the sum of its diagonal elements. 
    This turns out to be invariant with respect to changing basis.
    \[
        \tr(A) = A_{ii}
    \]
    \begin{eg}
        \[
            R_{ij} = \delta_{ij} - 2n_in_j \Rightarrow \tr(R) = R_{ii} = 3 -2 = 1  
        \]
    \end{eg}
    It has the commutative property
    \begin{equation}\label{eq:3-16}
        \tr(BC) = \tr(CB)
    \end{equation}
    \begin{proof}
        \begin{align*}
            \tr(BC) = (BC)_{ii} &= B_{ik}C_{ki} \tag{by \eqref{eq:3-13}}\\
            &= C_{ki}B_{ik} \\
            &= (CB)_{kk} = \tr(CB)\\
        \end{align*}
    \end{proof}

    \item[Identity Matrix]
    \[
        I_{ij} = \delta_{ij} = \begin{cases}
            1 & \text{ if } i = j \\
            0 & \text{ otherwise}
        \end{cases}
    \]

    \item[Interpretation of columns of a matrix]. Consider $\alpha: \R^n \rightarrow \R^n$, and use the standard basis.
    \[
        A \begin{pmatrix}
            1 \\ 0 \\ \vdots \\ 0
        \end{pmatrix} 
        =
        \begin{pmatrix}
            A_{11} \\ A_{21} \\ \vdots \\ A_{n1}
        \end{pmatrix}
    \]
    and in general, for the standard basis vector $e_i$
    \[
        A \begin{pmatrix}
            0 \\ \vdots \\ 1 \\ \vdots \\ 0
        \end{pmatrix} 
        =
        \begin{pmatrix}
            A_{1i} \\ A_{2i} \\ \vdots \\ A_{ni}
        \end{pmatrix}
    \]
    Therefore, the columns of the matrix are interpreted as the images of the basis vectors under $\alpha$.
    \begin{eg}
        Reflection in $\R^2$ on the line at an angle $\theta$ to $x$-axis.
        \begin{align*}
            \uvec i &\mapsto \uvec i' = \cos(2 \theta)\uvec i + \sin(2 \theta)\uvec j \\
            \uvec j &\mapsto \uvec j' = \sin(2 \theta)\uvec i - \cos(2 \theta)\uvec j
        \end{align*}
        Therefore, the matrix representing this map is given by
        \[
            A = \begin{pmatrix}
                \cos(2 \theta) & \sin(2 \theta) \\
                \sin(2 \theta) & -\cos(2 \theta)
            \end{pmatrix}  
        \]
    \end{eg}
\end{enumdescript}

\subsubsection{Decomposition of an $n \times n$ matrix}
Any $n \times n$ matrix can be written as a sum of symmetric and anti-symmetric parts. Write
\[
    B_{ij} = \frac{1}{2}(B_{ij} + B_{ji}) + \frac{1}{2}(B_{ij} - B_{ji})
\]
Where
\[
    S_{ij} = \frac{1}{2}(B_{ij} + B_{ji}) = S_{ji}
\]
and is symmetric, and
\[
    A_{ij} = \frac{1}{2}(B_{ij} - B_{ji}) = -A_{ji}
\]
is anti-symmetric. Therefore,
\begin{equation}\label{eq:3-17}
    B = \frac{1}{2}(B + B^T) + \frac{1}{2}(B - B^T)
\end{equation}
Furhermore, we can decompose $S$ into an Isotropic part $+$ Traceless part. Write
\[
    S_{ij} = \frac{1}{n}\tr(S)\delta_{ij} + (S_{ij} - \frac{1}{n}\tr(S)\delta_{ij})
\]
Where the first part is isotropic (a scalar multiple of $\delta_{ij}$) and the second half is traceless since
\begin{align*}
    T_{ij} = S_{ij} - \frac{1}{n}\tr(S)\delta_{ij} \\
    \Rightarrow T_{ii} = S_{ii} - \frac{1}{n}\tr(S)\delta_{ii} \\
    =\tr(s) - \frac{1}{n}\tr(s) \times n \\
    = 0
\end{align*}

Putting these all together
\begin{equation}\label{eq:3-18}
    B = \frac{1}{n}\tr(B)I + \crbkt{\frac{B + B^T}{2} - \frac{1}{n}\tr(B)I} + \frac{1}{2}(B - B^T)
\end{equation}

In three dimensions (only), the anti-symmetric part can be written in terms of a single vector. 
A general anti-symmetric matrix has the form
\[
    A = \begin{pmatrix}
    0 & a & b \\
    -a & 0 & c \\
    -b & -c & 0
    \end{pmatrix}
\]
Now consider
\[
    \varepsilon_{ijk}w_k = \begin{pmatrix}
        0 & w_3 & -w_2 \\
        -w_3 & 0 & w_1 \\
        w_2 & -w_1 & 0
    \end{pmatrix}
\]
Therefore, if we choose $w$ appropriately
\[
    w = \begin{pmatrix}
    c \\ -b \\ a
    \end{pmatrix}
    \Rightarrow A_{ij} = \varepsilon_{ijk}w_k
\]
\subsubsection{Matrix Inverse}
Consider an $m \times n$ matrix $A$, $n \times m$ matrix $B$ and $n \times m$ matrix $C$.
\begin{itemize}
    \item If $BA = I_{n \times n}$ then $B$ is the left inverse of $A$
    \item If $AC = I_{m \times m}$ then $C$ is the right inverse of $A$
    \item If $A$ is square ($n \times n$) matrix, then
    \[
        B = B(AC) = (BA)C = C  
    \]
    Therefore the left and right inverses are the same, and we can write $A^{-1}$
    \begin{equation}\label{eq:3-19}
        \Rightarrow AA^{-1} = A^{-1}A = I
    \end{equation}
\end{itemize}
\begin{remark}
    \begin{equation}\label{eq:3-20}
        (AB)^{-1} = B^{-1}A^{-1}
    \end{equation}
    \begin{proof}
        \[
            (B^{-1}A^{-1})(AB) = B^{-1}(A^{-1}A)B = I
        \]
    \end{proof}
\end{remark}

\begin{defi}[Orthogonal]
    A real $n \times n$ matrix is \emph{orthogonal} if
    \[
        A^TA = AA^T = I \Rightarrow A^T = A^{-1}  
    \]
\end{defi}

\begin{defi}[Unitary]
    A complex $n \times n$ matrix is \emph{unitary} if
    \[
        A^\dagger A = AA^\dagger = I \Rightarrow A^\dagger = A^{-1}  
    \]
\end{defi}
\begin{remark}\leavevmode
    \begin{enumerate}
        \item Let $A$ be an orthogonal matrix.
        \begin{align*}
            A_{ik}(A_{kj})^T = \delta_{ij} \\
            \Rightarrow A_{ik}A_{jk} = \delta_{ij}
        \end{align*}
        Which is the scalar product of row $i$ and row $j$.
        Therefore, we see that the rows of an orthogonal matrix from an orthonormal set.
        Similarly for $A^T$, the columns of an orthonormal matrix also form an orthonormal set.
        
%%%%%%%%%%%%%%%%%%%%%%
% Lecture 12: 01/11/2023
%%%%%%%%%%%%%%%%%%%%%%

        \item Similar result for unitary matrix $U$, but now
        \begin{align*}
            (U_{ik})(\overline{U_{kj}})^T = \delta_{ij} \\
            \Rightarrow  \overline{U_{jk}} U_{ik} = \delta_{ij}
        \end{align*}
        i.e; the rows are orthogonal.
    \end{enumerate}
\end{remark}
\begin{eg}\leavevmode
    \begin{enumerate}
        \item Reflection matrix ($3 \times 3$) is orthogonal. $R_{ij} = \delta_{ij} - 2 n_i n_j$
        \begin{align*}
            (RR^T)_{ij} &= R_{ik}R_{kj}^T \\
            &= R_{ik}R_{kj} = (\delta_{ik} - 2 n_i n_k)(\delta_{kj} - 2 n_k n_j) \\
            &= \delta_{ik}\delta_{kj} - 2\delta_{jk}n_in_k - 2\delta_{kj}n_kn_j + 4n_in_jn_kn_k \\
            &= \delta_{ij} - 2 n_in_j -2 n_in_j + 4 n_in_j \\
            &= \delta_{ij}
        \end{align*}
        \item Rotation matrix \eqref{eq:3-2} is orthogonal.
        \begin{ex}
            Express using monster expression for $R_{ij}$ and solve.
        \end{ex}
        Alternatively, denote rotation by $\theta$ about $\uvec n$ as $R(\theta, \uvec n)$. 
        Clearly,
        \[
            R(\theta, \uvec n)^{-1} = R(-\theta, \uvec n)    
        \]
        We check in \eqref{eq:3-2} that
        \[
            R(-\theta, \uvec n)_{ij} = R(\theta, \uvec n)_{ji}
        \]
        Hence
        \[
            R(-\theta, \uvec n) = R^T(\theta, \uvec n) \Rightarrow R^{-1} = R^T
        \]
        Therefore $R$ is an orthogonal matrix.
    \end{enumerate}
\end{eg}
\subsection{Determinants}
\begin{defi}[Determinants in $\R^3$]
    Consider a linear map $\alpha: \R^3 \rightarrow \R^3$ with standard basis $\vec e_1, \vec e_2, \vec e_3$, where
    \[
        \alpha(\vec e_1) = \vec e'_1 \text{ , } \alpha(\vec e_2) = \vec e'_2 \text{ , } \alpha(\vec e_3) = \vec e'_3
    \]
    or in general,
    \[
        \vec e'_i = A\vec e_i \text{ for } i = 1, 2, 3
    \]
    The unit cube is mapped by $\alpha$ to a parallelopiped with volume 
    \begin{align*}
        [\vec e'_1, \vec e'_2, \vec e'_3] &= \varepsilon_{ijk}(\vec e'_1)_i(\vec e'_2)_j(\vec e'_3)_k \\
        &= \varepsilon_{ijk} A_{il}(\vec e'_1)_l A_{jm}(\vec e'_2)_m A_{kn}(\vec e'_3)_n \\
        &= \varepsilon_{ijk} A_{il}\delta_{1l} A_{jm}\delta_{2m} A_{kn}\delta_{3n} \\
        &= \varepsilon_{ijk} A_{i1}A_{j2}A_{k3}
    \end{align*}
    This is the \emph{determinant} of $A$.
    \begin{align*}
        \det (A) &= \varepsilon_{ijk} A_{i1}A_{j2}A_{k3} \\
        &= \begin{vmatrix}
            A_{11} & A_{12} & A_{13} \\
            A_{21} & A_{22} & A_{23} \\
            A_{31} & A_{32} & A_{33} \\
        \end{vmatrix} \numberthis\label{eq:3-21}
    \end{align*}
\end{defi}

Let's extend this to $\R^n$!

\subsubsection{Permutations}
Consider the set $S_n$ of all permutations of numbers.
\[
    1, 2, \cdots, n
\]
which contains $n!$ elements. Consider $\rho \in S_n$ with $i \mapsto \rho(i)$, using the following notation
\begin{align*}
    \rho = \begin{pmatrix}
        1 & 2 & 3 & \cdots & n \\
        \rho(1) & \rho(2) & \rho(3) & \cdots & \rho(n)
    \end{pmatrix} \\[10pt]
    \rho^{-1} = \begin{pmatrix}
        \rho(1) & \rho(2) & \rho(3) & \cdots & \rho(n) \\
        1 & 2 & 3 & \cdots & n
    \end{pmatrix} \\
\end{align*}

\begin{defi}[Fixed Point]
    A \emph{fixed point} of $\rho$ is $k$ such that $\rho(k) = k$.
    \begin{eg}
        \[
            \rho = \begin{pmatrix}
                1 & 2 & 3 & 4 \\
                4 & 1 & 3 & 2
            \end{pmatrix}
        \]
        has a fixed point $3 \mapsto 3$. By convention, we omit the fixed points.
        \[
            \Rightarrow \rho = \begin{pmatrix}
                1 & 2 & 4 \\
                4 & 1 & 2
            \end{pmatrix}
        \]
    \end{eg}
\end{defi}

\begin{defi}[Disjoint]
    Two permutations are are \emph{disjoint} if numbers moved by one permutation are fixed by the other \& vice versa.
\end{defi}
\begin{eg}
    \[
        \begin{pmatrix}
            1 & 2 & 4 & 5 & 6 \\
            5 & 6 & 1 & 4 & 2
        \end{pmatrix} 
        = 
        \begin{pmatrix}
            2 & 6 \\
            6 & 2
        \end{pmatrix}
        \begin{pmatrix}
            1 & 4 & 5 \\
            5 & 1 & 4
        \end{pmatrix}
    \]
\end{eg}
\begin{remark}[Disjoint permutations commute]
    Disjoint permutations \emph{commute}, but in general, permutations do not.
\end{remark}

\begin{defi}[Cycles]
    $\begin{pmatrix}
        2 & 6 \\
        6 & 2
    \end{pmatrix}$ is a $2$-cycle, also called a \emph{transposition}. We write this as
    \[
        (2\ 6)  
    \]
    $\begin{pmatrix}
        1 & 4 & 5 \\
        5 & 1 & 4
    \end{pmatrix}$ is a $3$-cycle, we write
    \[
        (1\ 5\ 4)  
    \]
    For more details, refer to IA Groups Permutation Groups.
\end{defi}

\begin{defi}[Permutation Sign]
    We can write any $q$-cycle as a product of $2$-cycles. For example, 
    \begin{equation}\label{eq:3-22}
        (a_1\ a_2\ \cdots\ a_q) = (a_1\ a_2)(a_2\ a_3)\cdots(a_{q-1}\ a_q)  
    \end{equation}
    For a $q$ cycle, there are $q-1$ transpositions and we apply the maps from right to left.

    The \emph{sign} of a permutation $\epsilon(p)$ is defined to be
    \begin{equation}\label{eq:3-23}
        \epsilon(p) = (-1)^r
    \end{equation}
    where $r$ is the number of transpositions
    \[
        \epsilon(p) = \begin{cases}
            +1 & \text{even perm} \\
            -1 & \text{odd perm}
        \end{cases}
    \]
    Note that
    \begin{align*}
        \epsilon(\rho \sigma) = \epsilon(\rho) \epsilon(\sigma) \\
        \epsilon(\rho^{-1}) = \epsilon(\rho) \numberthis\label{eq:3-24}
    \end{align*}
\end{defi}

\begin{defi}[$n$-dimensional Levi-Cevita]
    Define the $n$-dimensional Levi-Cevita symbol by 
    \[
        \varepsilon_{J_1, J_2, \cdots, J_n} = \begin{cases}
            +1 & \text{ if } J_1, J_2, \cdots, J_n \text{ is an even perm} \\
            -1 & \text{ if } J_1, J_2, \cdots, J_n \text{ is an odd perm} \\
            0 & \text{ otherwise}
        \end{cases}
    \]
    This is the generalisation of $\varepsilon_{ijk}$. Can see from \eqref{eq:2-23} that $\rho \in S_n$
    \[
        \varepsilon_{\rho(1), \rho(2), \cdots, \rho(n)} = \epsilon(p)
    \]
\end{defi}

\begin{defi}[Determinants in $\R^n$]
    Now define the determinant of an $n \times n$ matrix $A$ as either
    \begin{subequations}
        \begin{equation}\label{eq:3-25}
            \det(A) = \sum_{\sigma \in S_n} \epsilon(\sigma) A_{\sigma(1)1}A_{\sigma(2)2} \cdots A_{\sigma(n)n}
        \end{equation}
        Or equivalently,
        \begin{equation}\label{eq:3-25a}
            \det(A) = \varepsilon_{J_1, J_2, \cdots, J_n} A_{J_11}A_{J_22} \cdots A_{J_nn}
        \end{equation}
        Here the suffix notation summing over $J_1, J_2, \cdots, J_n$ and there are $n^n$ terms, most of which have repeated indices and are therefore $0$.
    \end{subequations}
\end{defi}

\subsubsection{Properties of Determinants}
    \begin{enumerate}
        \item \begin{equation}\label{eq:3-26}
            \det A = \det A^T
        \end{equation}
        Consider a single term in \eqref{eq:3-25} and $\rho$ some other perm in $S_n$. 
        \[
            A_{\sigma(1)1}A_{\sigma(2)2} \cdots A_{\sigma(n)n} = A_{\rho(\sigma(1))\rho(1)}A_{\rho(\sigma(2))\rho(2)} \cdots A_{\rho(\sigma(n))\rho(n)}
        \]
        Applying $\rho$ to each term results in the same terms in a different order.
        Now, choose $\rho = \sigma^{-1}$ and note that from \eqref{eq:3-24},
        \[
            \epsilon(\rho) = \epsilon(\rho^{-1}) \Rightarrow \epsilon(\sigma) = \epsilon(\rho)
        \]
        \eqref{eq:3-25} gives 
        \begin{align*}
            \det(A) &= \sum_{\rho \in \S_n} \epsilon(\rho) A_{\rho(1)1}A_{\rho(2)2} \cdots A_{\rho(n)n} \\
            &= \det(A^T)
        \end{align*}

        \item If $B$ is formed by multiplying every element in a single row or column of $A$ by a scalar $\lambda$, then
        \begin{equation}\label{eq:3-27}
            \det (B) = \lambda \det (A)
        \end{equation}
        It follows that 
        \begin{equation}\label{eq:3-28}
            \det(\lambda B) = \lambda^n \det(A)
        \end{equation}
        This follows since, each term in \eqref{eq:3-25} is multiplied by $\lambda$ once to give \eqref{eq:3-27}.
%%%%%%%%%%%%%%%%%%%%%%
% Lecture 13: 03/11/2023
%%%%%%%%%%%%%%%%%%%%%%

        \item If $2$ rows or $2$ columns of $A$ are identicle then
        \begin{equation}\label{eq:3-29}
            \det (A) = 0
        \end{equation}
        \begin{proof}
            Suppose $A$ has columns $1$ and $2$ the same. In \eqref{eq:3-25} write
            \[
                \sigma = \rho(1\ 2)  
            \]
            From \eqref{eq:3-24}
            \[
                \epsilon{\sigma} = -\epsilon{\rho}  
            \]
            Applying the maps,
            \[
                A_{\sigma(1)1} = A_{\rho(2)1} \text{ and } A_{\sigma(2)1} = A_{\rho(1)2}
            \]
            while $A_{\sigma(i)i} = A_{\rho(i)i}$ for $i > 2$.
            \[
                \therefore \det(A) = \sum_{\rho \in S_n} - \epsilon(\rho) A_{\rho(2)1}A_{\rho(1)2} \cdots A_{\rho(n)n}
            \]
            Furthermore, since columns $1$ and $2$ are the same, therefore $A_{\rho(2)1} = A_{\rho(2)2}$ and $A_{\rho(1)2} = A_{\rho(1)1}$
            \begin{align*}
                \Rightarrow \det(A) = - \det(A) \\
                \therefore \det(A) = 0
            \end{align*}
        \end{proof}

        \item If $2$ rows or $2$ columns are linearly dependant then 
        \begin{equation}\label{eq:3-30}
            \det(A) = 0
        \end{equation}
        \begin{proof}
            Suppose $c_r + \lambda c_s = 0$ define
            \[
                B_{ij} = \begin{cases}
                    A_{ij} & j \neq r \\
                    A_{ij} + \lambda A_{is} & j = r
                \end{cases}  
            \]
            Then, 
            \begin{align*}
                \det(B) &= \sum_{\sigma \in S_n} \epsilon(\sigma) B_{\sigma(1)1} \cdots B_{\sigma(r)r} \cdots B_{\sigma(n)n} \\
                &= \sum_{\sigma \in S_n} \epsilon(\sigma) A_{\sigma(1)1} \cdots A_{\sigma(r)r} \cdots A_{\sigma(n)n} \\
                &= \sum_{\sigma \in S_n} \epsilon(\sigma) A_{\sigma(1)1} \cdots [A_{\sigma(r)r} + \lambda A_{\sigma(r)s}] \cdots A_{\sigma(n)n} \\
                &= \det(A) + \lambda \sum_{\sigma \in S_n} \epsilon(\sigma) A_{\sigma(1)1} \cdots A_{\sigma(r)s} \cdots A_{\sigma(n)n} \\
                &= \det(A) + \det(\text{matrix where columns } r = s) \\
                &= \det(A) + 0 \tag{by \eqref{eq:3-29}}
            \end{align*}
            But note that column $r$ of matrix $B$ is all zeros. Therefore
            \[
                \det(B) = 0 \Rightarrow \det(A) = 0
            \]
        \end{proof}

        Note that eeven if $2$ columns are not linearly dependant, we still have the first result where
        \[
            \det(B) = \det(A)  
        \]
        Therefore, we can add a linear multiple of one column (row) onto another column (row) without changing the determinant.

        \item Determinant of a product
        \begin{equation}\label{eq:3-31}
            \det(AB) = \det(A) \det(B)
        \end{equation}
        \begin{proof}
            First, note that for $\rho \in S_n$
            \begin{equation}\label{eq:3-32}
                \sum_{\sigma \in S_n} \epsilon(\sigma) A_{\sigma(1)\rho(1)} \cdots A_{\sigma(n)\rho(n)} = \epsilon(\rho) \det(A) 
            \end{equation}
            That is, swapping columns or rows an even or odd number of times, gives an extra factor of $\pm 1$ respectively. To see this, write
            \[
                \sigma = \mu \rho  
            \]
            then,
            \begin{align*}
                \sum_{\sigma \in S_n} \epsilon(\sigma) A_{\sigma(1)\rho(1)} \cdots A_{\sigma(n)\rho(n)} &= \sum_{\mu \in S_n} \epsilon(\mu \rho) A_{\mu \rho(1)\rho(1)} \cdots A_{\mu \rho(n)\rho(n)} \\
                &= \epsilon(\rho) \sum_{\mu \in S_n} \epsilon(\mu) A_{\mu \rho(1)\rho(1)} \cdots A_{\mu \rho(n)\rho(n)} \\
                &= \epsilon(\rho) \det(A)
            \end{align*}
            Now, 
            \begin{align*}
                \det(AB) &= \sum_{\sigma \in S_n} \epsilon(\sigma) (AB)_{\sigma(1)1} \cdots (AB)_{\sigma(n)n} \\
                &= \sum_{\sigma \in S_n} \epsilon(\sigma) \sum_{k_1, \cdots, k_n} A_{\sigma(1)k_1}B_{k_11} \cdots A_{\sigma(n)k_n}B_{k_nn} \\
                &= \sum_{k_1, \cdots, k_n} B_{k_11} \cdots B_{k_nn} \sum_{\sigma \in S_n} \epsilon(\sigma) A_{\sigma(1)k_1} \cdots A_{\sigma(n)k_n}
            \end{align*}
            In the latter sum, if $2$ terms in $k_1, \cdots, k_n$ are equal, then we have a determinant where two columns are equal and is therefore $0$.
            Therefore, the remaining terms are tuples of distinct $k_1, \cdots, k_n$ which is a re-arrangement of $1, 2, \cdots, n$.
            Hence, we can define a permutation $\rho$ such that $k_i = \rho(i)$. Hence, 
            \begin{align*}
                \det(AB) &= \sum_{\rho \in S_n} B_{\rho(1)1} \cdots B_{\rho(n)n} \sum_{\sigma \in S_n} \epsilon(\sigma) A_{\sigma(1)\rho(1)} \cdots A_{\sigma(n)\rho(n)} \\
                &= \sum_{\rho \in S_n} B_{\rho(1)1} \cdots B_{\rho(n)n} \epsilon(\rho) \det(A) \tag{by \ref{eq:3-32}} \\
                &= \det(B)\det(A)
            \end{align*}
        \end{proof}
        \begin{remark}
            From \eqref{eq:3-31} if $A$ is orthogonal
            \[
                \det(A \cdot A^T) = \det(A) \det(A^T)  
            \]
            By \eqref{eq:3-26} 
            \begin{align*}
                \det(A) = \det(A^T) &\Rightarrow \det(A)^2 = 1 \\
                &\therefore \det(A) = \pm 1 \numberthis\label{eq:3-33}
            \end{align*}

            If $U$ is unitary, 
            \begin{align*}
                \det(U) \cdot \det(U^\dagger ) = 1 \\
                \det(U^\dagger ) = (\det(U^T))^* = \det(U)^* \\
                \Rightarrow \det(U) \cdot \det(U)^* = 1 \\
                \Rightarrow \abs{\det(U)} = 1 \numberthis\label{eq:3-34}
            \end{align*}
            Hence, $\det(U)$ lies on the unit circle. 
            
            As we've seen before, in $\R^3$ orthogonal matrices represent a rotation ($\det = \pm 1$) or a reflection ($\det = -1$).
        \end{remark}
    \end{enumerate}

\subsubsection{Minors \& Cofactors}
For an $n \times n$ matrix $A$. 
Define $A^{ij}$ to be the $(n - 1) \times (n - 1)$ matrix in which row $i$ and column $j$ are removed.

\begin{defi}[Minor]
    The \emph{minor} of the $ij^{th}$ element of $A$ denoted $M_{ij}$ is
    \begin{equation}\label{eq:3-35}
        M_{ij} = \det(A^{ij})
    \end{equation}
\end{defi}

\begin{defi}[Cofactor]
    The \emph{cofactor} of the $ij^{th}$ element of $A$, denoted
    \begin{equation}\label{eq:3-36}
       \bigtriangleup_{ij} = (-1)^{i - j} M_{ij} 
    \end{equation}
\end{defi}

Now, we will determine an alternative expression for $\det A$ in terms of the cofactors.
Firstly, introduce more notation. Use $\overline{\cdot}$ to denote a symbol which is missing in a natural sequence.
\begin{eg}
    $(1, 2, 3, 5)$ is $(1, 2, 3, \overline{4}, 5)$
\end{eg}

For a fixed column $I$ equation \eqref{eq:3-25a} without suffix notation becomes,
\begin{equation}\label{eq:3-37}
    \det A = \sum_{J_I = 1}^{n} A_{J_II} \sum_{J_1, \cdots, \overline{J_I}, \cdots, J_n} \varepsilon_{J_1, \cdots, J_n} A_{J_11}A_{J_22} \cdots \overline{A_{J_II}} \cdots A_{J_nn}
\end{equation}
where we are expanding the determinant formula along the $I^{th}$ column. For instance,
\begin{itemize}
    \item $J_I = 1$ sums over all permutations that contain the first element of the $I^{th}$ row.
    \item $J_I = 2$ sums over all permutations that contain the second element of the $I^{th}$ row.
    \item etc.
\end{itemize}

Let $\sigma \in S_n$ be the permutation which moves $J_I$ to the $I^{th}$ position and leaves the remaining numbers in their natural order
\[
    \sigma = \begin{cases}
        \begin{pmatrix}
            1 & \cdots & I & I + 1 & \cdots & J_I - 1 & J_I & \cdots & n \\
            1 & \cdots & J_I & I & \cdots & J_I - 2 & J_I -1 & \cdots & n
        \end{pmatrix} & \text{ if } I < J_I \\[15pt]
        \begin{pmatrix}
            1 & \cdots & J_I & J_I + 1 & \cdots & I - 1 & I & \cdots & n \\
            1 & \cdots & J_I + 1 & J_I + 2 & \cdots & I & J_I & \cdots & n
        \end{pmatrix} & \text{ if } I > J_I \\[15pt]
        \text{ identity permutation } & \text{ if } I = J_I
    \end{cases}
\]
\begin{eg}\leavevmode
    \begin{enumerate}
        \item Take for $I = 2$, $J_I = 4$ i.e, $4$ to the $2^{nd}$ position
        \[
            \begin{pmatrix}
                1 & 2 & 3 & 4 & 5 & 6 \\
                1 & 4 & 2 & 3 & 5 & 6
            \end{pmatrix}
        \]
        \item $I = 5$, $J_I = 2$ i.e. $2$ to the $5^{th}$ position
        \[
            \begin{pmatrix}
                1 & 2 & 3 & 4 & 5 & 6 \\
                1 & 3 & 4 & 5 & 2 & 6
            \end{pmatrix}
        \]
    \end{enumerate}
\end{eg}

%%%%%%%%%%%%%%%%%%%%%%
% Lecture 14: 06/11/2023
%%%%%%%%%%%%%%%%%%%%%%

In all of the $3$ cases, $J_I < I$, $J_I > I$, $J_I = I$, $\sigma$ is a product of $\abs{J_I - I}$ such transpositions.
So from \eqref{eq:3-23}
\[
    \epsilon(\sigma) = (-1)^{I - J_I}
\]
Now consider another permutation $\rho \in S_n$ on the $n-1$ elements,
\[
    \rho = \begin{pmatrix}
        1 & \cdots & \cdots & \overline{J_I} & \cdots & n \\
        J_1 & \cdots & \overline{J_I} & \cdots & \cdots & J_n
    \end{pmatrix}
\]
We have written out $\sigma$ and $\rho$ such that the combination $\rho \sigma$ reorders.
\[
    (1, 2, \cdots, n) \mapsto (J_1, J_2, \cdots, J_n)
\]
To see that this is the case, we can expand $\rho$ into
\[
    \rho = \begin{cases}
        \begin{pmatrix}
            1 & \cdots & I-1 & I & \cdots & J_I -1 & J_I +1 & \cdots & n \\
            J_1 & \cdots & J_{I-1} & J_{I+1} & \cdots & J_{(J_I)} & J_{(J_I +1)} & \cdots & J_n
        \end{pmatrix} & \text{ if } I < J_I \\[15pt]
        \begin{pmatrix}
            1 & \cdots & J_I -1 & J_I +1 & \cdots & I & I + 1 & \cdots & n \\
            J_1 & \cdots & J_{(J_I -1)} & J_{(J_I)} & \cdots & J_{I-1} & J_{I+1} & \cdots & J_n
        \end{pmatrix} & \text{ if } I > J_I \\[15pt]
        \text{ identity permutation} & \text{ if } I = J_I
    \end{cases}
\]
Furthermore, we see that
\[
    \epsilon(\rho) = \varepsilon_{J_1, \cdots, \overline{J_I}, \cdots, J_n}
\]
and we can decompose
\begin{align*}
    \epsilon(\rho \sigma) = \varepsilon_{J_1, \cdots, J_n} &= \epsilon(\rho)\epsilon(\sigma) \\
    &= \varepsilon_{J_1, \cdots, \overline{J_I}, \cdots, J_n} \times (-1)^{I - J^I}
\end{align*}
Hence, \eqref{eq:3-37} becomes
\begin{align*}
    \det (A) &= \sum_{J_I = 1}^{n} A_{J_II} \sum_{J_1, \cdots, \overline{J_I}, \cdots, J_n} \varepsilon_{J_1, \cdots, \overline{J_I}, \cdots, J_n} (-1)^{I - J^I} A_{J_11} \cdots \overline{A_{J_II}} \cdots A_{J_nn} \\
    &= \sum_{J_I = 1}^{n} A_{J_II} (-1)^{I - J^I} \sum_{J_1, \cdots, \overline{J_I}, \cdots, J_n} \varepsilon_{J_1, \cdots, \overline{J_I}, \cdots, J_n} A_{J_11} \cdots \overline{A_{J_II}} \cdots A_{J_nn} \\
    &= \sum_{J_I = 1}^{n} A_{J_II} (-1)^{I - J^I} M_{J_II} \\
    &= \sum_{J_I = 1}^{n} A_{J_II} \bigtriangleup_{J_II} \numberthis\label{eq:3-38}
\end{align*}
Using \eqref{eq:3-26} we can do the same to $A^T$ to see that
\begin{equation}\label{eq:3-39}
    \det (A) = \sum_{J_I = 1}^{n} A_{IJ_I} \bigtriangleup_{IJ_I}
\end{equation}
These are the Laplace expansion formulae, and they allow us to expand along any column or row.

\begin{eg}\leavevmode
    \begin{enumerate}
        \item $2 \times 2$ determinant. We use \eqref{eq:3-25}
        \[
            \det (A) = \sum_{\sigma \in S_2} \epsilon(\sigma)A_{\sigma(1)1}A_{\sigma(2)2}  
        \]
        where the elements of $S_2$ are
        \begin{align*}
            S_2 = \crbkt{e, (1\ 2)} \\
            \epsilon(e) = 1 \\
            \epsilon(1\ 2) = -1
        \end{align*}
        Therefore,
        \begin{align*}
            \det (A) &=  \epsilon(e)A_{11}A_{22} + \epsilon(1\ 2)A_{21}A_{12} \\
            &= A_{11}A_{22} - A_{21}A_{12} \numberthis\label{eq:3-40}
        \end{align*}

        \item $3 \times 3$ determinant.
        \[
            A = \begin{vmatrix}
                2 & 4 & 1 \\
                3 & 2 & 1 \\
                2 & 0 & 1
            \end{vmatrix}  
        \]
        Use \eqref{eq:3-29} and choose $I = 1$ (First row)
        \begin{align*}
            \det(A) = 2 \begin{vmatrix}
                2 & 1 \\
                0 & 1
            \end{vmatrix} 
            -4 \begin{vmatrix}
                3 & 1 \\
                2 & 1
            \end{vmatrix}
            + 2 \begin{vmatrix}
                3 & 2 \\
                2 & 0
            \end{vmatrix} \\
            = 2 \times 2 - 4 \times 1 + 2 \times -4 = -8
        \end{align*}
        Or use \eqref{eq:3-38} and choose $I = 2$ (Second Column)
        \[
            \det(A) = -4 \begin{vmatrix}
                3 & 1 \\
                2 & 1
            \end{vmatrix} 
            2 \begin{vmatrix}
                2 & 2 \\
                2 & 1
            \end{vmatrix}
            + (-1) \cdot 0 = -8 
        \]
        In practical terms, we use combinations of properties of matrices together with appropriate choices of $I$ to evaluate $\det(A)$

        \item
        \begin{align*}
            \det(A) &= \begin{vmatrix}
                1 & a & a^2 \\
                1 & b & b^2 \\
                1 & c & c^2
            \end{vmatrix} \\[15pt]
            &= \begin{vmatrix}
                0 & a - b & a^2 - b^2\\
                1 & b & b^2 \\
                1 & c & c^2
            \end{vmatrix} \tag{row $1$ $\rightarrow$ row $1$ - row $2$} \\[15pt]
            &= (a - b)\begin{vmatrix}
                0 & 1 & a + b \\
                1 & b & b^2 \\
                1 & c & c^2
            \end{vmatrix} \tag{factor}\\[15pt]
            &= (a - b)\begin{vmatrix}
                0 & 1 & a + b \\
                0 & b - c & b^2 - c^2 \\
                1 & c & c^2
            \end{vmatrix} \tag{row $2$ $\rightarrow$ row $2$ - row $3$} \\[15pt]
            &= (a - b)(b - c)\begin{vmatrix}
                0 & 1 & a + b \\
                0 & 1 & b + c \\
                1 & c & c^2
            \end{vmatrix} \tag{factor}\\[15pt]
            &= (a - b)(b - c)\begin{vmatrix}
                0 & 0 & a - c \\
                0 & 1 & b + c \\
                1 & c & c^2
            \end{vmatrix} \tag{row $1$ $\rightarrow$ row $1$ - row $2$} \\[15pt]
            &= (a - b)(b - c)(a - c)\begin{vmatrix}
                0 & 0 & 1 \\
                0 & 1 & b + c \\
                1 & c & c^2
            \end{vmatrix} \tag{factor}
        \end{align*}
        Then we choose $I = 1$ and expand along the first row.
        \[
            \det(A) = (a - b)(b - c)(a - c) \begin{vmatrix}
                0 & 1 \\
                1 & c
            \end{vmatrix} = (a - b)(b - c)(c - a)
        \]
    \end{enumerate}
\end{eg}

\section{Matrices and Linear Equations}
\subsection{Simple Example $2 \times 2$ matrix}
\begin{align}
    A_{11}x_1 + A_{12}x_2 = d_1 \label{eq:4-1} \\
    A_{21}x_1 + A_{22}x_2 = d_2 \label{eq:4-2}
\end{align}
Solve $A \vec x = \vec d$ where $A$ is a $2 \times 2$ matrix. \eqref{eq:4-1} $\times A_{22} -$ \eqref{eq:4-2} $\times A_{11}$ gives, 
\begin{align*}
    (A_{11}A_{22} - A_{12}A_{21}) x_1 &= A_{22}d_1 - A_{12}d_2 \\
    (A_{11}A_{22} - A_{12}A_{21}) x_2 &= A_{11}d_2 - A_{21}d_1
\end{align*}
Where $A_{11}A_{22} - A_{12}A_{21} = \det A$ \eqref{eq:3-40}.
If $\det(A) \neq 0$
\[
    \begin{pmatrix}
        x_1 \\ x_2
    \end{pmatrix}
    = \frac{1}{\det(A)}
    \begin{pmatrix}
        A_{22} & -A_{12} \\
        -A_{21} & A_{11}
    \end{pmatrix}
    \begin{pmatrix}
        d_1 \\ d_2
    \end{pmatrix}
\]
Also, $A \vec x = \vec d$, so if $A^{-1}$ exists then $\vec x = A^{-1} \vec d$.
Hence we have constructed $A^{-1}$ in the $2 \times 2$ case.
\begin{equation}
    A^{-1} = \frac{1}{\det(A)}
    \begin{pmatrix}
        A_{22} & -A_{12} \\
        -A_{21} & A_{11}
    \end{pmatrix}
\end{equation}
Therefore we see that $A^{-1}$ exists if and only if $\det(A) \neq 0$
\subsection{Inverse of $n \times n$ matrix}
Consider a matrix $B$ which is identical to $A$ except that the $I^{th}$ row of $A$ is replaced with $i^{th}$ row of $A$ where $i \neq I$.
That is, $B$ is a copy of $A$ that has two identical rows - $i^{th}$ row and $I^{th}$ row.
\[
    \therefore \bigtriangleup_{Ij} \text{ of } B = \bigtriangleup_{Ij} \text{ of } A
\]
Since two rows of $B$ are the same, 
\begin{align*}
    \det(B) &= 0 \\
    &= \sum_{k=1}^{n} B_{Ik} \bigtriangleup_{Ik} \tag{from \eqref{eq:3-39}} \\
    &= \sum_{k=1}^{n} A_{ik} \bigtriangleup_{Ik} \numberthis\label{eq:4-4}
\end{align*}
For $i \neq I$. Similarly, consider matrix $C$ in which we take $A$ and replace the $J^{th}$ column with the $j^{th}$ column where $j \neq J$.
\begin{equation}\label{eq:4-5}
    \det(C) = 0 = \sum_{k=1}^{n} A_{kj} \bigtriangleup_{kJ}
\end{equation}
For $j \neq J$.
\begin{subequations}
    Together, \eqref{eq:3-39} and \eqref{eq:4-4} implies (with summation notation)
    \begin{align*}
        \Rightarrow A_{ik} \bigtriangleup_{jk} &= \begin{cases}
            0 & \text{ if } i \neq j \\
            \det(A) & \text{ otherwise}
        \end{cases} \\
        &= (\det(A)) \delta_{ij} \numberthis\label{eq:4-6}
    \end{align*}
    \eqref{eq:3-38} and \eqref{eq:4-5} implies
    \begin{align*}
        \Rightarrow A_{ki} \bigtriangleup_{kj} &= \begin{cases}
            0 & \text{ if } i \neq j \\
            \det(A) & \text{ otherwise}
        \end{cases} \\
        &= (\det(A)) \delta_{ij} \numberthis\label{eq:4-6a}
    \end{align*}
\end{subequations}

\begin{thm}
    If $\det A \neq 0$ then $A^{-1}$ exists and is given by 
    \begin{equation}\label{eq:4-7}
        (A^{-1})_{ij} = \frac{\bigtriangleup_{ji}}{\det(A)}
    \end{equation}
\end{thm}
\begin{proof}
    \begin{align*}
        (A^{-1})_{ik} A_{kj} = \frac{\bigtriangleup_{ki}}{\det(A)} \cdot A_{kj} &= \frac{\delta_{ij} \det(A)}{\det(A)} \tag{by \eqref{eq:4-6a}} \\
        &= \delta_{ij}
    \end{align*}
\end{proof}

\begin{eg}
    Shear matrix 
    \[
        S_\lambda = \begin{pmatrix}
            1 & \lambda & 0 \\
            0 & 1 & 0 \\
            0 & 0 & 1
        \end{pmatrix}  
    \]
    Using \eqref{eq:4-7} we can show that $S_\lambda^{-1} = S_{-\lambda}$
\end{eg}

\begin{remark}[Computational Complexity]
    How many arithmetic operations are involved in inverting and $n \times n$ matrix using \eqref{eq:4-7}?

    To simplify lets' just count the multiplications.
    \begin{enumerate}
        \item $\det A$, say takes $f_n$ operations, which involves calculating $n$ determinants of $(n-1) \times (n-1)$ matrices and then $n$ more multiplications to combine them.
        \begin{align*}
            &f_n = n f_{n-1} + n \\
            \Rightarrow& f_n = n! + n + n(n-1) + n(n-1)(n-2) + \cdots \\
            \Rightarrow& \frac{f_n}{n!} = 1 + \frac{1}{(n-1)!} + \frac{1}{(n-2)!} + \cdots + \frac{1}{(2)! + 1}
        \end{align*}
        Therefore as $n \rightarrow \infty \Rightarrow f_n \sim (e+ 1)n!$

        \item To populate matrix in \eqref{eq:4-7}, we need $n^2$ elements. We already have $n$ cofactors from (i) so $n(n-1)$ more cofactors needed.
    \end{enumerate}
    Therefore (i) + (ii) $= O(n \times n!)$
\end{remark}

%%%%%%%%%%%%%%%%%%%%%%
% Lecture 15: 08/11/2023
%%%%%%%%%%%%%%%%%%%%%%
\subsection{Homogenous and Inhomogenenous Equations}
Consider
\begin{equation}\label{eq:4-8}
    A \vec x = \vec b
\end{equation}
for $A$ an $n \times n$ matrix, $\vec x$ and $\vec b$ are $n \times 1$ columns.
If $\vec b = \vec 0$ the system is homogeneou, if $\vec b \neq \vec 0$ then we say it is Inhomogenenous.

Suppose $\det(A) \neq 0$ then by \eqref{eq:4-7} $A^{-1}$ exists and \eqref{eq:4-8} has a unique solution
\[
    \vec x =  A^{-1} \vec b
\]
where $\vec x = \vec 0$ if $\vec b = \vec 0$. 
Note that $\det A \neq 0$ implies that the columns of $A$ are linearly independant \eqref{eq:3-29}. 
The columns of $A$ are the images of the standard basis $\vec e'_i = A\vec e_i$.

Hence, $\det A \neq 0 \Longrightarrow$ $e'_i$ are linearly independant and there are $n$ of them
$\Longrightarrow$ they form a basis of $\R^n$ (or $\C^n$), $\therefore r(A) = n$.
By \eqref{eq:3-9}, $n(A) = 0$, $\therefore$ kernel of linear map is $\crbkt{\vec 0}$.

\subsubsection{Gaussian Elimination}
Consider
\begin{align*}
    A_{11}x_1 + A_{12}x_2 + \cdots + A_{1n}x_n &= d_1 \\
    A_{21}x_1 + A_{22}x_2 + \cdots + A_{2n}x_n &= d_1 \\
    \vdots & \\
    A_{m1}x_1 + A_{m2}x_2 + \cdots + A_{mn}x_n &= d_m \numberthis\label{eq:4-9}
\end{align*}
i.e; $m$ equations in $n$ unknowns. Assume $A_{11} \neq 0$ (otherwise reorder equations), then we may reduce the above to
\[
    \begin{array}{ccccccccccc}
        A_{11}x_1 & + & A_{12}x_2      & +      & \cdots          & \cdots & + & A_{1n}x_n       & = & d_1 \\
        0         & + & A^{(2)}_{22}x_2& +      & \cdots          & \cdots & + & A^{(2)}_{2n}x_n & = & d^{(2)}_2 \\
                  &   &                & \ddots &                 &        &   &                 &   & \vdots \\
                  &   &                &        & A^{(r)}_{rr}x_r & \cdots & + & A^{(r)}_{rn}x_n & = & d^{(r)}_{r} \\
                  &   &                &        &                 &        &   & 0               & = & d^{(r)}_{r + 1} \\
                  &   &                &        &                 &        &   & \vdots          &   & \vdots \\
                  &   &                &        &                 &        &   & 0               & = & d^{(r)}_{m}    
    \end{array}
\]
Here $A^{(i)}_{ii} \neq 0$ and the superfix refers to the version number of the coefficients. 
For instance, $A^{(2)}_{22}$ is the second version of $A_{22}$ after $x_1$ has been eliminated.
There are the following possibilities.
\begin{enumerate}
    \item $r \leq m$ and at least one of the $d^{(r)}_{r + 1}, \cdots, d^{(r)}_{m} \neq 0$ then the equations \eqref{eq:4-9} are inconsistent and there is no solution.
    We say that the system is \emph{overdetermined}.
    \begin{eg}
        \begin{align*}
            \begin{array}{ccccccc}
                3x_1 & + & 2x_2 & + & x_3 & = & 3 \\
                6x_1 & + & 3x_2 & + & 3x_3 & = & 0 \\
                6x_1 & + & 2x_2 & + & 4x_3 & = & 6 
            \end{array}\\[15pt]
            \Rightarrow \begin{array}{ccccccc}
                3x_1 & + & 2x_2 & + & x_3 & = & 3 \\
                0    & - &  x_2 & + & x_3 & = & -6 \\
                0    & - & 2x_2 & + & 2x_3 & = & 0 
            \end{array}\\[15pt]
            \Rightarrow \begin{array}{ccccccc}
                3x_1 & + & 2x_2 & + & x_3 & = & 3 \\
                     & - &  x_2 & + & x_3 & = & -6 \\
                     &   & 0    & + & 0   & = & 12  
            \end{array}
        \end{align*}
        Hence, no solutions.
    \end{eg}

    \item $r = n \leq m$ and $d^{(r)}_{i} = 0$ for $i = r + 1, \cdots, m$. 
    We have a unique solution for $x_n$ and can determine all $x_i$ by back substitution.
    We say that the system is \emph{determined}
    \begin{eg}
        \begin{align*}
            \begin{array}{ccccc}
                2x_1 & + & 5x_2 & = & 2 \\
                4x_1 & + & 3x_2 & = & 11
            \end{array}\\[15pt]
            \Rightarrow \begin{array}{ccccc}
                2x_1 & + & x_2 & = & 2 \\
                0    & - & 7x_2 & = & 7
            \end{array}\\[15pt]
            \Rightarrow x_2 = -1, \quad x_1 = \frac{7}{2}
        \end{align*}
    \end{eg}

    \item $r \leq n$ and $d^{(r)}_{i} = 0$ for $i = r + 1, \cdots, m$. 
    This implies that $x_{r + 1}, \cdots x_{n}$ can be chosen freely and therefore there are infinitely many solutions.
    We say that the system is \emph{underdetermined}
    \[
        \begin{array}{ccccc}
            x_1 & + & x_2 & = & 1 \\
            2x_1& + & 2x_2 & = & 1
        \end{array}
        \Rightarrow x_1 + x_2 = 1
    \]
    Hence, $x_2$ can be anything and $x_1 = 1 - x_2$
\end{enumerate}

\begin{remark}
    In the case that $n=m$, there are $O(n^3)$ operations required which is much faster than \eqref{eq:4-7} for solving \eqref{eq:4-8}.
\end{remark}

\subsubsection{Relation to Determinant}.
When $m = n$, $A$ is square,
\begin{equation}\label{eq:4-10}
    \det A = (-1)^k \abs{\begin{array}{cccccc}
        A_{11} & A_{12}                  & \cdots & \cdots       & \cdots & A_{1n} \\
        0      & A^{(2)}_{22}            & \cdots & \cdots       & \cdots & A^{2}_{2n} \\
               &                         & \ddots &              &        & \vdots      \\
               &                         &        & A^{(r)}_{rr} & \cdots & A^{(r)}_{rn} \\
               &                         &        &      0       & \cdots & 0      \\
               & \smash{\scalebox{4}{0}} &        &              & \ddots & \vdots \\
               &                         &        &              &        & 0 \\              
    \end{array}}
\end{equation}
where $k$ comes from the number of times we have reordered equations and $A^{(\alpha)}_{\alpha \alpha} \neq 0$

This is an upper triangular matrix. 
The determinant of an upper triangular matrix is the product of its diagonal elements (directly from \eqref{eq:3-25}),
the only $\sigma \in S_n$ such that $\sigma(i) \geq i \quad \forall i$ is the identity permutation.

Hence, if $r < n \Rightarrow \det A = 0$, if $r = n \Rightarrow \det A = (-1)^kA_{11}A^{(2)}_{22} \cdots A^{(n)}_{nn}$

\subsection{Matrix Rank}
Consider a linear map $\alpha: \R^n \rightarrow \R^n$ (works for $\C^n$ as well). 
Recall that the rank $r(\alpha)$ is the dimension of the image.
If $A$ is the associated matrix for $\alpha$, call rank $r(A)$.

Recall that if $\vec e_1, \cdots, \vec e_n$ are the standard basis of $\R^n$, 
then $A\vec e_1, \cdots, A\vec e_n$ span the image. (Not necessarily linearly independant, depends on the rank)!

Furthermore $A\vec e_1, \cdots, A\vec e_n$ are the columns of $A$. Hence
\begin{equation}\label{eq:4-11}
    r(A) = \text{\# of linearly independant columns}
\end{equation}

\begin{defi}[Column and row rank]\leavevmode
    \begin{itemize}
        \item Column rank of $A$ $=$ max number of linearly independant columns
        \item Row rank of $A$ $=$ max number of linearly independnat row
    \end{itemize}
\end{defi}

\begin{thm}[Column and row rank]
    Row rank $=$ Column rank for $m \times n$ matrix
\end{thm}
\begin{proof}
    For an $m \times n$ matrix, there are
    \begin{itemize}
        \item $m$ rows, each of which has $n$ elements
        \item $n$ columns, each of which has $m$ elements
    \end{itemize}

    Let $r$ be the row rank of $A$ and write the biggest set of linearly independant rows as
    \[
        \vec V_1^T, \vec V_2^T, \cdots, \vec V_r^T,   
    \]
    where $V_i$ are $n \times 1$ column vectors, or in component form,
    \[
        \vec V_k^T = (V_{k1}, V_{k2}, \cdots, V_{kn}) \quad \text{ for } k = 1, 2, \cdots, r
    \]
    Now, denote the $i^{th}$ row of $A$ as
    \[
        r_i^T = (A_{i1}, A_{i2}, \cdots, A_{in}) \quad \text{ for } i = 1, 2, \cdots, m
    \]
    Since every row of $A$ can be written as a linear combination of the $\vec V$'s,
    \[
        r_i^T = \sum_{k=1}^{r} C_{ik} V^T_k \quad \text{ for } i = 1, 2, \cdots, m
    \]
    Now the elements of $A$ are
    \[
        A_{ij} = (r_i^T)_j = \sum_{k=1}^{r} C_{ik} (V^T_k)_j  \quad \text{ for } j = 1, 2, \cdots, n
    \]
    Or, 
    \[
        \begin{pmatrix}
            A_{1j} \\ A_{2j} \\ \vdots \\ A_{mj}
        \end{pmatrix}
        =
        \sum_{k=1}^{r} (V^T_k)_j \begin{pmatrix}
            C_{1k} \\ C_{2k} \\ \vdots \\ C_{mk}
        \end{pmatrix}
    \]
    Therefore, any column of $A$ can be expressed as a linear combination of the $r$ ``coefficient'' column vectors.
    Hence, the column rank of $A$ $\leq$ row rank of $A$. 
    We apply the same argument to $A^T$ to show the converse that row rank of $A$ $\leq$ column rank of $A$.
    Therefore, 
    \[
        \text{row rank of } A = \text{ column rank of } A = r(A) 
    \]
\end{proof}

%%%%%%%%%%%%%%%%%%%%%%
% Lecture 16: 10/11/2023
%%%%%%%%%%%%%%%%%%%%%%

\subsection{Homogeneous problems}
We will restrict our attention to $A$ an $n \times n$ matrix, and consider
\[
    A \vec x = \vec 0  
\]
If $\det A \neq 0$ then we have a unique solution
\[
    \vec x = A^{-1} \vec 0 = \vec 0   
\]
Hence. 
\[
    A \vec x = \vec 0 \text{ and } x \neq 0 \Rightarrow \det A = 0
\]

\subsubsection{Geometrical Interpretation for the $3 \times 3$ case}
\[
    A = \begin{pmatrix}
        \vec r_1^T \\[10pt]
        \vec r_2^T \\[10pt]
        \vec r_3^T
    \end{pmatrix} 
    \Rightarrow \vec r_i \cdot x = 0 \quad \text{ for } i = 1, 2, 3
\]
The solution of $A \vec x = \vec 0$ corresponds to the intersection of $3$ planes through the origin.
There are 3 possibilities:
\begin{enumerate}
    \item intersect only at $\vec 0$
    \item intersect in a common line through $\vec 0$
    \item planes coincide
\end{enumerate}

\begin{description}
    \item[(i)] $\det(A) \neq \vec 0 \Rightarrow [\vec r_1, \vec r_2, \vec r_3] \neq \vec 0$
    \begin{align*}
        \Rightarrow \spnset{\vec r_1, \vec r_2, \vec r_3} = \R^3 \\
        \therefore r(A) = 3 \Rightarrow n(A) = 0
    \end{align*}
    Therefore the kernel is $\crbkt{\vec 0}$ and the image in $\R^3$.

    \item[(i), (ii)] $\det(A) = 0 \Rightarrow \dim(\spnset{\vec r_1, \vec r_2, \vec r_3}) = 1 \text{ or } 2$.
    \begin{itemize}
        \item If rank $=2$, take $\vec r_1, \vec r_2$ linearly independant. 
        Then $\vec x$ lies on the intersection of $\vec x \cdot \vec r_1 = 0$ and $\vec x \cdot \vec r_2 = 0$, which is the line
        \[
            \vec x \in \R^3 \mid \vec x = \lambda \vec t, \quad \lambda \in \R, \quad \vec t = \vec r_1 \times \vec r_2
        \]
        and all points on this line satisfy $\vec x \cdot \vec r_3 = 0$, since $\vec r_3$ is a linear combination of $\vec r_1$ and $\vec r_2$.
        Therefore, the kernel is a line,
        \[
            n(A) = 1 \text{ and } r(A) = 2  
        \]

        \item If rank $=1$, then $\vec r_1, \vec r_2, \vec r_3$ are parallel, So
        \[
            \vec x \cdot \vec r_1 = 0 \Rightarrow \vec x \cdot \vec r_2 = 0 \text{ and } \vec x \cdot \vec r_3 = 0   
        \]
        therefore the kernel is a plane,
        \[
            n(A) = 2 \text{ and } r(A) = 1
        \]

        \item If rank $=0$ we have $A_{ij} = 0$, therefore the $0$ map.
    \end{itemize}
\end{description}

\subsubsection{Linear mapping view of $A \vec x = \vec 0$}
Consider a linear map 
\begin{align*}
    \alpha: \R^n &\rightarrow \R^n \\
    \vec x &\mapsto \vec x' = A\vec x 
\end{align*}
The kernel of $\alpha$ is the solution space of $A \vec x = \vec 0$.
\begin{enumerate}
    \item If $n(A) = 0 \Rightarrow A(e_1), A(e_2), \cdots, A(e_n)$ form a linearly independant set of $n$ elements and $r(A) = n$
    \item If $n(A) > 0$, Let
    \[
        \crbkt{\vec u_i} \text{ for } i = 1, 2, \cdots, n(A)  
    \]
    be a basis for the kernel $\kappa (A)$. Hence, any solution of $A \vec x = \vec 0$ is 
    \[
        \vec x = \sum_{i=1}^{n(A)}\lambda_i \vec u_i  
    \]
    We can extend this to a basis of $\R^n$ by introducing new vectors
    \[
        \crbkt{\vec u_i} \text{ for } i = n(A)+1, n(A)+2, \cdots, n
    \]
    The vectors $A\vec u_i$ for $i = n(A)+1, \cdots, n$ are then a basis of the image of $\alpha$.
\end{enumerate}

\subsection{General solution of $A\vec x = \vec d$}
\begin{enumerate}
    \item $\det A \neq 0 \Rightarrow (A) = 0, r(A) = n$ and $A(\R^n) = \R^n$.
    So for any $\vec d \in \R^n$ $\exists \vec x$ such that
    \[
        \vec x \mapsto \vec d \Leftrightarrow \vec x = A^{-1} \vec d 
    \]
    
    \item $\det A = 0 \Rightarrow n(A) > 0, r(A) < n$, The image $I(A)$ is a proper subspace of $\R^n$.
    \begin{enumerate}
        \item If $\vec d \notin I(A)$, there are no solutions
        \item If $\vec d \in I(A)$ at least one solution exists.
    \end{enumerate}

    If $\vec d \in I(A)$ the general solution can be written as
    \[
        \vec x = \vec x_0 + \vec y  
    \]
    where $\vec x_0$ is a particular solution, i.e $A \vec x_0 = \vec d$ and $\vec y$ is any vector from the kernel,
    \[
        \vec y \in \ker (A) \Leftrightarrow A \vec y = \vec 0  
    \]
    \begin{enumerate}
        \item If $n(A) = 0 \Rightarrow \vec y = \vec 0$ and we have a unique solution
        \item If $n(A) > 0$ then let $u_1 \cdots, u_{n(A)}$ be a basis of the kernel,
        \[
            \vec y = \sum_{j=1}^{n(A)} \lambda_j u_j  
        \]
    \end{enumerate}
    and the most general solution is given by
    \[
        \vec x = \vec x_0 + \sum_{j=1}^{n(A)} \lambda_j u_j
    \]
\end{enumerate}

\begin{eg}\leavevmode
    \begin{enumerate}[label=(\arabic*)]
        \item $\begin{pmatrix}
            1 & 1 \\
            a & 1
        \end{pmatrix}
        \begin{pmatrix}
            x_1 \\ x_2
        \end{pmatrix}=
        \begin{pmatrix}
            1 \\ b
        \end{pmatrix}$

        We solve for the determinant
        \[
            \det A = 1 - a
        \]
        Therefore if $a \neq 1$, $A^{-1}$ exists and is
        \[
            A^{-1} = \frac{1}{1-a} \begin{pmatrix}
                1 & -1 \\
                -a & 1 
            \end{pmatrix}
        \]
        Then we have a unique solution
        \[
            \vec x = \frac{1}{1-a} \begin{pmatrix}
                1- b \\
                -a + b 
            \end{pmatrix}
        \]
        If $a = 1$, then,
        \[
            A \vec x = \begin{pmatrix}
                x_1 + x_2 \\
                x_1 + x_2
            \end{pmatrix} = (x_1 + x_2) \begin{pmatrix}
                1 \\ 1
            \end{pmatrix}
        \]
        Hence,
        \begin{align*}
            \im(A) = \spnset{\begin{pmatrix}
                1 \\ 1
            \end{pmatrix}}, \quad r(A) = 1 \\
            \ker(A) = \spnset{\begin{pmatrix}
                1 \\ -1
            \end{pmatrix}}, \quad n(A) = 1
        \end{align*}

        Hence, if $b \neq 1$ then $\begin{pmatrix}
            1 \\ b
        \end{pmatrix} \notin \im(A)$ therefore no solutions.
        Otherwise, $b = 1$ and $\begin{pmatrix}
            1 \\ 1
        \end{pmatrix} \in \im(A)$ and we can find a patricular solution $\begin{pmatrix}
            x_1 \\ x_2
        \end{pmatrix} = \begin{pmatrix}
            1 \\ 0
        \end{pmatrix} = \vec x_0$
        Therefore we can find a general solution
        \[
            \vec x = \begin{pmatrix}
                1 \\ 0
            \end{pmatrix} + \mu \begin{pmatrix}
                1 \\ -1
            \end{pmatrix}  
        \]

        \item Find the general solution of
        \[
            \begin{pmatrix}
                a & a & b \\
                b & a & a \\
                a & b & a
            \end{pmatrix}
            \begin{pmatrix}
                x \\ y \\ z
            \end{pmatrix}
            =
            \begin{pmatrix}
                1 \\ c \\ 1
            \end{pmatrix}
        \]
        We solve for the determinant
        \[
            \det A = (a - b)^2 (2a + b)
        \]
        So if $a \neq b$ and $2a \neq -b$ then $A^{-1}$ exists and we can find a unique solution
        \[
            \vec x = A^{-1} \begin{pmatrix}
                1 \\ c \\ 1
            \end{pmatrix}
        \]
        \begin{enumerate}
            \item $a = b$ and $a \neq \frac{-b}{2} \Rightarrow a \neq 0$.
            The kernel is $x + y + z = 0$ which is a plane, kernel is
            \[
                \spnset{
                    \begin{pmatrix} 
                        -1 \\ 1 \\ 0
                    \end{pmatrix},
                    \begin{pmatrix} 
                        -1 \\ 0 \\ 1
                    \end{pmatrix}
                }
            \]
            Extend this to a basis of $\R^3$
            \[
                \spnset{
                    \begin{pmatrix} 
                        -1 \\ 1 \\ 0
                    \end{pmatrix},
                    \begin{pmatrix} 
                        -1 \\ 0 \\ 1
                    \end{pmatrix},
                    \begin{pmatrix}
                        1 \\ 0 \\ 0
                    \end{pmatrix}
                }   
            \]
            Then we haev that the image,
            \[
                \im(A) = \spnset{\begin{pmatrix}
                    a \\ a \\ a
                \end{pmatrix}}
            \]
            Hence if $c \neq 1$ then $\begin{pmatrix}
                1 \\ c \\ 1
            \end{pmatrix} \notin \im(A)$ and we have no solutions.
            If $c = 1$ a particular solution is $\begin{pmatrix}
                \frac{1}{a} \\ 0 \\ 0
            \end{pmatrix}$ and a general solution is
            \[
                \begin{pmatrix}
                    \frac{1}{a} \\ 0 \\ 0
                \end{pmatrix} 
                + \lambda \begin{pmatrix}
                    -1 \\ 1 \\ 0
                \end{pmatrix}
                + \mu \begin{pmatrix}
                    -1 \\ 0 \\ 1
                \end{pmatrix}
            \]

            \item $a \neq b$, $b = -2a \Rightarrow a \neq 0$. Kenerl is $x = y = z$.
            \[
                \ker(A) = \spnset{\begin{pmatrix}
                    1 \\ 1 \\ 1
                \end{pmatrix}}
            \]
            We extend to a basis of $\R^3$
            \begin{align*}
                \R^3 = \spnset{
                    \begin{pmatrix}
                        1 \\ 1 \\ 1
                    \end{pmatrix},
                    \begin{pmatrix}
                        1 \\ 0 \\ 0
                    \end{pmatrix},
                    \begin{pmatrix}
                        0 \\ 0 \\ 1
                    \end{pmatrix}
                } \\
                \therefore \im(A) = \spnset{
                    \begin{pmatrix}
                        1 \\ -2 \\ 1
                    \end{pmatrix},
                    \begin{pmatrix}
                        -2 \\ 1 \\ 1
                    \end{pmatrix}
                }
            \end{align*}
            Is $\begin{pmatrix}
                1 \\ c \\ 1
            \end{pmatrix} \in \im(A)$ ? 
            \begin{align*}
                \begin{pmatrix}
                    1 \\ c \\ 1
                \end{pmatrix}
                = \lambda \begin{pmatrix}
                    1 \\ -2 \\ 1
                \end{pmatrix}
                + \mu \begin{pmatrix}
                    -2 \\ 1 \\ 1
                \end{pmatrix} \\
                \Rightarrow \mu = 0, \lambda =1 \\
                \Rightarrow c = -2
            \end{align*}

            Hence if $c \neq 2$ there are no solutions. 
            If $c = -2$ a particular solution is $\begin{pmatrix}
                \frac{1}{a} \\ 0 \\ 0
            \end{pmatrix}$. Therefore a general solution is
            \[
                \begin{pmatrix}
                    \frac{1}{a} \\ 0 \\ 0
                \end{pmatrix}
                + \lambda \begin{pmatrix}
                    1 \\ 1 \\ 1
                \end{pmatrix}  
            \]

            \item If $a = b$ and $b = -2a \Rightarrow a = b = 0$ and $A$ is a zero matrix.
            Hence $\ker(A) = \R^3$ and there are no solutions for any $c$.
        \end{enumerate}
    \end{enumerate}
\end{eg}

%%%%%%%%%%%%%%%%%%%%%%
% Lecture 17: 13/11/2023
%%%%%%%%%%%%%%%%%%%%%%
\section{Eigenvalues and Eigenvectors}

\subsection{Preliminaries}
\begin{thm}[Fundamental Theorem of Algebra]
    Consider a polynomail $p(z)$ of degree $m \geq 1$.
    \[
        p(z) = \sum_{r = 0}^{m} c_j z^J, \quad \text{ for } c_j \in \C, \text{ and } c_m \neq 0
    \]
    Then $p(z) = 0$ has precisely $m$ roots in $\C$ counting multiplicity.
    The root $z = w$ has multiplicity $k$ if $(z - w)^k$ is a factor of $p(z)$ but $(z - w)^{k + 1}$ is not a factor.
\end{thm}
\begin{eg}
    \begin{align*}
        p(z) = z^3 - z^2 - z + 1 = (z - 1)^2(z + 1) \\
        p(z) = 0 \Rightarrow z = 1, 1, -1
    \end{align*}
    Where $z = 1$ has multiplicity $2$.
\end{eg}

\begin{defi}[Eigenvectors and Eigenvalues]
    Consider a linear map $\alpha: \C^n \rightarrow \C^n$ with associate $n \times n$ matrix $A$.
    Then $\vec x \neq 0$ is an \emph{eigenvector} of $A$ and $\lambda$ is the associated \emph{eigenvalue} if
    \begin{equation}\label{eq:5-1}
        A \vec x = \lambda \vec x 
    \end{equation}
\end{defi}

\begin{defi}[Characteristic Equation and Polynomial]
    We think of the eigenvectors as the directions which are preserved by the linear map.
    \[
        \text{\eqref{eq:5-1}} \Rightarrow (A - \lambda I) \vec x = \vec 0
    \]
    Therefore, $\vec x \neq 0$ lies in the kernel of $A - \lambda I$.
    \begin{equation}\label{eq:5-2}
        \Rightarrow \det(A - \lambda I) = 0
    \end{equation}
    This it the \emph{characteristic equation} of $A$. Similarly,
    \[
        \mathcal{P}_A(\lambda) \equiv  \det(A - \lambda I)
    \]
    is the \emph{characteristic polynomail} of $A$.

    From \eqref{eq:3-25a},
    \begin{align*}
        P_A(\lambda) &= \det(A - \lambda I) \\
        &= \varepsilon_{J_1, \cdots, J_n}(A_{J_11} - \lambda I_{J_11})\cdots(A_{J_nn} - \lambda I_{J_nn})
    \end{align*}
    Therefore $P_A(\lambda)$ is a polynomial of order $n$.
    \begin{equation}\label{eq:5-3}
        P_A(\lambda) = C_0 + C_1 \lambda + \cdots + C_{n}\lambda^{n}
    \end{equation}
\end{defi}
\begin{remark}\leavevmode
    \begin{enumerate}
        \item \eqref{eq:5-3} and the fundamental theorem of algebra implies that an $n \times n$ matrix has $n$ eigenvalues counting for multiplicity.
        \item If $A$ is a real matrix then $c_i \in \R \Rightarrow \lambda \in \R$ or in complex conjugate pairs.
        \item In \eqref{eq:5-3}
        \begin{align*}
            C_n &= (-1)^n \\
            C_{n-1} &= (-1)^{n-1} (A_{11} + A_{22} + \cdots + A_{nn}) \\
            &= (-1)^{n-1} \tr(A) \\
            \text{also, } \frac{C_{n-1}}{C_n} &= -(\lambda_1 + \cdots + \lambda_n) \tag{by sum \& product of roots} \\
            \therefore & \tr(A) = \text{ sum of eigenvalues} \numberthis\label{eq:5-4}
        \end{align*}

        \item In \eqref{eq:5-3}
        \begin{align*}
            C_0 &= P_A(0) = \det(A) \\
            \text{and, } C_0 &= \lambda_1 \lambda_2 \cdots \lambda_n \\
            \therefore \det A = \text{ product of eigenvalues} \numberthis\label{eq:5-5}
        \end{align*}
    \end{enumerate}
\end{remark}

\begin{defi}[Eigenspace]
    The kernel of the matrix $A - \lambda I$ is the set 
    \[
        \crbkt{\vec 0} \cup \crbkt{\vec x \neq 0 \mid A \vec x = \lambda \vec x}  
    \]
    is a subspace denoted by $E_\lambda$ and is called the \emph{eigenspace} to eigenvalue $\lambda$.
\end{defi}

When we take all the eigenspaces together does it span the entire space? If so, the the matrix can be diagonalised.

\begin{defi}[Algebraic Multiplicity]
    The \emph{algebraic multiplicity} $M_\lambda$ is the multiplicity of $\lambda$ as a root of $P_A(\lambda) = 0$.
    Clearly, 
    \[
        \sum_{\lambda} M_\lambda = n  
    \]
    If $M_\lambda > 1$ the eigenvalue is said to be degenerate.
\end{defi}
\begin{defi}[Geometric Multiplicity]
    The \emph{geometric multiplicity} $m_\lambda$ is the maximum number of linearly independant eigenvectors corresponding to eigenvalue $\lambda$.
    \[
        m_\lambda = \dim(E_\lambda)  
    \]
\end{defi}
\begin{defi}[Defect]
    The \emph{defect} $\bigtriangleup_\lambda$ is 
    \[
        \bigtriangleup_\lambda = M_\lambda - m_\lambda  
    \]
\end{defi}
\begin{proof}
    Can be done! Try when you have time.
\end{proof}
\subsection{Linearly independant Eigenvectors}
\begin{thm}
    Suppose $n \times n$ matrix $A$ has $n$ distinct eigenvalues $\lambda_1, \cdots, \lambda_n$.
    Then, the corresponding eigenvectors
    \begin{equation}\label{eq:5-6}
        \vec x_1, \vec x_2, \cdots, \vec x_n
    \end{equation}
    are linearly independant.
    (w.r.t eigenvector basis, A is diagonal!)
\end{thm}
\begin{proof}
    We will proceed by contradiction. Suppose $\vec x_1, \cdots, \vec x_n$ are linearly dependant.
    Then we can find $d_i$, $i = 1, 2, \cdots, r \leq n$ not all zero such that
    \[
        d_1 \vec x_1 + \cdots + d_n \vec x_n = 0  
    \]
    is the shortest non trivial combination of linearly dependant vectors.
    Now, apply $A - \lambda_1 I$,
    \begin{align*}
        (A - \lambda_1 I)(d_1 \vec x_1 + d_2 \vec x_2 + \cdots + d_n \vec x_n = 0) \\
        \Rightarrow d_1 (\lambda_1 - \lambda_ 1) \vec x_1 + d_2 (\lambda_2 - \lambda_1) \vec x_2 + \cdots + d_r (\lambda_r - \lambda_1) \vec x_r = 0
    \end{align*}
    But $d_1 (\lambda_1 - \lambda_ 1) \vec x_1 = 0$, so we have constructed a shorter combination of linearly independant vectors $\vec x_2 \cdots \vec x_r$, therefore contradiction $\contradiction$. 
\end{proof}

\begin{eg}\leavevmode
    \begin{enumerate}
        \item $A = \begin{pmatrix}
            0 & 1 \\
            -1 & 0
        \end{pmatrix}$
        we can calculate the characteristic polynomial
        \begin{align*}
            \mathcal{P} _A(\lambda) = \lambda^2 + 1 = 0 \Rightarrow \lambda = i, -i  \\
            (A - \lambda_1 I)\vec x = \vec 0 \Rightarrow \begin{pmatrix}
                -i & 1 \\
                -1 & -i
            \end{pmatrix}
            \begin{pmatrix}
                x_1 \\ x_2
            \end{pmatrix}
            = \begin{pmatrix}
                0 \\ 0
            \end{pmatrix}
        \end{align*}
        One particular solution is $\begin{pmatrix}
            x_1 \\ x_2
        \end{pmatrix}=\begin{pmatrix}
            1 \\ i
        \end{pmatrix}$. Therefore
        \[
            E_{i} = \spnset{\begin{pmatrix}
                1 \\ i
            \end{pmatrix}} \Rightarrow m_i = M_i = 1
        \]
        Similarly,
        \begin{align*}
            (A - \lambda_2 I)\vec x = \vec 0 \Rightarrow \begin{pmatrix}
                i & 1 \\
                -1 & i
            \end{pmatrix}
            \begin{pmatrix}
                x_1 \\ x_2
            \end{pmatrix}
            = \begin{pmatrix}
                0 \\ 0
            \end{pmatrix} \\
            \Rightarrow \begin{pmatrix}
                x_1 \\ x_2
            \end{pmatrix}
            = \begin{pmatrix}
                1 \\ -i
            \end{pmatrix} \\
            \therefore E_{-i} = \spnset{\begin{pmatrix}
                1 \\ -i
            \end{pmatrix}} \Rightarrow m_{-i} = M_{-i} = 1
        \end{align*}

        Note that $\begin{pmatrix}
            1 \\ i
        \end{pmatrix}$ and $\begin{pmatrix}
            1 \\ -i
        \end{pmatrix}$ are linearly indeoendant and form a basis of $\C^2$

        \item $A = \begin{pmatrix}
            -2 & 2 & -3 \\
            2 & 1 & -6 \\
            -1 & -2 & 0
        \end{pmatrix}$ similarly, calculate the characteristic polynomail
        \begin{align*}
            \det(A - \lambda I) = 0 &\Rightarrow 45 + 21\lambda - \lambda^2 - \lambda^3 = 0 \\
            &\Rightarrow \lambda_1 = 5, \lambda_2 = \lambda_3 = -3
        \end{align*}

        \begin{description}
            \item[Case: $\lambda_1 = 5$] We have an eigenvector $\begin{pmatrix}
                1 \\ 2 \\ -1
            \end{pmatrix}$ 

            \item[Case: $\lambda_2 = \lambda_3 = -3$]
            \begin{align*}
                (A + 3I)\vec x = \begin{pmatrix}
                    1 & 2 & -3 \\
                    2 & 4 & -6 \\
                    -1 & -2 & 3
                \end{pmatrix}
                \begin{pmatrix}
                    x_1 \\ x_2 \\ x_3
                \end{pmatrix} 
                = \vec 0
            \end{align*}
            Solve to find $\vec x$
            \[
                \vec x = \begin{pmatrix}
                    -2x_2 + 3x_3 \\
                    x_2 \\
                    x_3
                \end{pmatrix}  
            \]
            Therefore we have $2$ linearly independant eigenvectors say
            \[
                \vec x_2 = \begin{pmatrix}
                    -2 \\ 1 \\ 0
                \end{pmatrix}
                \text{ and }
                \vec x_3 = \begin{pmatrix}
                    3 \\ 0 \\ 1
                \end{pmatrix} 
            \]
        \end{description}
        Hence,
        \begin{align*}
            M_5 = 1, m_5 = 1, \bigtriangleup_{5} = 0 \\
            M_{-3} = 2, m_{-3} = 2, \bigtriangleup_{-3} = 0
        \end{align*}

        \item $A = \begin{pmatrix}
            -3 & -1 & -1 \\
            -1 & -3 & 1 \\
            -2 & -2 & 0
        \end{pmatrix}$ solving the characteristic polynomail,
        \begin{align*}
            \mathcal{P} _A(\lambda) = 0 &\Rightarrow - (\lambda + 2)^3 = 0 \\
            &\Rightarrow \lambda = -2, -2, -2
        \end{align*}
        Solving for the eigenvectors
        \begin{align*}
            (A + 2I)\vec x = \vec 0 &\Rightarrow \begin{pmatrix}
                -1 & -1 & 1 \\
                -1 & -1 & 1 \\
                -2 & -2 & 1
            \end{pmatrix}
            \begin{pmatrix}
                x_1 \\ x_2 \\ x_3
            \end{pmatrix} = \vec 0 \\
            &\Rightarrow x_1 + x_2 + x_3 = 0
        \end{align*}
        We have a general solution $\vec x = \begin{pmatrix}
            -x_2 + x_3 \\
            x_2 \\
            x_3
        \end{pmatrix}$
        So, 
        \[
            E_{-2} = \spnset{
                \begin{pmatrix}
                    -1 \\ 1 \\ 0
                \end{pmatrix},                
                \begin{pmatrix}
                    1 \\ 0 \\ 1
                \end{pmatrix}
            }  
        \]
        Hence,
        \[
            M_{-2} = 3, m_{-2} = 2 \Rightarrow \bigtriangleup_{-2} = 1
        \]
        Note that the eigenspaces do not span $\C^3$ and the eigenvectors therefore do not form a basis.

%%%%%%%%%%%%%%%%%%%%%%
% Lecture 18: 13/11/2023
%%%%%%%%%%%%%%%%%%%%%%
        
        \item Reflection matrix \eqref{eq:3-3} - reflection in the plane with normal $\uvec n$.
        \[
            R\vec n = - \vec n
        \]
        Therefore we have eigenvalue $-1$ and eigenspace
        \[
            E_{-1} = \spnset{\vec n}
        \]
        If $\vec p$ lies in the plane, $R\vec p = \vec p$ and hence eigenvalue $1$,
        and the eigenspace is the plane with normal $\vec n$ through $\vec 0$.
        \[
            M_{1} = m_{1} = 2
        \]
        Therefore eigenvectors span whole $\R^3$

        \item Rotation by $\theta$ about $n$ \eqref{eq:3-2}, we have $R \vec n = \vec n$ with eigenvalue $+1$.
        No other real eigenvalues are possible. Turns out that the other eigenvalues are $e^{\pm i \theta}$.
        By theorem \eqref{eq:5-6}, eigenvectors form a basis. (of $\C^3$!)

        \item Shear $A = \begin{pmatrix}
            1 & \mu \\
            0 & 1
        \end{pmatrix}$
        Since this is an upper triangular matrix, we can apply cofactor decompositon to see that the eigenvalues are the diagonal elements.
        Therefore, eigenvalues $\lambda = 1, 1 \Rightarrow M_1 = 2$, with eigenvector $\begin{pmatrix}
            1 // 0
        \end{pmatrix} \rightarrow m_1 = 1$.
        Hence, eigenvectors do not span $\C^2$.
    \end{enumerate}
\end{eg}

\subsubsection{Eigenvector Basis}
If $n \times n$ matrix $A$ has $n$ distinct eigenvalues, then by \eqref{eq:5-6} we have $n$ linearly independant eigenvectors
\[
    \vec v_1, \vec v_2, \cdots, \vec v_n
\]
with respect to this basis, $A$ is diagonal
\[
    A \vec v_i = \lambda_i \vec v_i
\]
i.e the image of $\vec v_i$ is $\lambda_i \times \vec v_i$ and the columns of $A$ are the images of the basis vectors
\begin{equation}\label{eq:5-7}
    \bkt{\begin{array}{ccccc}
        \lambda_1              & 0         & \cdots    & \cdots &  \\
        0                      & \lambda_2 &           &        &\smash{\scalebox{3}{0}} \\
        \vdots                 &           & \lambda_3 &        & \vdots \\
                               &           &           & \ddots & \vdots  \\
        \smash{\scalebox{3}{0}}& \cdots    & \cdots    & \cdots & \lambda_n
    \end{array}} = \diag(\lambda_1, \lambda_2, \cdots, \lambda_n)
\end{equation}
Note that,
\begin{align*}
    \det(A) &= \lambda_1 \lambda_2 \cdots \lambda_n \\
    \tr(A) &= \lambda_1 + \lambda_2 + \cdots + \lambda_n
\end{align*}
and these are invariant with respect to basis

\subsection{Transformation Matrices}
We wish to consider how the components of a vector and a matrix change when when the basis change.
Let $\crbkt{\vec e_1, \vec e_2, \cdots, \vec e_n}$ and $\crbkt{\tilde{\vec{e_1}}, \vec{\tilde{e_2}}, \cdots, \vec{\tilde{e_n}}}$
be two different basis on $\R^n$ or $\C^n$. Then we can write
\begin{equation}\label{eq:5-8}
    \vec{\tilde{e_j}} = \sum_{i=1}^{n}P_{ij} \vec e_i \quad \text{ for } j = 1, \cdots, n
\end{equation}
where $P_{ij}$ is the $i^{th}$ component of $\vec{\tilde{e_j}}$ with respect to the basis $\crbkt{\vec e_1, \vec e_2, \cdots, \vec e_n}$.
So matrix $P$ has as it's columns, the vectors $\vec{\tilde{e_j}}$ relative to $\crbkt{\vec e_1, \vec e_2, \cdots, \vec e_n}$.
\[
    P = \begin{pmatrix}
        \uparrow & \uparrow & & \uparrow \\
        \vec{\tilde{e_1}} & \vec{\tilde{e_2}} & \cdots & \vec{\tilde{e_n}} \\
        \downarrow & \downarrow & & \downarrow
    \end{pmatrix}
\]
i.e; the columns of $P$ are the images of the mapping 
\[
    P(\vec e_i) = \vec{\tilde{e_i}}
\]
Note that this is only true if $P$ has basis $\vec e_1, \vec e_2, \cdots, \vec e_n$. 
Similarly, we can write
\begin{equation}\label{eq:5-9}
    \vec{e_i} = \sum_{k=1}^{n}Q_{ki} \vec{\tilde{e_k}} \quad \text{ for } i = 1, \cdots, n
\end{equation}
with matrix $Q$ such that $Q(\vec{\tilde{e_i}}) = \vec e_i$,
\[
    Q = \begin{pmatrix}
        \uparrow & \uparrow & & \uparrow \\
        \vec{e_1} & \vec{e_2} & \cdots & \vec{e_n} \\
        \downarrow & \downarrow & & \downarrow
    \end{pmatrix}
\]
Placing \eqref{eq:5-9} into \eqref{eq:5-8},
\begin{align*}
    \Rightarrow \vec{\tilde{e_j}} &= \sum_{i=1}^{n}P_{ij}\bkt{\sum_{k=1}^{n}Q_{ki} \vec{\tilde{e_k}}} \\
    &=  \sum_{k=1}^{n} \vec{\tilde{e_k}} \bkt{\sum_{i=1}^{n}P_{ij}Q_{ki}}
\end{align*}
But since $\vec{\tilde{e_j}}$ form a basis,
\begin{align*}
    \Rightarrow \sum_{i=1}^{n}P_{ij}Q_{ki} = \delta_{kj} &\Rightarrow QP = I \\
    &\Rightarrow Q = P^{-1} \numberthis\label{eq:5-10}
\end{align*}
we can also see this from how $P$ and $Q$ map the basis vectors. 

\subsubsection{Transformation Law for Vectors}
With respect to the basis $\vec e_1, \vec e_2, \cdots, \vec e_n$, 
\begin{equation}\label{eq:5-11}
    \vec u = \sum_{i=1}^{n} u_i \vec e_i
\end{equation}
and with respect to another basis $\tilde{\vec{e_1}}, \vec{\tilde{e_2}}, \cdots, \vec{\tilde{e_n}}$
\begin{equation}\label{eq:5-12}
    \vec u = \sum_{i=1}^{n} \tilde{u_i} \tvec e_i
\end{equation}
these are two differnet representations of the same vector, with respect to the different bases.
Use \eqref{eq:5-8} in \eqref{eq:5-12}
\begin{align*}
    \vec u &= \sum_{j=1}^{n} \tilde{u_j} \bkt{\sum_{i=1}^{n}P_{ij} \vec e_i} \\
    &= \sum_{i=1}^{n} \bkt{\sum_{j=1}^{n}P_{ij}\tilde{u_j}} \vec e_i
\end{align*}
Comparing to \eqref{eq:5-11}
\[
    u_i = \sum_{j=1}^{n}P_{ij}\tilde{u_j}
\]
Now, denote $\vec u$ as the vector with respect to $\vec e_1, \cdots, \vec e_n$,
and $\tvec u$ as the same vector with respect to $\tvec e_1, \cdots, \tvec e_n$ then,
\begin{equation}\label{eq:5-13}
    \vec u = P \tvec u \implies \tvec u = P^{-1} \vec u 
\end{equation}
\begin{remark}
    \[
        \tvec e_i = P_{ij} \vec e_j  
    \]
    BUT!
    \[
        \vec u_i = P_{ij} \tvec u_j
    \]
    Basis vectors and coordinates are the ``opposite way round''. 
    This is because in the former case $P$ moves one basis vector to the other in the vector space.
    Therefore, simply applying $P$ to $\vec u$ would move $\vec u$ from where it was in the old basis,
    to it's analagous position in the new basis. We can see this by writing $P \vec u$ in the $\tvec e_i$ basis.
    \[
       \text{($\vec e_i$ basis) } P \vec u \mapsto P^{-1} P \vec u = \vec u  \text{ ($\tvec e_i$ basis)}
    \]
    However, that is not what we want!
    Vectors exist in a vector space independant of the basis. 
    We want to leave the vector unchanged and find it's coordinates in the new basis.
    To do so, we write $\vec u$ in the new basis (simply by changing what the coordinates mean),
    and then apply the inverse map $P^{-1}$ to map back to the original vector.
\end{remark}
\begin{eg}
    Let
    \begin{align*}
        \vec e_1 = \begin{pmatrix} 
            1 \\ 0
        \end{pmatrix} & \qquad
        \tvec e_1 = \begin{pmatrix}
            1 \\ 1
        \end{pmatrix} \\
        \vec e_2 = \begin{pmatrix}
            0 \\ 1
        \end{pmatrix} & \qquad
        \tvec e_1 = \begin{pmatrix}
            -1 \\ 1
        \end{pmatrix}
    \end{align*}
    Therefore, the new basis vectors in terms of the old,
    \[
        \Rightarrow \tvec e_1 = \vec e_1 + \vec e_2 \quad \text{ and } \quad \tvec e_2 = - \vec e_1 + \vec e_2
    \]
    By \eqref{eq:5-8}, $P = \begin{pmatrix}
        1 & -1 \\
        1 & 1
    \end{pmatrix}$, take an arbitrary vector $\vec u$, then
    \[
        \tvec u = P^{-1} \vec u = \frac{1}{2}\begin{pmatrix}
            1 & 1 \\
            -1 & 1
        \end{pmatrix}
        \begin{pmatrix}
            u_1 \\ u_2
        \end{pmatrix}
        = \frac{1}{2} \begin{pmatrix}
            u_1 + u_2 \\
            -u_1 + u_2
        \end{pmatrix}
    \]
    
    Alternatively, by more elementary means,
    \begin{align*}
        \vec u &= u_1 \vec e_1 + u_2 \vec e_2 \\
        &= u_1 \frac{\tvec e_1 - \tvec e_2}{2} + u_2 \frac{\tvec e_1 + \tvec e_2}{2} \\
        &= \frac{1}{2}(u_1 + u_2) \tvec e_1 + \frac{1}{2}(-u_1 + u_2) \tvec e_2
    \end{align*}
\end{eg}

\subsubsection{Transformation Law for Matrices}
Consider a linear map $\alpha: \C^n \rightarrow \C^n$ and the corresponding $n \times n$ matrix $A$ such that
\[
    \vec u' = \alpha(\vec u) \Rightarrow \vec u' = A \vec u
\]
Denote $\vec u$ and $\vec u'$ as being expressed with respect to the basis $\vec e_1, \cdots, \vec e_n$ (same basis in both spaces),
and $\tvec u$ and $\tvec u'$ as being expressed with respect to $\tvec e_1, \cdots \tvec e_n$.
Using \eqref{eq:5-13}
\begin{align*}
    \vec u' = A \vec u &\Rightarrow P \tvec u' = A P \tvec u \\
    &\Rightarrow \tvec u' = P^{-1} A P \tvec u \\
    &\Rightarrow \tvec u' = \tilde{A} \tvec u
\end{align*}
Hence, with respect to the $\tvec e_1, \cdots \tvec e_n$ basis the matrix becomes
\begin{equation}\label{eq:5-14}
    \tilde{A} = P^{-1} A P
\end{equation}

%%%%%%%%%%%%%%%%%%%%%%
% Lecture 19: 17/11/2023
%%%%%%%%%%%%%%%%%%%%%%

\begin{eg}
    Shear $S_\lambda = \begin{pmatrix}
        1 & \lambda & 0 \\
        0 & 1 & 0 \\
        0 & 0 & 1
    \end{pmatrix}$ with respect to $\vec e_1, \cdots, \vec e_n$. 
    Then choose a new rotated basis such that
    \begin{align*}
        \tvec e_1 &= \vec e_1 \cos \theta + \vec e_2 \sin \theta \\
        \tvec e_2 &= - \vec e_1 \sin \theta + \vec e_2 \cos \theta \\
        \tvec e_2 &= \vec e_3
    \end{align*}
    \[
        \Rightarrow P = \begin{pmatrix}
            \cos \theta & -\sin \theta & 0 \\
            \sin \theta & \cos \theta & 0 \\
            0 & 0 & 1
        \end{pmatrix}
        \Rightarrow P^{-1} = \begin{pmatrix}
            \cos \theta & \sin \theta & 0 \\
            -\sin \theta & \cos \theta & 0 \\
            0 & 0 & 1
        \end{pmatrix}
    \]
    Now use \eqref{eq:5-14}
    \[
        \tilde{S_\lambda} = \begin{pmatrix}
            1 + \lambda \sin \theta \cos \theta & \lambda \cos^2 \theta & 0 \\
            -\lambda \sin^2 \theta & 1 - \lambda \sin \theta \cos \theta & 0 \\
            0 & 0 & 1
        \end{pmatrix}
    \]
\end{eg}

\subsubsection*{General case}
More generally, consider a map
\begin{align*}
    \alpha : \C^m &\rightarrow \C^n \\
    \vec x &\mapsto \vec x' = A\vec x
\end{align*}
Hence $\vec x \in \C^m$ and $\vec x' \in \C^n$ and $A$ is an $n \times m$ matrix.
Say our basis in $\C^m$ is $\vec e_1, \cdots, \vec e_m$ and the basis in $\C^n$ is $\vec f_1, \cdots, \vec f_n$.
Now we change basis to  $\tvec e_1, \cdots, \tvec e_m$ and $\tvec f_1, \cdots, \tvec f_n$.
From \eqref{eq:5-13},
\begin{align*}
    \vec x = P \tvec x \qquad \text{ for } P \text{ a } m \times m \text{ matrix} \\
    \vec x' = R \tvec x' \qquad \text{ for } R \text{ a } n \times n \text{ matrix}
\end{align*}
Then,
\begin{align*}
    \vec x' = A \vec x &\Rightarrow R \tvec x' = AP \tvec x \\
    &\Rightarrow \tvec x' = R^{-1}AP \tvec x
\end{align*}
\begin{equation}\label{eq:5-15}
    \therefore \tilde{A} = R^{-1}AP
\end{equation}

\begin{eg}
    Consider $\alpha: R^3 \rightarrow R^2$, with respect to standard basis vectors in both spaces,
    \[
        A = \begin{pmatrix}
            2 & 3 & 4 \\
            1 & 6 & 3
        \end{pmatrix}  
    \]
    use a new basis $\begin{pmatrix}
        2 \\ 1
    \end{pmatrix}$, $\begin{pmatrix}
        1 \\ 5
    \end{pmatrix}$ in $\R^2$ but keep the basis in $\R^3$. 
    Matrix $P$ relates basis in $\R^3$ so $P = I$ since it is unchanged.
    Matrix $R$ relates basis in $\R^2$ so $P$ so from \eqref{eq:5-8},
    \[
        R = \begin{pmatrix}
            2 & 1 \\
            1 & 5
        \end{pmatrix}
        \Rightarrow R^{-1} = \frac{1}{9}\begin{pmatrix}
            5 & -1 \\
            -1 & 2
        \end{pmatrix}
    \]
    \begin{align*}
        \Rightarrow \tilde{A} &= \frac{1}{9}\begin{pmatrix}
            5 & -1 \\
            -1 & 2
        \end{pmatrix}\begin{pmatrix}
            2 & 3 & 4 \\
            1 & 6 & 3
        \end{pmatrix} \\
        &= \begin{pmatrix}
            1 & 1 & \frac{17}{9} \\[10pt]
            0 & 1 & \frac{2}{9}
        \end{pmatrix}
    \end{align*}

    We may once again check by elementary methods, 
    if we write $\tvec f_1 = \begin{pmatrix}
        2 \\ 1
    \end{pmatrix}$, $\tvec f_2 = \begin{pmatrix}
        1 \\ 5
    \end{pmatrix}$, then
    \begin{alignat*}{2}
        \vec e_1 = \tvec e_1 & \longmapsto && \quad 2 \vec f_1 + \vec f_2 = 1 \tvec f_1 + 0 \tvec f_2 \\
        \vec e_2 = \tvec e_2 & \longmapsto && \quad 3 \vec f_1 + 6 \vec f_2 = 1 \tvec f_1 + 1 \tvec f_2 \\
        \vec e_3 = \tvec e_3 & \longmapsto && \quad 4 \vec f_1 + 3 \vec f_2 = \cdots = \frac{17}{9} \tvec f_1 + \frac{2}{9} \tvec f_2
    \end{alignat*}
\end{eg}

\subsection{Similar Matrices}
\begin{defi}
    Two $n \times n$ matrices $A$ and $B$ are similar if there exists an invertible matrix $P$ such that
    \begin{equation}\label{eq:5-16}
        B = P^{-1}AP
    \end{equation}
    In essence, these two matrices describe the same linear map but with different basis.
    From \eqref{eq:5-14}, matrices representing the same map with respect to different basis are similar.
\end{defi}
\begin{remark}\leavevmode
    \begin{enumerate}
        \item Similar matrices have the same determinant.
        \[
            \det(B) = \det(P^{-1}AP) = \det(P^{-1}) \det(A) \det(P) = \det(A)
        \]
        by \eqref{eq:3-31}.
        
        \item Also have the same trace
        \begin{align*}
            \tr(B) = B_{ii} &= P^{-1}_{ij}A_{jk}P_{ki} \\
            &= P_{ki}P^{-1}_{ij}A_{jk} \\
            &= \delta_{kj}A_{jk} \\
            &= A_{kk} = \tr(A)
        \end{align*}

        \item Similar matrices have the same characteristic equation.
        \begin{align*}
            P_B(\lambda) = \det(B - \lambda I) &= \det(P^{-1}AP - \lambda I) \\
            &= \det(P^{-1}) \det(AP - \lambda P) \\
            &= \det(P^{-1}) \det(P) \det(A - \lambda I) \\
            &= \det(A - \lambda I) = P_A(\lambda)
        \end{align*}
        Also, note that this implies (i) and (ii) by \eqref{eq:5-4} and \eqref{eq:5-5}
    \end{enumerate}
\end{remark}
\subsection{Diagonalisable Matrices}
A matrix $A$ ($n \times n$) is diagonalise if it is similar to a diagonal matrix. 
From \eqref{eq:5-7} this is equivalent to saying that the eigenvectors form a basis.
The condition that $A$ has $n$ distinct eigenvalues is sufficient but not necessary for $A$ to be diagonalisable.

\begin{eg} Section 5.2 example 2
    \[
        A = \begin{pmatrix}
            -2 & 2 & -3 \\
            2 & 1 & -6 \\
            -1 & 2 & 0
        \end{pmatrix}, 
        \Rightarrow \lambda = 5, -3, -3
    \]
    has three independant eigenvectors.
    \[
        \begin{pmatrix}
            1 \\ 2 \\ -1
        \end{pmatrix},
        \begin{pmatrix}
            -2 \\ 1 \\ 0
        \end{pmatrix},
        \begin{pmatrix}
            3 \\ 0 \\ 1
        \end{pmatrix}
        \Rightarrow P = \begin{pmatrix}
            1 & -1 & 3 \\
            2 & 1 & 0 \\
            -1 & 0 & 1
        \end{pmatrix}
    \]
    We find $P^{-1} = \frac{1}{8} \begin{pmatrix}
        1 & 2 & -3 \\
        -2 & 4 & 6 \\
        1 & 2 & 5
    \end{pmatrix}$ Hence, $\tilde{A} = P^{-1} A P$
    \begin{align*}
        \Rightarrow \frac{1}{8} \begin{pmatrix}
            1 & 2 & -3 \\
            -2 & 4 & 6 \\
            1 & 2 & 5
        \end{pmatrix}\begin{pmatrix}
            -2 & 2 & -3 \\
            2 & 1 & -6 \\
            -1 & 2 & 0
        \end{pmatrix}\begin{pmatrix}
            1 & -1 & 3 \\
            2 & 1 & 0 \\
            -1 & 0 & 1
        \end{pmatrix}\\
        = \begin{pmatrix}
            5 & 0 & 0 \\
            0 & -3 & 0 \\
            0 & 0 & -3
        \end{pmatrix}
    \end{align*}
    and $A$ is diagonalisable.
\end{eg}

\begin{thm}
    Let $\lambda_1, \cdots, \lambda_r$ with $r \leq n$ be the distinct eigenvalues of $A$,
    and $B_1, B_2, \cdots, B_r$ be the basis of the corresponding eigenspaces, $E_{\lambda_1}, E_{\lambda_2}, \cdots, E_{\lambda_r}$.
    Then the set 
    \begin{equation}\label{eq:5-17}
        B = \bigcup_{i=1}^{r}B_i
    \end{equation}
    is linearly independant.
\end{thm}
\begin{proof}
    Write $B_1 = \crbkt{\vec x_1^1, \vec x_2^1, \cdots, \vec x^1_{m_{\lambda_1}}}$, where $m_{\lambda_1}$ is the $\dim(E_{\lambda_1})$ etc.
    Consider a general linear combination of all elements of $B$.
    \begin{equation}\label{eq:5-18}
        \sum_{i=1}^{r} \sum_{j=1}^{m_{\lambda_i}} \alpha_{ij} \vec x^i_j = \vec 0
    \end{equation}
    for some coefficients $\alpha_{ij}$. Here we the first sum sums over each eigen space and the inner sum sums over all basis vectors in $B_i$.
    Now apply the matrix
    \[
        \prod_{k = 1, \cdots, \overline{t}, \cdots, r} (A - \lambda_kI) \quad \text{ for some } t
    \]
    to \eqref{eq:5-18}
    \[
        \Rightarrow \sum_{j=1}^{m_{\lambda_t}} \alpha_{tj} \prod_{k = 1, \cdots, \overline{t}, \cdots, r} (\lambda_t - \lambda_k) \vec x^t_j = 0
    \]
    We select only the $t^{th}$ eigenspace, since all other terms will contain a factor of $\lambda_i - \lambda_i$ in the product.
    The $\vec x_j^t$ are linearly independant themsleves since $B_t$ is a basis.
    Therefore $\alpha_{tj} = 0$ for all $j$.
    Repeat this for all other values of $t$ to get that all $\alpha_{ij} = 0$.
    Therefore $B$ is linearly independant.
\end{proof}

\subsubsection*{Test for Diagonalisability}
\begin{enumerate}
    \item Solve for $P_A(\lambda) = 0$ for distinct $\lambda_1, \cdots, \lambda_r$
    \item Find the basis $B_i$ for each eigenspace $E_{\lambda_i}$
    \item Calculate the size of $B$ by summing over all $B_i$
    \[
        \abs{B} = \sum_{i = 1}^{r} m_{\lambda_i}  
    \]
\end{enumerate}
If $\abs{B}$ = n then $A$ is diagonalisable. Otherwise $A$ is not diagonalisable. 
For instance, in section 5.2 example 3
\[
    \sum_{i=1}^{3} m_{\lambda_i} = m_{-3} = 2 < 3  
\]

%%%%%%%%%%%%%%%%%%%%%%
% Lecture 20: 20/11/2023
%%%%%%%%%%%%%%%%%%%%%%
\subsection{Cannonical form of $2 \times 2$ matrices}
Any $2 \times 2$ matrix $A$ is similar to one of
\begin{equation}\label{eq:5-19}
    \begin{array}{ccc}
        \text{(i)} \begin{pmatrix}
        \lambda_1 & 0 \\
        0 & \lambda_2
        \end{pmatrix}
        &
        \text{(ii)} \begin{pmatrix}
        \lambda & 0 \\
        0 & \lambda
        \end{pmatrix}
        &
        \text{(iii)} \begin{pmatrix}
        \lambda & 1 \\
        0 & \lambda
        \end{pmatrix}
    \end{array}
\end{equation}

\begin{proof}\leavevmode
    \begin{enumerate}
        \item If $A$ has distinct eigenvalues then the eigenvectors are linearly independant by \eqref{eq:5-6}.
        Use \eqref{eq:5-16} with $P$ formed from the eigenvectors as columns.
        
        \item If $\lambda_1 = \lambda_2 = \lambda$ and $\dim(E_\lambda) = 2$ then write $E_\lambda = \spnset{\vec u, \vec v}$.
        use $\vec u$ and $\vec v$ as a basis of $\C^2$ and change to this basis.
        \[
            \text{\eqref{eq:5-14}} \Rightarrow \tilde{A} = P^{-1} A P = \begin{pmatrix}
                \lambda & 0 \\
                0 & \lambda
            \end{pmatrix} = \lambda I
        \]
        Therefore $A = \lambda I$ with respect to any basis. (Isotropic)
        
        \item If $\lambda_1 = \lambda_2 = \lambda$ and $\dim(E_\lambda) = 1$. Write $E_\lambda = \spnset{\vec v}$. 
        Now, extend to a basis of $\C^2 = \spnset{\vec v, \vec w}$ where $\vec w \in \C^2 \setminus E_\lambda$.
        \[
            A \vec w = \vec w' \in \C^2 \Rightarrow \vec w' = \alpha \vec v + \beta \vec w
        \]
        Hence, change to a basis $\crbkt{\vec v, \vec w}$
        \[
            \tilde{A} = P^{-1} A P = \begin{pmatrix}
                \lambda \alpha \\
                0 & \beta
            \end{pmatrix}
        \]
        However, since $A$ and hence $\tilde{A}$ have eigenvalue $\lambda$ with multiplicity $2$ and $\tilde{A}$ is upper triangular
        \[
            \Rightarrow \beta = \lambda
        \]
        Further, introduce $\vec u = (\tilde{A} - \lambda I)\vec w$, $\vec u \neq 0$ because $\vec w \notin E_\lambda$.
        \begin{align*}
            (\tilde{A} - \lambda I)\vec u = (\tilde{A} - \lambda I)^2\vec w = \begin{pmatrix}
                0 & \alpha \\
                0 & 0
            \end{pmatrix}\begin{pmatrix}
                0 & \alpha \\
                0 & 0
            \end{pmatrix} \vec w = 0_{2 \times 2} \vec w = \vec 0
        \end{align*}
        Therefore $\vec u$ is an eigenvector of $\tilde{A}$ with eigenvalue $\lambda$. Expanding
        \[
            \vec u = \tilde{A}\vec w - \lambda \vec w \Rightarrow \tilde{A}\vec w = \vec u + \lambda \vec w
        \]
        Now we use basis $\crbkt{\vec u, \vec w}$ and with respect to this basis the matrix becomes
        \[
            \begin{pmatrix}
                \lambda & 1 \\
                0 & \lambda
            \end{pmatrix}  
        \]
        Therefore this is a $2$-step process. P sends the basis to $\crbkt{\vec v, \vec w}$ and $Q$ sends the basis to $\crbkt{\vec u, \vec w}$.
        So the similarity transform is
        \[
            Q^{-1}(P^{-1}AP)Q = (PQ)^{-1}A(PQ) = \begin{pmatrix}
                \lambda & 1 \\
                0 & \lambda
            \end{pmatrix}  
        \]
    \end{enumerate}
\end{proof}

\subsubsection*{$n \times n$ case:}
The cannonical or \emph{Jordan normal} form exists for any complex $n \times n$ matrix, and is $\tilde{A} = P^{-1}AP$
where $\tilde{A}_{\alpha \alpha} = \lambda_\alpha$ and $\tilde{A}_{\alpha \alpha + 1} = 0 \text{ or } 1$, $\tilde{A}_{ij} = 0$ otherwise.
\[
    \bkt{\begin{array}{ccccc}
        \lambda_1              & (0,1)     & \cdots    & \cdots    &  \\
        0                      & \lambda_2 & (0,1)     &           &\smash{\scalebox{3}{0}} \\
        \vdots                 &           & \lambda_3 & (0,1)     & \vdots \\
                               &           &           & \ddots    & (0,1)  \\
        \smash{\scalebox{3}{0}}& \cdots    & \cdots    & \cdots    & \lambda_n
    \end{array}}    
\]
We see that there is a basis where a matrix takes it's cleanest form! All the information lives on the diagonal and its super diagonal.
\begin{eg}
    Section 5.2 example 3
    \[
        A = \begin{pmatrix}
            -3 & -1 & 1 \\
            -1 & -3 & 1 \\
            2 & -2 & 0
        \end{pmatrix} \Rightarrow \lambda = -2, -2, -2
    \]
    We have eigenvectors
    \[
        \begin{pmatrix}
            -1 \\ 1 \\ 0
        \end{pmatrix}, 
        \begin{pmatrix}
            1 \\ 0 \\ 1
        \end{pmatrix}
    \]
    Take $\vec w = \begin{pmatrix}
        1 \\ 0 \\ 0
    \end{pmatrix}$ for $\vec w \notin E_\lambda$. Find $\vec u$
    \[
        \vec u = (A + 2I)\vec w = \cdots = \begin{pmatrix}
            -1 \\ -1 \\ -2
        \end{pmatrix}
    \]
    Note that $A \vec u = -2 \vec u$ and $A \vec w = \vec u - 2 \vec w$
    Take as a basis, $\vec u, \vec w, \vec v$ where $\vec v$ is an eigenvector which is linearly independant of $\vec u$.
    For example, $\vec v = \begin{pmatrix}
        1 \\ 0 \\ 1
    \end{pmatrix}$
    \[
        \Rightarrow P = \begin{pmatrix}
            -1 & 1 & 1 \\
            -1 & 0 & 0 \\
            -2 & 0 & 1
        \end{pmatrix}  
    \]
    then the Jordan normal form is $\tilde{A} = P^{-1} A P$
    \[
        \tilde{A} = \begin{pmatrix}
            -2 & 1 & 0 \\
            0 & -2 & 0 \\
            0 & 0 & -2
        \end{pmatrix}  
    \]
\end{eg}
\subsection{Cayley-Hamilton Theorem}
\begin{thm}[Cayler-Hamilton Theorem]
    Every $n \times n$ matrix satisfies its own characteristic equation.
    \begin{equation}\label{eq:5-20}
        \mathcal{P}_A(A) = 0
    \end{equation}
\end{thm}
\begin{proof}
    We will only cover the case for diagonalisable matrices. Hence, there exists $P$ such that
    \[
        D = \diag(\lambda_1, \lambda_2, \cdots, \lambda_n) = P^{-1}AP
    \]
    Note that
    \[
        D^i = P^{-1}AP \cdots P^{-1}AP = P^{-1}A^iP
    \]
    Hence,
    \[
        \mathcal{P}_D(D) = \mathcal{P}_D(P^{-1}AP) = P^{-1}\mathcal{P}_D(A)P  
    \]
    But similar matrices have the same characteristic polynomial, $\mathcal{P}_D(\lambda) = \mathcal{P}_A(\lambda)$.
    \[
        \therefore \mathcal{P}_A(D) = P^{-1}\mathcal{P}_A(A)P  
    \]
    Also, $D^i = \diag(\lambda_1^i, \lambda_2^i, \cdots, \lambda_n^i)$
    \begin{align*}
        \Rightarrow \mathcal{P}_A(D) &= \diag(\mathcal{P}_A(\lambda_1), \mathcal{P}_A(\lambda_2), \cdots, \mathcal{P}_A(\lambda_n)) \\
        &= \diag(0, 0, \cdots, 0) \\
        \therefore \quad& P^{-1}\mathcal{P}_A(A)P = 0_{n \times n} \\
        \Rightarrow \quad& \mathcal{P}_A(A) = 0
    \end{align*}
\end{proof}
\begin{remark}\leavevmode
    \begin{enumerate}
        \item For $2 \times 2$ matrices which are similar to $B = \begin{pmatrix}
            \lambda & 1 \\
            0 & \lambda
        \end{pmatrix}$ see that
        \[
            \mathcal{P}_B(z) = \det(B - zI) = (\lambda - z)^2
        \]
        then,
        \[
            \mathcal{P}_B(B) = (\lambda I - B)^2 = \begin{pmatrix}
                0 & -1 \\
                0 & 0
            \end{pmatrix}\begin{pmatrix}
                0 & -1 \\
                0 & 0
            \end{pmatrix} = 0
        \]
        Hence, all $2 \times 2$ matrices satisfy Cayley-Hamilton.
        
        \item Any function that can be expressed as a power series can take a matrix as an input.
        For example, let $A$ be an $n \times n$ matrix,
        \[
            \exp(A) = I + A + \frac{A^2}{2} + \frac{A^3}{3!} + \cdots + \frac{A^n}{n!} + \cdots
        \]
        which converges for all matrices! If $A$ is diagonalisable $D = P^{-1}AP$ then
        \begin{align*}
            P^{-1} \exp(A) P &= P^{-1}IP + P^{-1}AP + \cdots + \frac{P^{-1}A^nP}{n!} + \cdots \\
            &= I + D + \cdots + \frac{D^n}{n!} + \cdots \\
            &= \diag(e^{\lambda_1}, e^{\lambda_2}, \cdots, e^{\lambda_n})
        \end{align*}
        \[
            \therefore \quad \exp(A) = P \diag(e^{\lambda_1}, e^{\lambda_2}, \cdots, e^{\lambda_n}) P^{-1}
        \]
        Solution of $\diff{\vec x}{t} = A \vec x \Longrightarrow \vec x(t) = \exp(At)\vec x_0$.
        Likewise, 
        \[
            \log (I - A) = - A - \frac{A^2}{2} - \frac{A^3}{3} - \cdots  
        \]
        where we have a ``spectral radius'' which is similar to the radius of convergence.

        \item If $A^{-1}$ exists
        \begin{align*}
            A^{-1}\mathcal{P}_A(A) &= A^{-1}(c_0 I + c_1 A + \cdots + c_n A^n) \\
            &= 0
        \end{align*}
        Since $\det A \neq 0 \Longrightarrow c_0 \neq 0$
        \[
            \Rightarrow A^{-1} = \frac{-1}{c_0}(c_1 I + c_2 A + \cdots + c_n A^{n-1})
        \]
    \end{enumerate}
\end{remark}

%%%%%%%%%%%%%%%%%%%%%%
% Lecture 21: 22/11/2023
%%%%%%%%%%%%%%%%%%%%%%
\subsection{Eigenvalues and Eigenvectors of a hermition matrix}
\begin{thm}
    Eigenvalues of a Hermition Matrix are real
    \begin{equation}\label{eq:5-21}
        H = H^\dagger \Rightarrow \lambda \in \R
    \end{equation}
    \begin{proof} Let $H$ be a hermition matrix with eigenvalue $\lambda$ and eigenvector $\vec v \neq 0$.
        \begin{align*}
            \Rightarrow H\vec v = \lambda \vec v \\
            \Rightarrow \vec v^\dagger H\vec v = \vec v^\dagger \lambda \vec v \tag{($\star$)}
        \end{align*}
        We can take the hermition conjugate of the LHS of ($\star$)
        \begin{align*}
            (vec v^\dagger H\vec v)^\dagger &= vec v^\dagger H^\dagger \vec v \\
            &= vec v^\dagger H\vec v
        \end{align*}
        We apply the same to the RHS of ($\star$) to see that
        \begin{align*}
            (\vec v^\dagger \lambda \vec v)^\dagger &= (\lambda \vec v^\dagger \vec v)^\dagger \\
            &= \lambda^* \vec v^\dagger \vec v \tag{($\star^\dagger$)}
        \end{align*}
        Take ($\star$) $-$ ($\star^\dagger$)
        \[
            (\lambda - \lambda^\star)\vec v^\dagger \vec v = 0
        \]
        Since, $\vec v \neq 0 \Rightarrow \vec v^\dagger \neq 0 \Rightarrow \vec v^\dagger \vec v \neq 0$
        \[
            \therefore \lambda = \lambda^\star
        \]
    \end{proof}
\end{thm}

Recall from section 2.6.4,
\[
    \vec v^\dagger \vec v = \braket{\vec v}{\vec v}
\]

\begin{thm}
    The eigenvectors of a Hermition Matrix corresponding to distinct eigenvalues are orthogonal.
    Let $\vec v_i$ and $\vec v_j$ be eigenvectors corresponding to different eigenvalues, then
    \begin{equation}\label{eq:5-22}
        \braket{\vec v_i}{\vec v_j} = 0
    \end{equation}
\end{thm}
\begin{proof}
    For $H$ a hermition matris, let
    \begin{enumerate}
        \item $H \vec v_i = \lambda_i \vec v_i$
        \item $H \vec v_j = \lambda_j \vec v_j$
    \end{enumerate}
    We pre multiply (i) and (ii) by the hermition conjugate of the other eigenvector,
    \begin{enumerate}
        \item $\vec v_j^\dagger H \vec v_i = \vec v_j^\dagger \lambda_i \vec v_i$
        \item $\vec v_i^\dagger H \vec v_j = \vec v_i^\dagger \lambda_j \vec v_j$
    \end{enumerate}
    Then consider
    \[
        (\vec v_j^\dagger H \vec v_i)^\dagger = \vec v_i^\dagger H \vec v_j
    \]
    Therefore,
    \begin{align*}
        (\vec v_j^\dagger \lambda_i \vec v_i)^\dagger = \vec v_i^\dagger \lambda_j \vec v_j \\
        \Rightarrow \lambda_i^* \vec v_i^\dagger\vec v_j = \lambda_j \vec v_i^\dagger\vec v_j \tag{by \eqref{eq:5-21}}\\
        \Rightarrow (\lambda_i - \lambda_j)\vec v_i^\dagger\vec v_j = 0
    \end{align*}
    Therefore, if $\lambda_i - \lambda_j \neq 0$ i.e, the eigenvalues are distinct, then
    \[
        \vec v_i^\dagger\vec v_j = \braket{\vec v_i}{\vec v_j} = 0
    \]
\end{proof}

Therefore, if a hermition matrix $H$ has $n$ distinct eigenvalues, then the eigenvectors form an orthonormal basis:
\begin{enumerate}
    \item form a basis by \eqref{eq:5-6}
    \item orthogonal by \eqref{eq:5-22}
    \item we can normalise each eigenvector By
    \[
        \vec v_i \rightarrow \frac{\vec v_i}{\braket{\vec v_i}{\vec v_i}}
    \]
    so that
    \[
        \vec v_i^\dagger\vec v_j = \delta_{ij}
    \]
\end{enumerate}

\subsubsection{Gram-Schmidt Orthogonalisation}
Suppose we have a set $B = \crbkt{\vec w_1, \vec w_2, \cdots, \vec w_r}$ of linearly independant but not necessarily orthogonal vectors. 
We can form an orthogonal set $\tilde{B} = \crbkt{\vec v_1, \vec v_2, \cdots, \vec v_r}$.

Define the projection of $\vec w$ onto $\vec v$ by
\[
    \wp_{\vec{v}}(\vec w) = \frac{\braket{\vec v}{\vec w}}{\braket{\vec v}{\vec v}} \vec v   
\]
Now, we construct $\tilde{B}$ iteratively:
\begin{enumerate}
    \item $\vec v_1 = \vec w_1$
    \item $\vec v_2 = \vec w_2 - \wp_{v_1}(\vec w_2)$
    \item $\vec v_3 = \vec w_3 - \wp_{v_1}(\vec w_3) - \wp_{v_2}(\vec w_3)$
    \item[{$\vdots$}] 
    \item[(r)] $\vec v_r = \vec w_r - \sum_{j=1}^{r-1}\wp_{v_j}(\vec w_r)$
\end{enumerate}
At each stage, $\vec v_k$ is defined to be $\vec w_k$ minus the projection of $\vec w_k$ onto the subspace
\[
    \spnset{\vec v_1, \vec v_2, \cdots, \vec v_{k-1}}
\]
\subsubsection{Unitary Transformation}
Suppose $U$ is the transformation matrix between one orthonormal basis and a new otrhonormal basis
\[
    \crbkt{\vec u_1, \vec u_2, \cdots, \vec u_n}
\]
where
\[
    \braket{\vec u_i}{\vec u_i} = \delta_{ij}
\]
Then, 
\begin{align*}
    U &= \begin{pmatrix}
        \uparrow & \uparrow & & \uparrow \\
        \vec u_1 & \vec u_2 & \cdots & \vec u_n \\
        \downarrow & \downarrow & & \downarrow
    \end{pmatrix} \\[10pt]
    &= \begin{pmatrix}
        (\vec u_1)_1 & (\vec u_2)_1 & \cdots & (\vec u_3)_1 \\
        (\vec u_1)_2 & (\vec u_2)_2 & \cdots & (\vec u_n)_2 \\
        \vdots & \vdots & & \vdots \\
        (\vec u_1)_n & (\vec u_2)_n & \cdots & (\vec u_n)_n
    \end{pmatrix}
\end{align*}
Therefore
\begin{align*}
    (U^\dagger U)_{ij} = (U^\dagger)_{ik} U_{kj} &= U^*_{ki} U_{kj} \\
    &= (u^*_i)_k (u_j)_k \\
    &= \braket{\vec u_i}{\vec u_i} = \delta_{ij}
\end{align*}
Therefore,
\[
    U^\dagger U = I \quad \therefore U \text{ is unitary}
\]

\subsubsection{Diagonalisation of $n \times n$ Hermition Matrix}
\begin{thm}
    A $n \times n$ Hermition Matrix has $n$ orthogonal eigenvectors and is therefore diagonalisable
    \begin{equation}\label{eq:5-23}
        H = P^{-1}DP \qquad \text{ for } D = \diag(\lambda_1, \cdots, \lambda_n)
    \end{equation}
\end{thm}
\begin{proof} (Non - examinable)
    Let $\lambda_1, \lambda_2, \cdots, \lambda_r$ be the distinct eigenvalues of $H$ $(r \leq n)$ with corresponding set of orthonormal (by \eqref{eq:5-22}) eigenvectors.
    \[
        B = \crbkt{\vec v_1, \cdots, \vec v_r} 
    \]
    Extend $B$ to a basis of the whole of $\C^n$
    \[
        B = \crbkt{\vec v_1, \cdots, \vec v_r, \vec w_1, \cdots, \vec w_{n-r}}  
    \]
    and then use Gram-Schmidt to create an orthonormal basis of $\C^n$
    \[
        \tilde{B} = \crbkt{\vec v_1, \cdots, \vec v_r, \vec u_1, \cdots, \vec u_{n-r}}
    \]
    now write 
    \[
        P = \begin{pmatrix}
            \uparrow & & \uparrow & \uparrow & & \uparrow \\
            \vec v_1 & \cdots & \vec v_r & \vec u_1 & \cdot & \vec u_{n-r} \\
            \downarrow & & \downarrow & \downarrow & & \downarrow
        \end{pmatrix}_{n \times n} 
    \]
    Since this is a transformation matrix between orthonormal basis,
    \[
        P^{-1} = P^\dagger \Rightarrow P^{-1}HP = P^\dagger H P =
    \]
    \[
        \begin{pmatrix}
            \lambda_1 & 0         & \cdots & 0         & 0          & \cdots & \cdots & 0 \\
            0         & \lambda_2 &        & \vdots    & \vdots     & \ddots &        & \vdots \\
            \vdots    &           & \ddots & \vdots    & \vdots     &        & \ddots & \vdots \\
            0         & \cdots    & \cdots & \lambda_r & 0          & \cdots & \cdots & 0 \\
            0         & \cdots    & \cdots & 0         & c_{11}     & \cdots & \cdots & c_{1, n-r} \\  
            \vdots    & \ddots    &        & \vdots    & \vdots     & \ddots &        & \vdots \\
            \vdots    &           & \ddots & \vdots    & \vdots     &        & \ddots & \vdots \\
            0         & \cdots    & \cdots & 0         & c_{n-r, 1} & \cdots & \cdots & c_{n-r, n-r}
        \end{pmatrix}
    \]
    The resulting matrix has block zeros since $\tilde{B}$ is orthonormal so
    \[
        \vec{u}_i^\dagger H \vec v_j = \lambda \vec{u}_i^\dagger \vec v_j = 0
    \]
    Furthermore,
    \[
        (P^\dagger H P)^\dagger = P^\dagger H^\dagger P = P^\dagger H P
    \]
    Therefore, $C$ is a $(n-r) \times (n-r)$ hermition matrix as well.
    The eigen values of $C$ are also eigenvalues of $H$ because
    \begin{align*}
        \det(H - \lambda I) = \det(P^\dagger H P - \lambda I) \\
        = (\lambda_1 - \lambda) \cdots (\lambda_r - \lambda) \det(C - \lambda I)
    \end{align*}
    by cofactor decomposition.
    Now suppose eigenvalues of $C$ are distinct, (otherwise repeat the process).
    Hence we have $n - r$ orthonormal eigenvectors $\vec{\mathcal{W}}_j$ for $j = 1, \cdots, n-r$ by \eqref{eq:5-22} and let
    \[
        Q = \begin{pmatrix}
            1      & 0      & \cdots & 0      & 0                   & \cdots              & \cdots & 0 \\
            0      & 1      &        & \vdots & \vdots              & \ddots              &        & \vdots \\
            \vdots &        & \ddots & \vdots & \vdots              &                     & \ddots & \vdots \\
            0      & \cdots & \cdots & 1      & 0                   & \cdots              & \cdots & 0 \\
            0      & \cdots & \cdots & 0      & \uparrow            & \uparrow            &        & \uparrow \\  
            \vdots & \ddots &        & \vdots & \vec{\mathcal{W}}_1 & \vec{\mathcal{W}}_2 & \cdots & \vec{\mathcal{W}}_{n-r} \\
            0      & \cdots & \cdots & 0      & \downarrow          & \downarrow          &        & \downarrow
        \end{pmatrix}
    \]
    Similarly, the columns of $Q$ are orthonormal and therefore $Q$ is unitary and
    \begin{align*}
        Q^{-1}(P^\dagger H P)Q = Q^\dagger(P^\dagger H P)Q &= (PQ)^\dagger H (PQ) \\
        &= \diag(\lambda_1, \cdots, \lambda_r, \lambda_{r+1}, \cdots, \lambda_n)
    \end{align*}
    Then the $n$ linearly independant eigenvectors will be the columns of the matrix $PQ$.
    
    %%%%%%%%%%%%%%%%%%%%%%
    % Lecture 22: 24/11/2023
    %%%%%%%%%%%%%%%%%%%%%%
    It follows that $H$ is diagonalisable because it is similar to the diagonal matrix $D$ via the transform $U = PQ$.
    $U$ is unitary since $P$ and $Q$ are
    \begin{equation}\label{eq:5-24}
        D = U^\dagger H U \Leftrightarrow H = U D U^\dagger
    \end{equation}
\end{proof}
\begin{remark}
    A real symmetric matrix $S$ is a special case of a Hermition Matrix.
    Now, have
    \begin{equation}\label{eq:5-25}
        D = Q^T S Q \Leftrightarrow S = Q D Q^T
    \end{equation}
    where $Q$ is an orthogonal matrix.
\end{remark}

\begin{eg}
    Take $S = \begin{pmatrix}
        1 & \beta \\
        \beta & 1
    \end{pmatrix}$, for $\beta \in \R$. It has eigenvalues,
    \begin{align*}
        \det(S - \lambda I) = 0 \\
        \Rightarrow (1 - \lambda)^2 - \beta^2 = 0
    \end{align*}
    Therefore $\lambda = 1 \pm \beta$ and the corresponding Eigenvectors
    \[
        \frac{1}{\sqrt{2}}\begin{pmatrix}
            1 \\ \pm 1
        \end{pmatrix}  
    \]
    We change from the standard basis to the eigenvectors
    \[
        Q = \begin{pmatrix}
            \frac{1}{\sqrt{2}} & \frac{1}{\sqrt{2}} \\
            \frac{1}{\sqrt{2}} & - \frac{1}{\sqrt{2}}
        \end{pmatrix}
    \]
    Then by \eqref{eq:5-25},
    \begin{align*}
        S = Q D Q^T \\
        \Rightarrow D = \begin{pmatrix}
            1 + \beta & 0 \\
            0 & 1 - \beta
        \end{pmatrix}
    \end{align*}
\end{eg}

\subsubsection{Normal Matrices}
\begin{defi}[Normal Matrix]
    A \emph{normal matrix} commutes with its Hermition conjugate
    \begin{equation}\label{eq:5-26}
        N^\dagger N = N N^\dagger
    \end{equation}
    Hermition, Symmetric, Skew-Hermition, Anti-symmetric, Orthogonal, Unitary are all special cases of normal matrices.
\end{defi}
It can be shown that
\begin{enumerate}
    \item The eigenvectors corresponding to distinct eigenvalues are orthogonal (exactly as \eqref{eq:5-22} a similar proof)
    \item A normal matrix can always be diagonalised with an orthonormal basis of eigen vectors (including the degenerate case) (as in \eqref{eq:5-23})
    \item If matrix $A$ is similar to a Diagonal matrix via a unitary transform, then $A$ is normal.
    \begin{proof}
        Let $A = U D U^\dagger$ for $U$ unitary.
        \begin{align*}
            A A^\dagger &= U D U^\dagger (U D U^\dagger)^\dagger \\
            &= U D U^\dagger U D^\dagger U^\dagger \\
            &= U DD^\dagger U^\dagger \\
            &= U D^\dagger D U^\dagger \\
            &= A A^\dagger
        \end{align*}
    \end{proof}
\end{enumerate}
\subsection{Gershgorin Circle Theorem}
\begin{thm}[Gershgorin Circle Theorem]
    Consider an $n \times n$ matrix A. For each row define
    \[
        R_i = \sum_{j \neq i}\abs{A_{ij}} \text{  for } i = 1, 2, \cdots, n
    \]
    and the disc centered around the $i^{th}$ diagonal element,
    \[
        \mathcal{D}_i(A_{ii}, R_i) \subset \C
    \]
    with radius $R_i$.
    Then every eigenvalue lies within a disk.
\end{thm}
\begin{proof}
    Consider $A \vec x = \lambda \vec x$. Take the $k^{th}$ component.
    \[
        (A \vec x)_k = \lambda \vec x_k \tag{$\star$}
    \]
    and choose $k$ such that $\vec x_k$ is the element of $\vec x$ with the biggest modulus ($\star$), so
    \begin{align*}
        \sum_{j=1}^{n} A_{kj} \vec x_j = \lambda \vec x_k \\
        \Rightarrow \sum_{j\neq k}^{n} A_{kj} \vec x_j = (\lambda - A_{kk}) \vec x_k
    \end{align*}
    Now divide by $\vec x_k$ and take modulus,
    \[
        \abs{\sum_{j\neq k}^{n} A_{kj} \frac{\vec x_j}{\vec x_k}} = \abs{\lambda - A_{kk}}  
    \]
    Comparing,
    \[
        \text{LHS} \leq \sum_{j\neq k}^{n} \abs{A_{kj}} \abs{\frac{\vec x_j}{\vec x_k}} \leq \sum_{j\neq k}^{n} \abs{A_{kj}} = R_{k}
    \]
    Since 
    \begin{itemize}
        \item By the triangle inequality $\abs{a + b} \leq \abs{a} + \abs{b}$
        \item $\vec x_k$ is chosen to be the largest magnitutde component so $\abs{\frac{\vec x_j}{\vec x_k}} \leq 1$
    \end{itemize}
    Therefore
    \[
        \abs{\lambda - A_{kk}} \leq R_{k}
    \]
    and the eigenvalue $\lambda$ lies in $\mathcal{D}_k(A_{kk}, R_k)$
\end{proof}
\begin{remark}\leavevmode
    \begin{enumerate}
        \item If $A$ is diagonal then the disks are just points
        \item Theorem also applies to $A^T$
        \item Can extend to degeneracy
    \end{enumerate}
\end{remark}
\begin{eg}
    \[
        \begin{pmatrix}
        10 & 0 & 0.2 \\
        1 & 3 & 0.1 \\
        1 & -1 & 2
        \end{pmatrix}  
    \]
    \begin{alignat*}{4}
        \text{eigenvalues:} &\qquad 10.0213 &&\qquad 2.9473 &\qquad 2.05393 \\
        \text{from $A$:} &\qquad \mathcal{D}(10, 0.2) &&\qquad \mathcal{D}(3, 1.1) &\qquad \mathcal{D}(2, 2) \\
        \text{from $A^T$:} &\qquad \mathcal{D}(10, 2) &&\qquad \mathcal{D}(3, 1) &\qquad \mathcal{D}(2, 0.3)
    \end{alignat*}
\end{eg}

\section{Quadratic forms \& Conics}
\begin{defi}[Sesquilinear forms]
    The quantity
    \begin{equation}\label{eq:6-1}
        F = \vec x^\dagger A \vec x = \vec{x}_i^\dagger A_{ij} \vec x_i
    \end{equation}
    is called a \emph{sesquilinear form}.
    \begin{itemize}
        \item If $A = H$ a hermition matrix then it is a \emph{hermition form}
        \item If $A = S$ a symmetric matrix then it is a \emph{quadratic form} 
    \end{itemize}
\end{defi}

Hermition forms are real,
\begin{align*}
    (\vec x^\dagger H \vec x)^* &= (\vec x^\dagger H \vec x)^\dagger \tag{since it is a ``scalar''}\\
    &= \vec x^\dagger H^\dagger \vec x = \vec x^\dagger H \vec x
\end{align*}
We know that any hermition matrix can be diagonalised with a unitary transformation. Hence,
\begin{align*}
    F(\vec x) = \vec x^\dagger H \vec x &= \vec x^\dagger U D U^\dagger \vec x \tag{by \eqref{eq:5-24}} \\
    &= (\vec x')^\dagger D (\vec x')
\end{align*}
For $D = \diag(\lambda_1, \lambda_2, \cdots, \lambda_n)$ a diagonal matrix and $\vec x' = U^\dagger \vec x$.
From \eqref{eq:5-13} we know that $\vec x'$ is the vector $\vec x$ relative to the eigenvector basis.
\begin{equation}\label{eq:6-2}
    F(\vec x) = \sum_{i=1}^{n} \lambda_i \abs{x'_i}^2
\end{equation}
The eigenvectors are known as the principle axes. 
We can think of the $\lambda_i$'s as scalars which are rank $0$ tensors which are invariant.

\begin{eg}
    Consider $F = 2x^2 - 4xy + 5y^2$. 
    Write 
    \[
        F = \vec x^T S \vec x
    \] 
    where $S$ is symmetric.
    Note that if $A$ is anti-symmetric $A_{ij}x_ix_j = 0$. Also, all matrices can be written as a sum of symmetric and anti-symmetric matrices.
    \[
        F = \begin{pmatrix}
            x \\ y
        \end{pmatrix}
        \begin{pmatrix}
            2 & -2 \\
            -2 & 5
        \end{pmatrix}
        \begin{pmatrix}
            x \\ y
        \end{pmatrix}
    \]
    The eigenvalues and corresponding eigenvectors of $S$ are
    \[
        \lambda = 1, \quad \vec x_{1} = \frac{1}{\sqrt{5}}\begin{pmatrix}
            2 \\ 1
        \end{pmatrix}, \text{  and  },
        \lambda = 6, \quad \vec x_{6} = \frac{1}{\sqrt{5}}\begin{pmatrix}
            1 \\ -2
        \end{pmatrix}
    \]
    So, $S = Q D Q^T$ where 
    \[
        Q = \begin{pmatrix}
            \frac{2}{\sqrt{5}} & \frac{1}{\sqrt{5}} \\
            \frac{1}{\sqrt{5}} & \frac{-2}{\sqrt{5}}
        \end{pmatrix}
    \]
    Therefore,
    \[
        F = (x')^2 + 6(y')^2
    \]
    where
    \[
        x' = Q^T \vec x = \begin{pmatrix}
            \frac{2x + y}{\sqrt{5}} \\
            \frac{x - 2y}{\sqrt{5}}
        \end{pmatrix}
    \]
    % sketch elipse on an angle
    Therefore, we have an ellipse where the major and minor axes are given by
    \begin{align*}
        F &= (x')^2 + 6(y')^2 \\
        &= \bkt{\frac{2x + y}{\sqrt{5}}}^2 + 6 \bkt{\frac{x - 2y}{\sqrt{5}}}^2
    \end{align*}
\end{eg}

%%%%%%%%%%%%%%%%%%%%%%
% Lecture 23: 27/11/2023
%%%%%%%%%%%%%%%%%%%%%%
\subsection{Quadratics and Conics}
A quadratic is an $n$-dimensional surface defined by the zero of a real quadratic polynomial. For instance,
\[
    \vec x^T A \vec x + \vec b^T \vec x + c = 0
\]
A is a real $n \times n$ matrix, $\vec x$ and $\vec b$ are $n$-dimensional column vectors and $c$ is a constant.
As noted, if $A$ is antisymmetric then $\vec x^T A \vec x = 0$ so split general $A$ into symmetric and anti-symmetric bits.
\[
    S = \frac{1}{2}(A + A^T)  
\]
then $S = S^T$ and the equation becomes
\begin{equation}\label{eq:6-3}
    \vec x^T S \vec x + \vec b^T \vec x + c = 0
\end{equation}

$S$ is a symmetric matrix and therefore diagonalisable \eqref{eq:5-25} therefore $S = Q D Q^T$ and write $\vec x' = Q^T \vec x$ and $\vec b' = Q^T \vec b$
\begin{equation}\label{eq:6-4}
    \vec{x}'^T D \vec x' + \vec{b}'^T \vec x' + c = 0
\end{equation}
If $S$ has no zero eigenvalues then $D$ is invertible and write
\[
    \vec{x}{''} = \vec x' + \frac{1}{2} D^{-1} \vec b' 
\]
into \eqref{eq:6-4}. This shifts the origin to eliminate the linear term. Finally, drop the ${''}$ superfix
\begin{equation}\label{eq:6-5}
    \vec x^T D \vec x = k
\end{equation}
\subsubsection*{Conic Sections $n = 2$}
If we let $D = \diag(\lambda_1, \lambda_2)$
\[
    \text{\eqref{eq:6-5}} \Rightarrow \lambda_1 x^2 + \lambda_2 x^2 = k    
\]
\begin{enumerate}
    \item $\lambda_1 \lambda_2 > 0$ (non zero and same sign), we require $\sgn(\lambda_i) = \sgn(k)$.
    This produces an ellipse with axes coinciding with eigenvectors of $S$

    \item $\lambda_1 \lambda_2 < 0$ (hyperbola) Wlog, let
    \[
        \lambda_1 = \frac{k}{a^2} > 0 \text{ and } \lambda_2 = -\frac{k}{b^2} < 0  
    \]
    Then,
    \[
        \bkt{\frac{x_1}{a}}^2 + \bkt{\frac{x_2}{a}}^2 = 1
    \]

    \item $\lambda_1 \lambda_2 = 0$ Wlog, take $\lambda_1 \neq 0$, $\lambda_2 = 0$ then $D$ is not invertible, so we take a step back.
    \[
        \text{\eqref{eq:6-4}} \Rightarrow \lambda(x'_1)^2 + b'_1 x'_1 + b'_2 x'_2 + c = 0
    \]
    Shift the origin again,
    \[
        x_1{''} = x'_1 + \frac{b'_1}{2 \lambda_1} \text{ and } x_2{''} = x'_2 + \frac{c}{b'_2} - \frac{b_1^2}{4 \lambda_1 b'_2}
    \]
    Finally, we drop ${''}$ and
    \[
        \lambda_1 x_1^2 + b'_2 x_2 = 0
    \]
    which is a parabola. 
    If $b_2' = 0$ then find
    \[
        x_1^2 = \frac{b_1^2}{4 \lambda_1^2} - \frac{c}{\lambda_1}  
    \]
    which is either 2 straight lines or not solution depending on the sign of the RHS.
\end{enumerate}
\subsection{Focus-Directrix Property}
Consider two control parameters
\begin{enumerate}
    \item $e$ the eccentricity of orbit
    \item $a$ the scale
\end{enumerate}
Then define the foci $\equiv (\pm ae, 0)$ points and the directrices $x = \pm \frac{a}{e}$ lines.

\begin{defi}[Conic Secion]
    A \emph{conic section} is the set of points satisfying the property that distance from the focus = $e \times$ distance from closer directrix,
    except when $e=1$, then take the distance to the other directrix.
\end{defi}

\begin{enumerate}[cases]
    \item ($e < 1$)
    \begin{align*}
        \sqrt{[(x - ae) + y]} = e \bkt{\frac{a}{e} - x} \\
        \Rightarrow \frac{x^2}{a^2} + \frac{y^2}{a^2(1 - e^2)} = 1
    \end{align*}
    which is an ellipse with semi-major axis $a$ and semi-minor axis $a \sqrt{1 - e^2}$
    \item ($e > 1$) 
    \begin{align*}
        \sqrt{[(x - ae) + y]} = e \bkt{x - \frac{a}{e}} \\
        \Rightarrow \frac{x^2}{a^2} - \frac{y^2}{a^2(e^2 - 1)} = 1
    \end{align*}
    which is a hyperbola.
    \item ($e = 1$) 
    \begin{align*}
        \sqrt{[(x - ae) + y]} = (x + a) \\
        \Rightarrow y^2 = 4ax
    \end{align*}
    which is a quadratic
\end{enumerate}

\subsubsection*{Work in polar coordinates centered on a focus}
Call the distance between the focus and directrix
\[
    \frac{l}{e} \equiv \abs{\frac{a}{e} - ae} \Rightarrow l = a\abs{1 - e^2}  
\]
From the focus directrix property
\[
    r = e(\frac{l}{e} - r \cos \theta)
\]
Therefore $\forall e$
\begin{equation}\label{eq:6-6}
    r = \frac{l}{1 + e \cos \theta}
\end{equation}
We see that $r \rightarrow \infty$ if $\theta \rightarrow \cos^{-1}(\frac{-1}{e})$.
This is only possible if $e \geq 1$ i.e, parabola or hyperbola.
We get an ellipse or circle if $e < 1$ and the orbit is bounded. 

%%%%%%%%%%%%%%%%%%%%%%
% Lecture 24: 29/11/2023
%%%%%%%%%%%%%%%%%%%%%%

\section{Transformation Groups}
\subsection{Groups of orthogonal matrices}
The set of all $n \times n$ orthogonal matrices $P$ form a group under matrix multiplication.
\begin{enumerate}
    \item Closure: if $P$ and $Q$ are orthogonal, $R = PQ$
    \[
        R^TR = (PQ)^T(PQ) = Q^TP^TPQ = I  
    \]
    \item Associative: matrix multiplication is assocaitive
    \item Identity: $I$ is orthogonal
    \item Inverse: if $P$ is orthogonal, then $P^{-1} = P^T$ which is orthogonal.
\end{enumerate}
The group is called $\Or(n)$. 
The group of orthogonal matrices with $\det = 1$ is the special orthogonal group $\SO(n)$.

We can show that any matrix in $\Or(2)$ is of the form
\[
    \begin{pmatrix}
        \cos \theta & - \sin \theta \\
        \sin \theta & \cos \theta
    \end{pmatrix}
    \text{ or }
    \begin{pmatrix}
        \cos \theta & \sin \theta \\
        \sin \theta & -\cos \theta
    \end{pmatrix}
\]
\subsection{Length preserving matrices}
Let $P \in \Or(n)$ then the following are equivalent.
\begin{enumerate}
    \item $P$ is orthogonal
    \item $\abs{P\vec x} = \abs{\vec x}$
    \item $(P\vec x)^T (P \vec y) = \vec x^T \vec y$
    \item If $(\vec v_1, \cdots, \vec v_n)$ are orthonormal, so are $(P\vec v_1, \cdots, P\vec v_n)$
    \item Columns of $P$ are orthonormal
\end{enumerate}
\begin{proof}\leavevmode
    \begin{description}
        \item[(i) $\Rightarrow$ (ii)] 
        \[
            \abs{P\vec x}^2 = (P\vec x)^T(P\vec x) = \vec x^T P^TP \vec x = \abs{x}^2
        \]
        \item[(ii) $\Rightarrow$ (iii)] From (ii), $\abs{P(\vec x + \vec y)}^2 = \abs{\vec x + \vec y}^2$
        \begin{align*}
            \text{RHS: } (x + y)^T(x + y) &= x^Tx + y^Ty + 2x^Ty \\
            &= \abs{x}^2 + \abs{y}^2 + 2x^Ty \\ 
            \text{LHS: } \abs{P\vec x + p\vec y}^2 &= \abs{P\vec x}^2 + \abs{P\vec y}^2 + 2(P \vec x)^T(P \vec y) \\
            &= \abs{x}^2 + \abs{y}^2 + 2(P \vec x)^T(P \vec y)
        \end{align*}
        Hence $(P \vec x)^T(P \vec y) = \vec x^T \vec y$
        \item[(iii) $\Rightarrow$ (iv)] From (iii) $(P \vec v_i)^T(P \vec v_j) = \vec v_i^T \vec v_j$.
        So if 
        \[
            \vec v_i^T \vec v_j = \delta_{ij} \Rightarrow (P \vec v_i)^T(P \vec v_j) = \delta_{ij}
        \]
        \item[(iv) $\Rightarrow$ (v)] From (iv) take $\crbkt{\vec v_1, \cdots, \vec v_n}$ to be the standard basis, then 
        \[
            \crbkt{P\vec v_1, \cdots, P\vec v_n}  
        \]
        are the columns of $P$
        \item[(v) $\Rightarrow$ (i)] Columns of $P$ are orthonormal $\Rightarrow$ $P^TP = I$
    \end{description}
\end{proof}
\subsection{Lorentz Transformation}
Consider \emph{Minkowski} (1 + 1) dimensional spacetime. 
The Minkwoski inner product of $2$ vectors $\vec x$, $\vec y$ is 
\begin{equation}\label{eq:7-1}
    \braket{\vec x}{\vec y} \equiv \vec x^T J \vec y
\end{equation}
Where $J = \begin{pmatrix}
    1 & 0 \\
    0 & -1
\end{pmatrix}$. It follows that
\begin{equation}\label{eq:7-2}
    \braket{\vec x}{\vec y} = x_1y_1 - x_2y_2
\end{equation}

Now consider a transformation matrix $M$ which preserves the minkowski inner product,
\begin{align*}
    \braket{\vec x}{\vec y} &= \braket{M\vec x}{M\vec y} \\
    \vec x^T J \vec y &= (M\vec x)^T J (M\vec y) \\
    &= \vec x^T M^T J M \vec y
\end{align*}
Since this must be true $\forall \vec x, \vec y$
\begin{equation}\label{eq:7-3}
    J = M^T J M
\end{equation}
We can show that $M$ is of the form
\[
    H_\alpha \equiv \begin{pmatrix}
        \cosh \alpha & \sinh \alpha \\
        \sinh \alpha & \cosh \alpha
    \end{pmatrix}
\]
We call this a hyperbolic rotation, or 
\[
    K_\frac{\alpha}{2} \equiv \begin{pmatrix}
        \cosh \alpha & - \sinh \alpha \\
        \sinh \alpha & - \cosh \alpha
    \end{pmatrix}
\]
which is a hyperbolic reflection.
Note that in order to derive these, we must insist that $M_{11} > 0$ (to remain in the forwards / backwards light cone)

\begin{defi}[Lorentz Matrix]
    Define the Lorentz Matrix
    \[
        B_v = \frac{1}{\sqrt{1 - v^2}}\begin{pmatrix}
            1 & v \\
            v & 1
        \end{pmatrix}  
    \]
    $B_v$ is called a Lorentz Boost
    \begin{equation}\label{eq:7-4}
        B_v = H_{\tanh^{-1} v}
    \end{equation}
\end{defi}

The set of all $B_v$ forms a group, the Lorentz Group. 
Note that closure follows from $B_{v_1}B_{v_2} = B_{v_3}$
\[
    v_3 = \tanh \bkt{\tanh^{-1} v_1 + \tanh^{-1} v_2} = \frac{v_1 + v_2}{1 + v_1 v_2}    
\]
$B_v$ is a group of transformations which preserve Minkowski inner product.
\end{document}